% !TEX program = lualatex
% ---------------------------------------------------------------------------------------------
%  Readback note 46: Discrete series coefficients near the identity, Laplace's method for two coefficients, and Bessel functions with a peak
%
%  Build:  python3 tools/rb.py note <slug>
%
%  One unit of a Lean artifact, rewritten as ordinary mathematics. The prose is hand-written. Lean names,
%  line ranges and counts come from data.tex; quoted Lean statements are copied from the artifact by the build.
%
%  Rules of the text: standard terminology; every symbol is explained where it first appears and listed again
%  in the notation at the end; no letter has two meanings; a proof is complete or says what it takes for granted;
%  Lean names appear only in the LEAN lines and in the concordance.
%
%  The signature rule (the build enforces it): a function, relation or predicate that the artifact defines is
%  written only through its notation macro, which the build makes from the spec (`notation`) with one slot for
%  each explicit argument of the Lean definition. No argument is fixed and dropped, also not "for brevity" and
%  not for one section. Every argument and every hypothesis of a stated declaration stands inside its
%  statement, or in a Setting that the statement cites by number; only implicit and instance arguments (the
%  ambient objects) may stand in the conventions.
% ---------------------------------------------------------------------------------------------
\documentclass[11pt,a4paper]{article}
\ifdefined\pdfvariable\pdfvariable suppressoptionalinfo \numexpr 1+2+4+8\relax\fi
\usepackage[a4paper,hmargin=35mm,top=22mm,bottom=27mm,footskip=12mm,nohead]{geometry}

\usepackage{fontspec}
\directlua{
  luaotfload.add_fallback("gfallback", {"STIXTwoMath-Regular.otf:mode=node;", "DejaVuSans.ttf:mode=node;"})
  luaotfload.add_fallback("gmonofallback", {"DejaVuSansMono.ttf:mode=node;", "STIXTwoMath-Regular.otf:mode=node;"})
}
\setmainfont{NewCM10}[Extension=.otf,UprightFont=*-Book,ItalicFont=*-BookItalic,
  BoldFont=*-Bold,BoldItalicFont=*-BoldItalic,SmallCapsFont=*-Book,SmallCapsFeatures={Letters=SmallCaps},
  RawFeature={fallback=gfallback}]
\setsansfont{NewCMSans10}[Extension=.otf,UprightFont=*-Book,ItalicFont=*-BookOblique,
  BoldFont=*-Bold,BoldItalicFont=*-BoldOblique,RawFeature={fallback=gfallback}]
\setmonofont{JetBrainsMono}[Extension=.otf,UprightFont=*-Regular,ItalicFont=*-Italic,
  BoldFont=*-SemiBold,BoldItalicFont=*-SemiBoldItalic,Scale=MatchLowercase,
  RawFeature={-calt;-liga;fallback=gmonofallback}]
\usepackage{amsmath}
\usepackage{amsthm}
\usepackage[math-style=TeX]{unicode-math}
\setmathfont{NewCMMath-Book.otf}
\emergencystretch=2.5em\relax

\usepackage{xcolor}
\usepackage{tikz}
\usetikzlibrary{calc,arrows.meta,positioning,backgrounds}
\usepackage[protrusion=true,expansion=true]{microtype}
\usepackage{enumitem}
\usepackage{needspace}
\usepackage{booktabs}
\usepackage{array}
\usepackage{longtable}
\usepackage{fancyhdr}
\usepackage{lastpage}
\usepackage{multicol}
\usepackage[hidelinks,pdfusetitle]{hyperref}

\definecolor{ink}{HTML}{0F172A}
\definecolor{inksoft}{HTML}{475569}
\definecolor{inkfaint}{HTML}{94A3B8}
\definecolor{rulegrey}{HTML}{CBD5E1}
\definecolor{arith}{HTML}{C2410C}
\definecolor{est}{HTML}{7C3AED}
\definecolor{ana}{HTML}{0891B2}
\definecolor{ok}{HTML}{047857}
\definecolor{leancol}{HTML}{3730A3}

\title{Discrete series coefficients near the identity, Laplace's method for two coefficients, and Bessel functions with a peak}
\author{grh-compact}
\date{}

% ------------------------------------------------------------------------------------- macros
\newcommand{\leanfont}{\ttfamily\fontsize{8.6}{10.5}\selectfont\color{leancol}}
\newcommand{\lean}[1]{\csname LN@#1\endcsname}
\newcommand{\fillin}[1]{{\color{arith}\itshape[#1]}}  % a place of the template that is not written yet; the build reports every one that is left
\newcommand{\leanquote}[1]{\csname LQ@#1\endcsname}
% generated by rb.py note; do not edit
\newcommand{\ConcDecls}{35}
\newcommand{\Pin}{72fead10}
\newcommand{\PinText}{that of commit {\leanfont 72fead10}}
\newcommand{\ArtifactName}{grh-compact}
\expandafter\newcommand\csname LN@Frame\endcsname{\href{https://github.com/lsernau/grh-lean/blob/72fead102f663d5c7a407964a13dc0cf9d977207/grh-compact/src/New/DiscreteSeries/BesselExistence.lean\#L16-L17}{{\leanfont Frame}}}
\expandafter\newcommand\csname LN@IsBessel\endcsname{\href{https://github.com/lsernau/grh-lean/blob/72fead102f663d5c7a407964a13dc0cf9d977207/grh-compact/src/New/DiscreteSeries/BesselExistence.lean\#L21-L26}{{\leanfont Filling.\allowbreak{}IsBessel}}}
\expandafter\newcommand\csname LN@besselGraph\endcsname{\href{https://github.com/lsernau/grh-lean/blob/72fead102f663d5c7a407964a13dc0cf9d977207/grh-compact/src/New/DiscreteSeries/BesselExistence.lean\#L34-L37}{{\leanfont Shape.\allowbreak{}besselGraph}}}
\expandafter\newcommand\csname LN@besselGraphChart\endcsname{\href{https://github.com/lsernau/grh-lean/blob/72fead102f663d5c7a407964a13dc0cf9d977207/grh-compact/src/New/DiscreteSeries/BesselExistence.lean\#L45-L47}{{\leanfont Shape.\allowbreak{}besselGraphChart}}}
\expandafter\newcommand\csname LN@besselGraphMeasure\endcsname{\href{https://github.com/lsernau/grh-lean/blob/72fead102f663d5c7a407964a13dc0cf9d977207/grh-compact/src/New/DiscreteSeries/BesselExistence.lean\#L51-L51}{{\leanfont Shape.\allowbreak{}besselGraphMeasure}}}
\expandafter\newcommand\csname LN@besselGraphMeasurefiniteOnCompacts\endcsname{\href{https://github.com/lsernau/grh-lean/blob/72fead102f663d5c7a407964a13dc0cf9d977207/grh-compact/src/New/DiscreteSeries/BesselExistence.lean\#L56-L57}{{\leanfont Shape.\allowbreak{}besselGraphMeasure\_finiteOnCompacts}}}
\expandafter\newcommand\csname LN@besselGraphMeasureisMulRightInvariant\endcsname{\href{https://github.com/lsernau/grh-lean/blob/72fead102f663d5c7a407964a13dc0cf9d977207/grh-compact/src/New/DiscreteSeries/BesselExistence.lean\#L59-L63}{{\leanfont Shape.\allowbreak{}besselGraphMeasure\_isMulRightInvariant}}}
\expandafter\newcommand\csname LN@besselGraphfstmembesselS\endcsname{\href{https://github.com/lsernau/grh-lean/blob/72fead102f663d5c7a407964a13dc0cf9d977207/grh-compact/src/New/DiscreteSeries/BesselExistence.lean\#L67-L68}{{\leanfont besselGraph\_fst\_mem\_besselS}}}
\expandafter\newcommand\csname LN@lintegralunfoldbesselGraphMeasure\endcsname{\href{https://github.com/lsernau/grh-lean/blob/72fead102f663d5c7a407964a13dc0cf9d977207/grh-compact/src/New/DiscreteSeries/BesselExistence.lean\#L73-L77}{{\leanfont lintegral\_unfold\_besselGraphMeasure}}}
\expandafter\newcommand\csname LN@besselGraphtracemulsubeqzero\endcsname{\href{https://github.com/lsernau/grh-lean/blob/72fead102f663d5c7a407964a13dc0cf9d977207/grh-compact/src/New/DiscreteSeries/BesselExistence.lean\#L83-L87}{{\leanfont besselGraph\_trace\_mul\_sub\_eq\_zero}}}
\expandafter\newcommand\csname LN@besselPairHom\endcsname{\href{https://github.com/lsernau/grh-lean/blob/72fead102f663d5c7a407964a13dc0cf9d977207/grh-compact/src/New/DiscreteSeries/BesselExistence.lean\#L91-L93}{{\leanfont besselPairHom}}}
\expandafter\newcommand\csname LN@HasPeakAt\endcsname{\href{https://github.com/lsernau/grh-lean/blob/72fead102f663d5c7a407964a13dc0cf9d977207/grh-compact/src/New/DiscreteSeries/BesselExistence.lean\#L98-L103}{{\leanfont HasPeakAt}}}
\expandafter\newcommand\csname LN@existsisTwoFactorbesselPairHom\endcsname{\href{https://github.com/lsernau/grh-lean/blob/72fead102f663d5c7a407964a13dc0cf9d977207/grh-compact/src/New/DiscreteSeries/BesselExistence.lean\#L108-L173}{{\leanfont IsSteep.\allowbreak{}exists\_isTwoFactor\_besselPairHom}}}
\expandafter\newcommand\csname LN@isProperMapfstbesselPairHom\endcsname{\href{https://github.com/lsernau/grh-lean/blob/72fead102f663d5c7a407964a13dc0cf9d977207/grh-compact/src/New/DiscreteSeries/BesselExistence.lean\#L177-L179}{{\leanfont isProperMap\_fst\_besselPairHom}}}
\expandafter\newcommand\csname LN@existsisBesselFunction\endcsname{\href{https://github.com/lsernau/grh-lean/blob/72fead102f663d5c7a407964a13dc0cf9d977207/grh-compact/src/New/DiscreteSeries/BesselExistence.lean\#L185-L225}{{\leanfont exists\_isBesselFunction}}}
\expandafter\newcommand\csname LN@gapWeightedNormSq\endcsname{\href{https://github.com/lsernau/grh-lean/blob/72fead102f663d5c7a407964a13dc0cf9d977207/grh-compact/src/New/DiscreteSeries/GaussianShape.lean\#L16-L16}{{\leanfont gapWeightedNormSq}}}
\expandafter\newcommand\csname LN@sumnormliecoordProjsq\endcsname{\href{https://github.com/lsernau/grh-lean/blob/72fead102f663d5c7a407964a13dc0cf9d977207/grh-compact/src/New/DiscreteSeries/GaussianShape.lean\#L18-L25}{{\leanfont sum\_norm\_lie\_coordProj\_sq}}}
\expandafter\newcommand\csname LN@normweightMinorssqleexp\endcsname{\href{https://github.com/lsernau/grh-lean/blob/72fead102f663d5c7a407964a13dc0cf9d977207/grh-compact/src/New/DiscreteSeries/GaussianShape.lean\#L27-L41}{{\leanfont norm\_weightMinors\_sq\_le\_exp}}}
\expandafter\newcommand\csname LN@dsCoeffgaussianshape\endcsname{\href{https://github.com/lsernau/grh-lean/blob/72fead102f663d5c7a407964a13dc0cf9d977207/grh-compact/src/New/DiscreteSeries/GaussianShape.lean\#L48-L99}{{\leanfont dsCoeff\_gaussian\_shape}}}
\expandafter\newcommand\csname LQ@dsCoeffgaussianshape\endcsname{\leanquoteline{0}{lemma dsCoeff\_gaussian\_shape [Nonempty n]}\leanquoteline{4}{(hg : l.HasLargeGap A) (hL : ∀ a, |l a| ≤ L) \{g : USign ε\} (hZ : ‖g.1.1}\leanquoteline{2}{- 1‖ ≤ δ) (hT : 1 ≤ T) (a : card n \textasciicircum{} 2 ≤ T := by nlinarith) (b : card n \textasciicircum{} 2 * Real.log (card n +}\leanquoteline{2}{L + 1) ≤ T := by bound) (hδ : δ * T ≤ 1e-12 := by bound) : ∃ κ ∈ Set.Icc (0 : ℝ) 1, ‖l.dsCoeff g /}\leanquoteline{2}{(intDiagonal l.lowestKType * logOneAdd (g.1.1 - 1)).trace.exp - κ‖ ≤ κ / 10 + (-100 * T).exp}}
\expandafter\newcommand\csname LN@existsdetcompresseqexplogOneAdd\endcsname{\href{https://github.com/lsernau/grh-lean/blob/72fead102f663d5c7a407964a13dc0cf9d977207/grh-compact/src/New/DiscreteSeries/MinorExpansion.lean\#L18-L44}{{\leanfont exists\_det\_compress\_eq\_exp\_logOneAdd}}}
\expandafter\newcommand\csname LN@existsprincipalMinoreqexp\endcsname{\href{https://github.com/lsernau/grh-lean/blob/72fead102f663d5c7a407964a13dc0cf9d977207/grh-compact/src/New/DiscreteSeries/MinorExpansion.lean\#L49-L109}{{\leanfont exists\_principalMinor\_eq\_exp}}}
\expandafter\newcommand\csname LN@normpoincareperiodsuble\endcsname{\href{https://github.com/lsernau/grh-lean/blob/72fead102f663d5c7a407964a13dc0cf9d977207/grh-compact/src/New/DiscreteSeries/RestrictionLemma.lean\#L25-L124}{{\leanfont norm\_poincare\_period\_sub\_le}}}
\expandafter\newcommand\csname LN@laplacemethod\endcsname{\href{https://github.com/lsernau/grh-lean/blob/72fead102f663d5c7a407964a13dc0cf9d977207/grh-compact/src/New/DiscreteSeries/RestrictionLemma.lean\#L134-L192}{{\leanfont IsTwoFactor.\allowbreak{}laplace\_method}}}
\expandafter\newcommand\csname LQ@laplacemethod\endcsname{\leanquoteline{0}{lemma IsTwoFactor.laplace\_method \{n m X S\} [Fintype n] [Fintype m] [Fintype X] [Group S]}\leanquoteline{2}{[TopologicalSpace S] [MeasurableSpace S] [MeasurableMul S] [BorelSpace S] (ν : Measure S)}\leanquoteline{2}{[ν.IsMulRightInvariant] [IsFiniteMeasureOnCompacts ν] \{ε l η k τ\} \{B β : ℝ\} \{L M h : ℕ\} [Nonempty}\leanquoteline{2}{n] [Nonempty m] \{e : X ⊕ Fin M ≃ n\} \{j : S →* USign ε × \_\} \{φ\} (H : IsTwoFactor l k j τ B e φ)}\leanquoteline{2}{(hg : l.HasLargeGap h) (hg' : @HasLargeGap m \_ η k h) (hB : 1 ≤ B) (hL : ∀ a, |l a| ≤ L) (hK :}\leanquoteline{2}{∀ a, |k a| ≤ L) (hβ : 10 \textasciicircum{} 12 ≤ β) (hn : card n + card m + L ≤ β) (hh : 10 \textasciicircum{} 30 * ((card n + card}\leanquoteline{2}{m) \textasciicircum{} 2 * β.log) \textasciicircum{} 3 ≤ h) (hj : IsProperMap (j · |>.1)) (hj' : Continuous (j · |>.2)) (lip :}\leanquoteline{2}{∀ s, ‖(subOnePair j s).1‖ ≤ 1 → ‖(j s).2.1.1‖ ≤ β + 1) : ν.real \{s | ‖(subOnePair j s).1‖ ≤ 1\} ≤}\leanquoteline{2}{β \textasciicircum{} (20 * (card n + card m) \textasciicircum{} 2) * ‖∫ s, l.laplaceIntegrand k j φ s ∂ν‖}}
\expandafter\newcommand\csname LN@existsisTwoFactorrestriction\endcsname{\href{https://github.com/lsernau/grh-lean/blob/72fead102f663d5c7a407964a13dc0cf9d977207/grh-compact/src/New/DiscreteSeries/RestrictionLemma.lean\#L200-L243}{{\leanfont IsSteep.\allowbreak{}exists\_isTwoFactor\_restriction}}}
\expandafter\newcommand\csname LN@phase\endcsname{\href{https://github.com/lsernau/grh-lean/blob/72fead102f663d5c7a407964a13dc0cf9d977207/grh-compact/src/New/DiscreteSeries/TwoFactorSetup.lean\#L19-L20}{{\leanfont HCParam.\allowbreak{}phase}}}
\expandafter\newcommand\csname LN@IsSteep\endcsname{\href{https://github.com/lsernau/grh-lean/blob/72fead102f663d5c7a407964a13dc0cf9d977207/grh-compact/src/New/DiscreteSeries/TwoFactorSetup.lean\#L24-L29}{{\leanfont IsSteep}}}
\expandafter\newcommand\csname LN@subOnePair\endcsname{\href{https://github.com/lsernau/grh-lean/blob/72fead102f663d5c7a407964a13dc0cf9d977207/grh-compact/src/New/DiscreteSeries/TwoFactorSetup.lean\#L33-L33}{{\leanfont subOnePair}}}
\expandafter\newcommand\csname LN@laplaceIntegrand\endcsname{\href{https://github.com/lsernau/grh-lean/blob/72fead102f663d5c7a407964a13dc0cf9d977207/grh-compact/src/New/DiscreteSeries/TwoFactorSetup.lean\#L37-L37}{{\leanfont HCParam.\allowbreak{}laplaceIntegrand}}}
\expandafter\newcommand\csname LN@IsTwoFactor\endcsname{\href{https://github.com/lsernau/grh-lean/blob/72fead102f663d5c7a407964a13dc0cf9d977207/grh-compact/src/New/DiscreteSeries/TwoFactorSetup.lean\#L42-L50}{{\leanfont IsTwoFactor}}}
\expandafter\newcommand\csname LN@existsunitarysuboneeqinclMatrix\endcsname{\href{https://github.com/lsernau/grh-lean/blob/72fead102f663d5c7a407964a13dc0cf9d977207/grh-compact/src/New/DiscreteSeries/TwoFactorSetup.lean\#L54-L57}{{\leanfont exists\_unitary\_sub\_one\_eq\_inclMatrix}}}
\expandafter\newcommand\csname LN@normsqaddretracele\endcsname{\href{https://github.com/lsernau/grh-lean/blob/72fead102f663d5c7a407964a13dc0cf9d977207/grh-compact/src/New/DiscreteSeries/TwoFactorSetup.lean\#L59-L68}{{\leanfont norm\_sq\_add\_re\_trace\_le}}}
\expandafter\newcommand\csname LN@phaselogOneAdd\endcsname{\href{https://github.com/lsernau/grh-lean/blob/72fead102f663d5c7a407964a13dc0cf9d977207/grh-compact/src/New/DiscreteSeries/TwoFactorSetup.lean\#L72-L84}{{\leanfont IsTwoFactor.\allowbreak{}phase\_logOneAdd}}}
\expandafter\newcommand\csname LN@dsCoeffmuleqofzeroweightblock\endcsname{\href{https://github.com/lsernau/grh-lean/blob/72fead102f663d5c7a407964a13dc0cf9d977207/grh-compact/src/New/DiscreteSeries/TwoFactorSetup.lean\#L88-L100}{{\leanfont dsCoeff\_mul\_eq\_of\_zero\_weight\_block}}}
\expandafter\newcommand\csname LN@normdsCoeffpowfourle\endcsname{\href{https://github.com/lsernau/grh-lean/blob/72fead102f663d5c7a407964a13dc0cf9d977207/grh-compact/src/New/DiscreteSeries/TwoFactorSetup.lean\#L102-L114}{{\leanfont norm\_dsCoeff\_pow\_four\_le}}}
\expandafter\newcommand\csname LN@existsnearone\endcsname{\href{https://github.com/lsernau/grh-lean/blob/72fead102f663d5c7a407964a13dc0cf9d977207/grh-compact/src/New/DiscreteSeries/TwoFactorSetup.lean\#L119-L148}{{\leanfont IsTwoFactor.\allowbreak{}exists\_near\_one}}}
\newcommand{\ConcordanceRows}{\crow{Frame}{\href{https://github.com/lsernau/grh-lean/blob/72fead102f663d5c7a407964a13dc0cf9d977207/grh-compact/src/New/DiscreteSeries/BesselExistence.lean\#L16-L17}{{\leanfont Frame}}}{BesselExistence.lean}{16--17}
\crow{IsBessel}{\href{https://github.com/lsernau/grh-lean/blob/72fead102f663d5c7a407964a13dc0cf9d977207/grh-compact/src/New/DiscreteSeries/BesselExistence.lean\#L21-L26}{{\leanfont Filling.\allowbreak{}IsBessel}}}{BesselExistence.lean}{21--26}
\crow{besselGraph}{\href{https://github.com/lsernau/grh-lean/blob/72fead102f663d5c7a407964a13dc0cf9d977207/grh-compact/src/New/DiscreteSeries/BesselExistence.lean\#L34-L37}{{\leanfont Shape.\allowbreak{}besselGraph}}}{BesselExistence.lean}{34--37}
\crow{besselGraphChart}{\href{https://github.com/lsernau/grh-lean/blob/72fead102f663d5c7a407964a13dc0cf9d977207/grh-compact/src/New/DiscreteSeries/BesselExistence.lean\#L45-L47}{{\leanfont Shape.\allowbreak{}besselGraphChart}}}{BesselExistence.lean}{45--47}
\crow{besselGraphMeasure}{\href{https://github.com/lsernau/grh-lean/blob/72fead102f663d5c7a407964a13dc0cf9d977207/grh-compact/src/New/DiscreteSeries/BesselExistence.lean\#L51-L51}{{\leanfont Shape.\allowbreak{}besselGraphMeasure}}}{BesselExistence.lean}{51--51}
\crow{besselGraphMeasurefiniteOnCompacts}{\href{https://github.com/lsernau/grh-lean/blob/72fead102f663d5c7a407964a13dc0cf9d977207/grh-compact/src/New/DiscreteSeries/BesselExistence.lean\#L56-L57}{{\leanfont Shape.\allowbreak{}besselGraphMeasure\_finiteOnCompacts}}}{BesselExistence.lean}{56--57}
\crow{besselGraphMeasureisMulRightInvariant}{\href{https://github.com/lsernau/grh-lean/blob/72fead102f663d5c7a407964a13dc0cf9d977207/grh-compact/src/New/DiscreteSeries/BesselExistence.lean\#L59-L63}{{\leanfont Shape.\allowbreak{}besselGraphMeasure\_isMulRightInvariant}}}{BesselExistence.lean}{59--63}
\crow{besselGraphfstmembesselS}{\href{https://github.com/lsernau/grh-lean/blob/72fead102f663d5c7a407964a13dc0cf9d977207/grh-compact/src/New/DiscreteSeries/BesselExistence.lean\#L67-L68}{{\leanfont besselGraph\_fst\_mem\_besselS}}}{BesselExistence.lean}{67--68}
\crow{lintegralunfoldbesselGraphMeasure}{\href{https://github.com/lsernau/grh-lean/blob/72fead102f663d5c7a407964a13dc0cf9d977207/grh-compact/src/New/DiscreteSeries/BesselExistence.lean\#L73-L77}{{\leanfont lintegral\_unfold\_besselGraphMeasure}}}{BesselExistence.lean}{73--77}
\crow{besselGraphtracemulsubeqzero}{\href{https://github.com/lsernau/grh-lean/blob/72fead102f663d5c7a407964a13dc0cf9d977207/grh-compact/src/New/DiscreteSeries/BesselExistence.lean\#L83-L87}{{\leanfont besselGraph\_trace\_mul\_sub\_eq\_zero}}}{BesselExistence.lean}{83--87}
\crow{besselPairHom}{\href{https://github.com/lsernau/grh-lean/blob/72fead102f663d5c7a407964a13dc0cf9d977207/grh-compact/src/New/DiscreteSeries/BesselExistence.lean\#L91-L93}{{\leanfont besselPairHom}}}{BesselExistence.lean}{91--93}
\crow{HasPeakAt}{\href{https://github.com/lsernau/grh-lean/blob/72fead102f663d5c7a407964a13dc0cf9d977207/grh-compact/src/New/DiscreteSeries/BesselExistence.lean\#L98-L103}{{\leanfont HasPeakAt}}}{BesselExistence.lean}{98--103}
\crow{existsisTwoFactorbesselPairHom}{\href{https://github.com/lsernau/grh-lean/blob/72fead102f663d5c7a407964a13dc0cf9d977207/grh-compact/src/New/DiscreteSeries/BesselExistence.lean\#L108-L173}{{\leanfont IsSteep.\allowbreak{}exists\_isTwoFactor\_besselPairHom}}}{BesselExistence.lean}{108--173}
\crow{isProperMapfstbesselPairHom}{\href{https://github.com/lsernau/grh-lean/blob/72fead102f663d5c7a407964a13dc0cf9d977207/grh-compact/src/New/DiscreteSeries/BesselExistence.lean\#L177-L179}{{\leanfont isProperMap\_fst\_besselPairHom}}}{BesselExistence.lean}{177--179}
\crow{existsisBesselFunction}{\href{https://github.com/lsernau/grh-lean/blob/72fead102f663d5c7a407964a13dc0cf9d977207/grh-compact/src/New/DiscreteSeries/BesselExistence.lean\#L185-L225}{{\leanfont exists\_isBesselFunction}}}{BesselExistence.lean}{185--225}
\crow{gapWeightedNormSq}{\href{https://github.com/lsernau/grh-lean/blob/72fead102f663d5c7a407964a13dc0cf9d977207/grh-compact/src/New/DiscreteSeries/GaussianShape.lean\#L16-L16}{{\leanfont gapWeightedNormSq}}}{GaussianShape.lean}{16--16}
\crow{sumnormliecoordProjsq}{\href{https://github.com/lsernau/grh-lean/blob/72fead102f663d5c7a407964a13dc0cf9d977207/grh-compact/src/New/DiscreteSeries/GaussianShape.lean\#L18-L25}{{\leanfont sum\_norm\_lie\_coordProj\_sq}}}{GaussianShape.lean}{18--25}
\crow{normweightMinorssqleexp}{\href{https://github.com/lsernau/grh-lean/blob/72fead102f663d5c7a407964a13dc0cf9d977207/grh-compact/src/New/DiscreteSeries/GaussianShape.lean\#L27-L41}{{\leanfont norm\_weightMinors\_sq\_le\_exp}}}{GaussianShape.lean}{27--41}
\crow{dsCoeffgaussianshape}{\href{https://github.com/lsernau/grh-lean/blob/72fead102f663d5c7a407964a13dc0cf9d977207/grh-compact/src/New/DiscreteSeries/GaussianShape.lean\#L48-L99}{{\leanfont dsCoeff\_gaussian\_shape}}}{GaussianShape.lean}{48--99}
\crow{existsdetcompresseqexplogOneAdd}{\href{https://github.com/lsernau/grh-lean/blob/72fead102f663d5c7a407964a13dc0cf9d977207/grh-compact/src/New/DiscreteSeries/MinorExpansion.lean\#L18-L44}{{\leanfont exists\_det\_compress\_eq\_exp\_logOneAdd}}}{MinorExpansion.lean}{18--44}
\crow{existsprincipalMinoreqexp}{\href{https://github.com/lsernau/grh-lean/blob/72fead102f663d5c7a407964a13dc0cf9d977207/grh-compact/src/New/DiscreteSeries/MinorExpansion.lean\#L49-L109}{{\leanfont exists\_principalMinor\_eq\_exp}}}{MinorExpansion.lean}{49--109}
\crow{normpoincareperiodsuble}{\href{https://github.com/lsernau/grh-lean/blob/72fead102f663d5c7a407964a13dc0cf9d977207/grh-compact/src/New/DiscreteSeries/RestrictionLemma.lean\#L25-L124}{{\leanfont norm\_poincare\_period\_sub\_le}}}{RestrictionLemma.lean}{25--124}
\crow{laplacemethod}{\href{https://github.com/lsernau/grh-lean/blob/72fead102f663d5c7a407964a13dc0cf9d977207/grh-compact/src/New/DiscreteSeries/RestrictionLemma.lean\#L134-L192}{{\leanfont IsTwoFactor.\allowbreak{}laplace\_method}}}{RestrictionLemma.lean}{134--192}
\crow{existsisTwoFactorrestriction}{\href{https://github.com/lsernau/grh-lean/blob/72fead102f663d5c7a407964a13dc0cf9d977207/grh-compact/src/New/DiscreteSeries/RestrictionLemma.lean\#L200-L243}{{\leanfont IsSteep.\allowbreak{}exists\_isTwoFactor\_restriction}}}{RestrictionLemma.lean}{200--243}
\crow{phase}{\href{https://github.com/lsernau/grh-lean/blob/72fead102f663d5c7a407964a13dc0cf9d977207/grh-compact/src/New/DiscreteSeries/TwoFactorSetup.lean\#L19-L20}{{\leanfont HCParam.\allowbreak{}phase}}}{TwoFactorSetup.lean}{19--20}
\crow{IsSteep}{\href{https://github.com/lsernau/grh-lean/blob/72fead102f663d5c7a407964a13dc0cf9d977207/grh-compact/src/New/DiscreteSeries/TwoFactorSetup.lean\#L24-L29}{{\leanfont IsSteep}}}{TwoFactorSetup.lean}{24--29}
\crow{subOnePair}{\href{https://github.com/lsernau/grh-lean/blob/72fead102f663d5c7a407964a13dc0cf9d977207/grh-compact/src/New/DiscreteSeries/TwoFactorSetup.lean\#L33-L33}{{\leanfont subOnePair}}}{TwoFactorSetup.lean}{33--33}
\crow{laplaceIntegrand}{\href{https://github.com/lsernau/grh-lean/blob/72fead102f663d5c7a407964a13dc0cf9d977207/grh-compact/src/New/DiscreteSeries/TwoFactorSetup.lean\#L37-L37}{{\leanfont HCParam.\allowbreak{}laplaceIntegrand}}}{TwoFactorSetup.lean}{37--37}
\crow{IsTwoFactor}{\href{https://github.com/lsernau/grh-lean/blob/72fead102f663d5c7a407964a13dc0cf9d977207/grh-compact/src/New/DiscreteSeries/TwoFactorSetup.lean\#L42-L50}{{\leanfont IsTwoFactor}}}{TwoFactorSetup.lean}{42--50}
\crow{existsunitarysuboneeqinclMatrix}{\href{https://github.com/lsernau/grh-lean/blob/72fead102f663d5c7a407964a13dc0cf9d977207/grh-compact/src/New/DiscreteSeries/TwoFactorSetup.lean\#L54-L57}{{\leanfont exists\_unitary\_sub\_one\_eq\_inclMatrix}}}{TwoFactorSetup.lean}{54--57}
\crow{normsqaddretracele}{\href{https://github.com/lsernau/grh-lean/blob/72fead102f663d5c7a407964a13dc0cf9d977207/grh-compact/src/New/DiscreteSeries/TwoFactorSetup.lean\#L59-L68}{{\leanfont norm\_sq\_add\_re\_trace\_le}}}{TwoFactorSetup.lean}{59--68}
\crow{phaselogOneAdd}{\href{https://github.com/lsernau/grh-lean/blob/72fead102f663d5c7a407964a13dc0cf9d977207/grh-compact/src/New/DiscreteSeries/TwoFactorSetup.lean\#L72-L84}{{\leanfont IsTwoFactor.\allowbreak{}phase\_logOneAdd}}}{TwoFactorSetup.lean}{72--84}
\crow{dsCoeffmuleqofzeroweightblock}{\href{https://github.com/lsernau/grh-lean/blob/72fead102f663d5c7a407964a13dc0cf9d977207/grh-compact/src/New/DiscreteSeries/TwoFactorSetup.lean\#L88-L100}{{\leanfont dsCoeff\_mul\_eq\_of\_zero\_weight\_block}}}{TwoFactorSetup.lean}{88--100}
\crow{normdsCoeffpowfourle}{\href{https://github.com/lsernau/grh-lean/blob/72fead102f663d5c7a407964a13dc0cf9d977207/grh-compact/src/New/DiscreteSeries/TwoFactorSetup.lean\#L102-L114}{{\leanfont norm\_dsCoeff\_pow\_four\_le}}}{TwoFactorSetup.lean}{102--114}
\crow{existsnearone}{\href{https://github.com/lsernau/grh-lean/blob/72fead102f663d5c7a407964a13dc0cf9d977207/grh-compact/src/New/DiscreteSeries/TwoFactorSetup.lean\#L119-L148}{{\leanfont IsTwoFactor.\allowbreak{}exists\_near\_one}}}{TwoFactorSetup.lean}{119--148}}
% the notations of the defined functions: one slot for each explicit argument of the Lean declaration
\newcommand{\frm}[2]{\operatorname{Frm}(#1,#2)}
\newcommand{\isbes}[3]{\operatorname{BesType}(#1;#2,#3)}
\newcommand{\bgr}[1]{\mathsf{HU}(#1)}
\newcommand{\bch}[1]{\operatorname{ch}_{#1}}
\newcommand{\bms}[1]{\nu_{#1}}
\newcommand{\bph}[2]{\jmath_{#1}(#2)}
\newcommand{\gwn}[2]{\operatorname{gw}_{#1}(#2)}
\newcommand{\phs}[3]{\operatorname{ph}_{#1,#2}#3}
\newcommand{\steep}[3]{\operatorname{Steep}(#1;#2,#3)}
\newcommand{\lapint}[5]{\theta_{#1,#2}^{#3,#4}(#5)}
\newcommand{\twof}[7]{\operatorname{TwoFactor}(#1,#2;#3,#7;#6,#4;#5)}
\newcommand{\UnitMembers}{37}
\newcommand{\UnitInNote}{35}
\newcommand{\UnitProofs}{17}
\newcommand{\UnitSteps}{235}
\newcommand{\UnitAllSteps}{260}


\linespread{1.08}                       % 11 pt on 14.7 pt
\setlength{\parindent}{0pt}
\setlength{\parskip}{5.5pt plus 1.5pt minus 0.5pt}
\setlength{\abovedisplayskip}{7.5pt plus 2pt minus 3pt}
\setlength{\belowdisplayskip}{7.5pt plus 2pt minus 3pt}
\setlength{\abovedisplayshortskip}{3pt plus 2pt}
\setlength{\belowdisplayshortskip}{6pt plus 2pt minus 2pt}
\emergencystretch=1.5em
\clubpenalty=10000 \widowpenalty=10000 \displaywidowpenalty=10000
\relpenalty=7000 \binoppenalty=9000
\raggedbottom
\pagestyle{fancy}\fancyhf{}\renewcommand{\headrulewidth}{0pt}
\fancyfoot[C]{\sffamily\fontsize{7.6}{9}\selectfont\color{inkfaint}\thepage\,/\,\pageref*{LastPage}}
\setlist{topsep=2pt,partopsep=0pt,parsep=0pt,itemsep=4pt,leftmargin=2.1em,labelsep=0.55em}

\newcommand{\Q}{\mathbb{Q}}
\newcommand{\Z}{\mathbb{Z}}
\newcommand{\F}{\mathbb{F}}
\newcommand{\len}{\operatorname{len}}
\newcommand{\cok}{\operatorname{cok}}
\newcommand{\End}{\operatorname{End}}
\newcommand{\Hom}{\operatorname{Hom}}
\newcommand{\Gal}{\operatorname{Gal}}
\newcommand{\diag}{\operatorname{diag}}
\newcommand{\kw}[1]{\emph{#1}}
\newcommand{\m}[1]{\mbox{$#1$}}        % a short formula that must not be broken across lines
\newcommand{\head}[1]{\par\addvspace{17pt plus 4pt minus 2pt}\Needspace{6\baselineskip}%
  \noindent{\fontsize{13.2}{16}\selectfont\bfseries\color{ink}#1}\par\nopagebreak\vspace{-1pt}}
\newcommand{\sect}[2]{\refstepcounter{section}\label{#2}\par\addvspace{19pt plus 4pt minus 2pt}\Needspace{7\baselineskip}%
  \noindent\llap{{\fontsize{13.2}{16}\selectfont\bfseries\color{inkfaint}\thesection}\hspace{9pt}}%
  {\fontsize{13.2}{16}\selectfont\bfseries\color{ink}#1}\par\nopagebreak\vspace{-1pt}}
\newcommand{\subhead}[1]{\par\addvspace{9pt plus 2pt}\Needspace{4\baselineskip}\noindent\textbf{#1}\enspace\ignorespaces}
% where a statement lives in the Lean archive
\newcommand{\leantag}{{\sffamily\bfseries\fontsize{6.3}{8}\selectfont\color{leancol}\addfontfeatures{LetterSpace=7}LEAN}}
\newcommand{\sep}{\unskip\nobreak\hspace{5pt}{\color{inkfaint}\textperiodcentered}\hspace{5pt plus 2pt}\ignorespaces}
\newcommand{\leanline}[1]{\par\nopagebreak\vspace{-2.5pt}{\fontsize{8.9}{12.4}\selectfont\color{inksoft}\raggedright
  \noindent\llap{\leantag\hspace{9pt}}#1\par}\vspace{1pt}}
% a displayed statement in a tinted box: #1 colour, #2 text
\newcommand{\statement}[2]{\par\addvspace{7pt}\begingroup\setlength{\fboxsep}{8.5pt}\setlength{\fboxrule}{0.6pt}%
  \noindent\fcolorbox{#1!30!white}{#1!3!white}{\parbox{\dimexpr\linewidth-2\fboxsep-2\fboxrule\relax}{#2}}%
  \endgroup\par\addvspace{7pt}}
\newcommand{\nt}[2]{\par\noindent\hangindent=1.1em\hangafter=1 #1\enspace{\color{inksoft}#2}\par\vspace{1pt}}
\newcommand{\legdot}[1]{\tikz[baseline=-0.62ex]\fill[#1] (0,0) circle (2.5pt);\hspace{3pt}}
\newcommand{\crow}[4]{\csname CS@#1\endcsname & #2 & \textcolor{inksoft}{#3} & \textcolor{inksoft}{#4}\\}
\newcommand{\leanquoteline}[2]{\par\noindent\hangindent=\dimexpr#1\fontdimen2\font+2.2em\relax\hangafter=1
  \hspace*{#1\fontdimen2\font}#2}
\newcommand{\quoted}[1]{\par\vspace{3pt}{\ttfamily\fontsize{7.5}{10.4}\selectfont\color{leancol}\raggedright
  \setlength{\parskip}{0pt}\leanquote{#1}\par}\vspace{4pt}}
\newcommand{\bx}[3]{{\sffamily\bfseries\fontsize{7.3}{9.2}\selectfont\color{#1!82!black}#2}\\[1.5pt]#3}
\newcommand{\cs}[2]{\expandafter\newcommand\csname CS@#1\endcsname{#2}}

% theorems
\newtheoremstyle{plainx}{7pt plus 2pt}{7pt plus 2pt}{\itshape}{}{\bfseries}{.}{0.5em}{}
\newtheoremstyle{uprightx}{7pt plus 2pt}{7pt plus 2pt}{}{}{\bfseries}{.}{0.5em}{}
\theoremstyle{plainx}
\newtheorem{theorem}{Theorem}[section]
\newtheorem{lemma}[theorem]{Lemma}
\newtheorem{proposition}[theorem]{Proposition}
\newtheorem{corollary}[theorem]{Corollary}
\theoremstyle{uprightx}
\newtheorem{definition}[theorem]{Definition}
\newtheorem{setting}[theorem]{Setting}          % hypotheses shared by several statements; each of them cites it by number
\newtheorem{remark}[theorem]{Remark}
\newtheorem{example}[theorem]{Example}
\renewcommand{\qedsymbol}{{\setlength{\fboxsep}{0pt}\setlength{\fboxrule}{0.5pt}\fbox{\rule{0pt}{0.46em}\hspace{0.46em}}}}
% what holds throughout the note: the ambient objects only (implicit and instance arguments), never a hypothesis
\newenvironment{conventions}{\par\addvspace{5pt}\noindent\textbf{Conventions.}\enspace\ignorespaces}{\par\addvspace{3pt}}
\makeatletter
\renewenvironment{proof}[1][Proof]{\par\pushQED{\qed}\normalfont\topsep0pt\partopsep0pt
  \trivlist\item[\hskip\labelsep\itshape #1.]\ignorespaces}{\popQED\endtrivlist\@endpefalse}
\makeatother

% where each declaration of the concordance is treated: \cs{key}{place}, the place naming a \label or a lettered item
\cs{Frame}{Definition~\ref{d:frame}}
\cs{IsBessel}{Definition~\ref{d:bestype}}
\cs{besselGraph}{Definition~\ref{d:graph}}
\cs{besselGraphChart}{Definition~\ref{d:chart}}
\cs{besselGraphMeasure}{Definition~\ref{d:gmeasure}}
\cs{besselGraphMeasurefiniteOnCompacts}{Lemma~\ref{l:gmeasure}}
\cs{besselGraphMeasureisMulRightInvariant}{Lemma~\ref{l:gmeasure}}
\cs{besselGraphfstmembesselS}{Lemma~\ref{l:fst}}
\cs{lintegralunfoldbesselGraphMeasure}{Lemma~\ref{l:unfold}}
\cs{besselGraphtracemulsubeqzero}{Lemma~\ref{l:gtrace}}
\cs{besselPairHom}{Definition~\ref{d:pairhom}}
\cs{HasPeakAt}{Definition~\ref{d:peak}}
\cs{existsisTwoFactorbesselPairHom}{Theorem 4}
\cs{isProperMapfstbesselPairHom}{Lemma~\ref{l:proper}}
\cs{existsisBesselFunction}{Theorem 5}
\cs{gapWeightedNormSq}{Definition~\ref{d:gw}}
\cs{sumnormliecoordProjsq}{Lemma~\ref{l:levelsums}}
\cs{normweightMinorssqleexp}{Lemma~\ref{l:minorsK}}
\cs{dsCoeffgaussianshape}{Theorem 1}
\cs{existsdetcompresseqexplogOneAdd}{Proposition~\ref{p:compress}}
\cs{existsprincipalMinoreqexp}{Proposition~\ref{p:minors}}
\cs{normpoincareperiodsuble}{Theorem 6}
\cs{laplacemethod}{Theorem 2}
\cs{existsisTwoFactorrestriction}{Theorem 3}
\cs{phase}{Definition~\ref{d:phase}}
\cs{IsSteep}{Definition~\ref{d:steep}}
\cs{subOnePair}{Definition~\ref{d:pair}}
\cs{laplaceIntegrand}{Definition~\ref{d:theta}}
\cs{IsTwoFactor}{Definition~\ref{d:twofactor}}
\cs{existsunitarysuboneeqinclMatrix}{Lemma~\ref{l:block}}
\cs{normsqaddretracele}{Lemma~\ref{l:trace}}
\cs{phaselogOneAdd}{Lemma~\ref{l:phaselog}}
\cs{dsCoeffmuleqofzeroweightblock}{Lemma~\ref{l:zeroblock}}
\cs{normdsCoeffpowfourle}{Lemma~\ref{l:fourth}}
\cs{existsnearone}{Proposition~\ref{p:nearone}}
% ---------------------------------------------------------------------------- macros of this note (written by hand)
\newcommand{\CC}{\mathbb{C}}
\newcommand{\RR}{\mathbb{R}}
\newcommand{\NN}{\mathbb{N}}
\newcommand{\ZZ}{\mathbb{Z}}
\newcommand{\QQ}{\mathbb{Q}}
\newcommand{\fs}{\mathfrak{s}}
\newcommand{\fp}{\mathsf{p}}
\newcommand{\abs}[1]{\lvert #1\rvert}
\newcommand{\bigabs}[1]{\bigl\lvert #1\bigr\rvert}
\newcommand{\Bigabs}[1]{\Bigl\lvert #1\Bigr\rvert}
\newcommand{\hs}[1]{\lVert #1\rVert_{\mathrm{HS}}}
\newcommand{\op}[1]{\lVert #1\rVert_{\mathrm{op}}}
\newcommand{\nrm}[1]{\lVert #1\rVert}
\newcommand{\tr}{\operatorname{tr}}
\newcommand{\re}{\operatorname{Re}}
\newcommand{\vol}{\operatorname{vol}}
\newcommand{\GL}{\operatorname{GL}}
\newcommand{\Mat}[1]{\mathrm{M}_{#1}(\CC)}
\newcommand{\comm}[2]{[#1,#2]}
\newcommand{\ip}[2]{\langle #1,#2\rangle}
% unitary groups and matrices (background library)
\newcommand{\Us}[1]{\mathrm{U}(#1)}                      % USign ε
\newcommand{\Ks}[1]{\mathsf{K}(#1)}                      % maxCompactU ε
\newcommand{\bd}[2]{\operatorname{bd}_{#1}(#2)}          % blockDiagPart ε g
\newcommand{\inc}[1]{\operatorname{in}_{#1}}             % inclMatrix ℂ σ
\newcommand{\lgo}[1]{\operatorname{lg}(#1)}              % logOneAdd X
\newcommand{\cpr}[1]{p_{#1}}                             % coordProj p
\newcommand{\pmn}[2]{\Delta_{#1}(#2)}                    % principalMinor y p
\newcommand{\dg}[1]{\operatorname{diag}(#1)}             % intDiagonal w, signMatrix ε
% parameters and coefficients (background library)
\newcommand{\lowk}[1]{\Lambda(#1)}                       % lowestKType l
\newcommand{\ups}[1]{\Upsilon(#1)}                       % paramAddRho l
\newcommand{\dsc}[2]{\psi_{#1}(#2)}                      % dsCoeff l g
\newcommand{\fj}[3]{\operatorname{FJ}_{#1}(#2,#3)}       % flenstedJensen l g z
\newcommand{\wm}[2]{\mathfrak{m}_{#1}(#2)}               % weightMinors w g
\newcommand{\ktn}[2]{\lVert\mathfrak{m}_{#2}\rVert^{2}_{\mathsf{K}(#1)}}   % kTypeNormSq (signs ε, weight w)
\newcommand{\fdeg}[1]{\deg(#1)}                          % formalDegree l
% the shape and its groups (background library)
\newcommand{\rhs}[1]{\varrho(#1)}                        % Shape.ϱ
\newcommand{\kps}[1]{\kappa(#1)}                         % Shape.κ
\newcommand{\nsh}[1]{n(#1)}                              % Shape.n
\newcommand{\rsh}[1]{r(#1)}                              % Shape.r
\newcommand{\Ind}[1]{\operatorname{Ind}(#1)}             % Shape.BesselIndex
\newcommand{\Jone}[1]{J_{1}(#1)}                         % Shape.J₁
\newcommand{\Gone}[1]{G_{1}(#1)}                         % Shape.G₁
\newcommand{\Uu}[1]{\mathcal{U}(#1)}                     % besselU (complex form)
\newcommand{\Hh}[1]{\mathcal{H}(#1)}                     % besselH
\newcommand{\Qq}[1]{\mathcal{Q}(#1)}                     % besselS
\newcommand{\emb}[2]{\operatorname{emb}_{#1}(#2)}        % embedH h
\newcommand{\sd}[1]{\operatorname{sd}(#1)}               % simpleRootSum
\newcommand{\lev}[1]{\operatorname{lev}(#1)}             % FI.level
\newcommand{\hU}[1]{\eta_{#1}}                           % Shape.besselHaar
\newcommand{\hq}[1]{\dot\eta_{#1}}                       % besselQuotMeasure
\newcommand{\bsz}[1]{\operatorname{size}_{\mathcal{Q}}(#1)}  % besselSize
\newcommand{\peak}[5]{\operatorname{Peak}_{#1,#2}(#3;#4,#5)} % HasPeakAt l rep W x P
\newcommand{\lf}[1]{{\leanfont #1}}
\newcommand{\inp}[1]{\textup{(I#1)}}
\begin{document}
\color{ink}

{\raggedright
{\sffamily\bfseries\fontsize{7.2}{9}\selectfont\color{leancol}\addfontfeatures{LetterSpace=8}READBACK NOTE 46\par}
\vspace{5pt}
{\fontsize{21.5}{25.5}\selectfont\bfseries Discrete series coefficients near the identity, Laplace's method for two coefficients, and Bessel functions with a peak\par}
\vspace{5pt}
{\fontsize{12.4}{16}\selectfont\itshape\color{inksoft} One unit of the Lean artifact \ArtifactName, read back into
mathematics\par}}
\vspace{8pt}
{\color{rulegrey}\hrule height 0.6pt}
\vspace{7pt}
\thispagestyle{empty}


\textbf{Summary.} For signs $\varepsilon$ on a finite set, $\Us{\varepsilon}$ is a unitary group $\mathrm{U}(p,q)$, and
Flensted-Jensen's function $\dsc{l}{\,\cdot\,}$ is the matrix coefficient of the lowest $K$-type of the discrete
series whose doubled Harish-Chandra parameter is the integer vector~$l$. This note proves five things about it.
Near the identity $\dsc{l}{g}$ is an explicit phase, $\exp\tr(\dg{\lowk{l}}\lgo{g-1})$, times a factor
$\kappa_{\mathrm{G}}\in[0,1]$, up to a small relative error (Theorem~1); the proof expands compressed determinants to second
order. For an integrand made of two such coefficients on a group~$\mathfrak{g}$, with phases that cancel, a version of
Laplace's method bounds the volume of a unit ball of~$\mathfrak{g}$ by a power of a parameter times the absolute value of its integral
(Theorem~2). The hypotheses of Theorem~2 hold in two cases, both built from Lagrange interpolation: a restriction
from $\mathrm{U}(p,q)$ to a conjugate of itself inside a group one rank larger (Theorem~3), and the Bessel subgroup of
a unitary group at a point of controlled size (Theorem~4). In the second case the integral is the value at
that point of a Bessel function, which therefore has a peak there (Theorem~5). Last, the period of an automorphic form
against a Poincaré series of deep level is its main term up to an explicit error (Theorem~6).

\statement{leancol}{\fontsize{9.6}{13.2}\selectfont\textbf{What this note is.} The artifact \ArtifactName{} is a
body of Lean code: the compact edition of a formal proof, due to L.~Alpöge with the AI system Claude, of the
Riemann hypothesis and of its generalisation to Dirichlet $L$-functions. That headline is a claim. The authors
report that every file compiles with Lean v4.35.0 against a pinned version of the library Mathlib and that Lean's
kernel accepted the proofs; a replay of the build is recorded in the handout's record of checks when it reports.
No mathematician has read the argument through, and nothing here says that the Riemann hypothesis is proved. This
note treats one unit of the artifact, the folder on discrete series: all \UnitInNote{} named declarations of its
\UnitMembers{} (the other two are unnamed instances, stated in Lemma~\ref{l:gmeasure}). The results of the unit are
proved here in full from the Inputs of §\ref{s:inputs}: results of the background library of the artifact
(Flensted-Jensen's formula and its invariance, decay of the coefficient, estimates for Haar measures, Schur
orthogonality, Witt's theorem, unfolding), quoted with their owners and, where a reference is confirmed, with
it; and one result of another unit of the argument. Every statement was compared with its Lean text, argument by
argument; what Lean's kernel checked is the formal text, and what is written here is a reading of it. The Lean
text read is \PinText; Appendix~A is the concordance of \ConcDecls{} declarations.}

\sect{The statements}{s:intro}
The notions used in the six statements are defined in §\ref{s:setting}; here is what a reader needs to parse them.
For signs \m{\varepsilon\colon\Omega\to\{\pm1\}} on a finite set, $\Us{\varepsilon}$ is the group of the complex
matrices $g$ with \m{g^{*}\dg{\varepsilon}g=\dg{\varepsilon}}. The function $\dsc{l}{\,\cdot\,}$ on $\Us{\varepsilon}$
is Flensted-Jensen's function of the integer vector~$l$ (Definition~\ref{d:coeff}), $\lowk{l}$ is the highest weight
of the lowest $K$-type (Definition~\ref{d:param}), and $\lgo{X}$ is the logarithm of \m{1+X}
(Definition~\ref{d:lg}). Two norms of matrices occur: the Hilbert--Schmidt norm $\hs{\,\cdot\,}$ and the
operator norm $\op{\,\cdot\,}$. A vector~$l$ \kw{has large gap~$A$ for~$\varepsilon$} if its entries are distinct,
of the parity of \m{\abs{\Omega}-1}, and differ by at least \m{2(500\abs{\Omega}^{3}+A)} across the two signs
(Definition~\ref{d:gap}).

\statement{arith}{\phantomsection\label{t:1}\textbf{Theorem 1} (Gaussian shape near the identity).\enspace{\itshape Let
$\Omega$ be a non-empty finite set, \m{\varepsilon\colon\Omega\to\{\pm1\}} signs, \m{l\in\ZZ^{\Omega}}, let
$\delta$ and $T$ be real numbers, $L$ and $A$ natural numbers, and \m{g\in\Us{\varepsilon}}. Suppose that $l$ has
large gap~$A$ for~$\varepsilon$, that \m{\abs{l_i}\le L} for every \m{i\in\Omega}, and that
\[\begin{gathered}\hs{g-1}\le\delta,\qquad 1\le T,\qquad \abs{\Omega}^{2}\le T,\\
\abs{\Omega}^{2}\log(\abs{\Omega}+L+1)\le T,\qquad \delta\,T\le 10^{-12}.\end{gathered}\]
Then there is a real number \m{\kappa_{\mathrm{G}}\in[0,1]} with
\[\Bigabs{\frac{\dsc{l}{g}}{\exp\tr\bigl(\dg{\lowk{l}}\,\lgo{g-1}\bigr)}-\kappa_{\mathrm{G}}}\ \le\ \frac{\kappa_{\mathrm{G}}}{10}+e^{-100\,T}.\]}}
\leanline{\lean{dsCoeffgaussianshape}}

The proof gives $\kappa_{\mathrm{G}}=\exp\bigl(-\tfrac14\sum_{i,i'}\abs{\ups{l}_i-\ups{l}_{i'}}\,\abs{(g-1)_{ii'}}^{2}\bigr)$, a
Gaussian in the entries of \m{g-1} whose widths are the differences of the entries of
\m{\ups{l}=\lambda+\rho} (Definition~\ref{d:param}). The factor $\kappa_{\mathrm{G}}$ (the index G for ``Gaussian'') is small as soon as $g$ leaves the
diagonal torus by more than the reciprocal square roots of these differences; on the torus the coefficient is the
phase alone. The error is uniform in~$l$: only the size~$L$ and the rank enter, through~$T$.

The next statement concerns a product of two coefficients. For a homomorphism
$j=(j_1,j_2)\colon \mathfrak{g}\to\Us{\varepsilon}\times\Us{\varepsilon'}$ and a function \m{\phi\colon \mathfrak{g}\to\CC} put
\[\lapint{l}{l'}{j}{\phi}{\sigma}=\dsc{l}{j_1(\sigma)}\;\overline{\dsc{l'}{j_2(\sigma)}}\;e^{-\phi(\sigma)}\qquad(\sigma\in \mathfrak{g}).\]
The predicate \m{\twof{l}{l'}{j}{e'}{B}{e}{\phi}} (Definition~\ref{d:twofactor}) says, in short: the weights
$\lowk{l}$ vanish on a block $e(\Xi)$ of positive coordinates and are steep, between $-3B$ and $-B$, on the other
coordinates, whose signs alternate; $\phi$ is additive, purely imaginary, and equal to
$\tr(\dg{\lowk{l}}(j_1-1))-\tr(\dg{\lowk{l'}}(j_2-1))$; and $j(\mathfrak{g})$ contains a copy of the unitary group of
$\CC^{\Xi}$, embedded in both factors.

\statement{arith}{\phantomsection\label{t:2}\textbf{Theorem 2} (Laplace's method for two coefficients).\enspace{\itshape Let
$\Omega$ and $\Omega'$ be non-empty finite sets, $\Xi$ a finite set, and $\mathfrak{g}$ a group with a topology, with its
Borel $\sigma$-algebra, for which all left and right translations are measurable. Let $\nu$ be a measure on~$\mathfrak{g}$
that is invariant under right translations and finite on compact sets. Let \m{\varepsilon\colon\Omega\to\{\pm1\}} and
\m{\varepsilon'\colon\Omega'\to\{\pm1\}} be signs, \m{l\in\ZZ^{\Omega}}, \m{l'\in\ZZ^{\Omega'}},
\m{e'\colon\Xi\to\Omega'} a map, $B$ and $\beta$ real numbers, $L$, $M$ and $A$ natural numbers,
\m{e\colon\Xi\sqcup\{0,\dots,M-1\}\to\Omega} a bijection,
$j=(j_1,j_2)\colon \mathfrak{g}\to\Us{\varepsilon}\times\Us{\varepsilon'}$ a homomorphism of groups and
\m{\phi\colon \mathfrak{g}\to\CC} a function. Suppose that
\begin{enumerate}[label=\textup{(\alph*)},leftmargin=2.3em,itemsep=1pt]
\item \m{\twof{l}{l'}{j}{e'}{B}{e}{\phi}} holds;
\item $l$ has large gap~$A$ for~$\varepsilon$, and $l'$ has large gap~$A$ for~$\varepsilon'$;
\item \m{1\le B}; \m{\abs{l_i}\le L} for every \m{i\in\Omega}; \m{\abs{l'_i}\le L} for every \m{i\in\Omega'};
\item \m{10^{12}\le\beta}, \m{\abs{\Omega}+\abs{\Omega'}+L\le\beta} and
  $10^{30}\bigl((\abs{\Omega}+\abs{\Omega'})^{2}\log\beta\bigr)^{3}\le A$;
\item the map \m{j_1\colon \mathfrak{g}\to\Us{\varepsilon}} is proper, the map \m{j_2\colon \mathfrak{g}\to\Us{\varepsilon'}} is continuous, and
  \m{\op{j_2(\sigma)}\le\beta+1} for every \m{\sigma\in \mathfrak{g}} with \m{\op{j_1(\sigma)-1}\le1}.
\end{enumerate}
Then the set on the left has finite measure, $\lapint{l}{l'}{j}{\phi}{\,\cdot\,}$ is $\nu$-integrable, and
\[\nu\bigl\{\sigma\in \mathfrak{g}:\ \op{j_1(\sigma)-1}\le1\bigr\}\ \le\ \beta^{\,20(\abs{\Omega}+\abs{\Omega'})^{2}}\cdot
\Bigabs{\int_{\mathfrak{g}}\lapint{l}{l'}{j}{\phi}{\sigma}\,d\nu(\sigma)}.\]}}
\leanline{\lean{laplacemethod}}
The formal statement is the inequality for the real-valued measure of the set and the integral of the function; it
does not assert the finiteness of the measure or the integrability, which are proved in Step~3 of the proof below and
make the inequality meaningful as written here.

So the integral of the oscillating integrand is not much smaller than the volume of the region where the integrand
is not small. The reason is that the integrand has non-negative real part up to a tiny error, everywhere: far from the
block-diagonal matrices both coefficients are tiny, and near them a translation by the embedded unitary group
of~$\CC^{\Xi}$, which changes neither the integrand nor the measure, moves the point close to~$1$, where Theorem~1
applies to both factors and the two phases cancel against~$e^{-\phi}$.

\statement{arith}{\phantomsection\label{t:3}\textbf{Theorem 3} (the restriction to a conjugate subgroup).\enspace{\itshape Let
$\Omega$ and $\Xi$ be finite sets, \m{\varepsilon\colon\Omega\to\{\pm1\}} signs, $B$ and $H$ real numbers, $M$ a
natural number, \m{l\in\ZZ^{\Omega}} and \m{l'\in\ZZ^{\Omega\sqcup\{*\}}}, and let
\m{e\colon\Xi\sqcup\{0,\dots,M-1\}\to\Omega} and $\tilde e\colon\Xi\sqcup\{0,\dots,M\}\to\Omega\sqcup\{*\}$ be
bijections. Write \m{\varepsilon\oplus1} for the signs on \m{\Omega\sqcup\{*\}} that extend $\varepsilon$ by $+1$
at~$*$. Suppose that $M$ is even, that \m{\steep{l}{B}{e}} holds for~$\varepsilon$ and \m{\steep{l'}{B}{\tilde e}}
holds for \m{\varepsilon\oplus1}, that $l'$ has gap~$H$ for \m{\varepsilon\oplus1}, and that \m{1\le B}. Then
there is \m{\omega\in\Us{\varepsilon\oplus1}} such that the homomorphism
\[j^{\omega}\colon\Us{\varepsilon}\to\Us{\varepsilon}\times\Us{\varepsilon\oplus1},\qquad
j^{\omega}(g)=\bigl(g,\ \omega^{-1}(g\oplus1)\,\omega\bigr),\]
has the following three properties:
\begin{enumerate}[label=\textup{(\roman*)},leftmargin=2.3em,itemsep=1pt]
\item $\op{\omega}\le(\abs{\Omega}+1)^{2}\,(6B)^{M+(M+1)}$;
\item $\twof{l}{-l'}{j^{\omega}}{\tilde e|_{\Xi}}{B}{e}{0}$ holds, with the function $0$ as phase;
\item for every \m{g\in\Us{\varepsilon}}, the second component of the pair
  $\bigl(j^{\omega}_1(g)-1,\ j^{\omega}_2(g)-1\bigr)$ is \m{\omega^{-1}\bigl((g-1)\oplus0\bigr)\omega}.
\end{enumerate}}}
\leanline{\lean{existsisTwoFactorrestriction}}

Here \m{g\oplus1} and \m{(g-1)\oplus0} are the block matrices on \m{\Omega\sqcup\{*\}} with the entry $1$,
respectively~$0$, at $(*,*)$. With Theorem~2 this gives a lower bound for the integral over $\Us{\varepsilon}$ of
$\dsc{l}{g}$ against the conjugate of $\dsc{-l'}{\omega^{-1}(g\oplus1)\omega}$: a period for the restriction from
$\mathrm{U}(p+1,q)$ to $\mathrm{U}(p,q)$, at a position~$\omega$ of explicit size.

Theorems 4 and 5 concern a \kw{shape} $\fs$ (§\ref{s:shape}): a hermitian space $\CC^{\Ind{\fs}}$ with $\rhs{\fs}$
hyperbolic planes, a positive definite part with basis indexed by $\kps{\fs}$ and a negative vector, of rank
$\rsh{\fs}$, with unitary group $\Gone{\fs}$; its unipotent group $\Uu{\fs}$ with the character
\m{u\mapsto e^{\sd{u}-\overline{\sd{u}}}}; the compact group $\Us{1}$ of $\CC^{\kps{\fs}}$, embedded by $\emb{\fs}{\,\cdot\,}$;
the group $\bgr{\fs}$ of the pairs \m{(g,h)} with \m{g\in\emb{\fs}{h}\,\Uu{\fs}}, with its measure $\bms{\fs}$; and
the homomorphism $\bph{x}{\,\cdot\,}$ from $\bgr{\fs}$ to \m{\Us{\varepsilon}\times\Us{1}}, which conjugates the
first component by a point~$x$ and by a frame~$\fp$ (Definitions \ref{d:frame} to~\ref{d:pairhom}).

\statement{arith}{\phantomsection\label{t:4}\textbf{Theorem 4} (the Bessel subgroup at a point of controlled size).\enspace{\itshape Let
$\fs$ be a shape, \m{\varepsilon\colon\Ind{\fs}\to\{\pm1\}} signs, \m{\fp\in\frm{\fs}{\varepsilon}} a frame, $\Xi$ a
finite set, $B$ a real number, \m{l\in\ZZ^{\Ind{\fs}}}, and
$e\colon\Xi\sqcup\{0,\dots,2\rhs{\fs}+1\}\to\Ind{\fs}$ a bijection. Let
\m{\tilde e\colon\Xi\sqcup\{*\}\to\kps{\fs}} be a bijection, \m{l'\in\ZZ^{\kps{\fs}}}, $a$ a natural number and $b$
a real number. Suppose that \m{\steep{l}{B}{e}} holds for~$\varepsilon$; that, for the signs~$1$ on $\kps{\fs}$,
\m{\lowk{l'}_i=a} if \m{i=\tilde e(*)} and \m{\lowk{l'}_i=0} otherwise; and that
\[1\le B,\qquad 3B+a\le 2\,b^{2},\qquad 8\,(\rhs{\fs}+1)\le b.\]
Then there is a point \m{x\in\Gone{\fs}} with \m{\bsz{x}\le b^{\,5\,\nsh{\fs}}} such that
\[\twof{l}{l'}{\bph{x}{\,\cdot\,}}{\tilde e|_{\Xi}}{B}{e}{\phi_{\fs}}\quad\text{holds, where}\quad
\phi_{\fs}(\sigma)=\sd{\sigma_1}-\overline{\sd{\sigma_1}}\]
for \m{\sigma=(\sigma_1,\sigma_2)\in\bgr{\fs}}.}}
\leanline{\lean{existsisTwoFactorbesselPairHom}}

Here \m{\bsz{x}=\inf_{q\in\Qq{\fs}}\op{q\,x}+1} measures the distance of~$x$ from the Bessel subgroup $\Qq{\fs}$. The
predicate $\peak{l}{\mathrm{rep}}{W}{x}{P}$ of the next statement (Definition~\ref{d:peak}) says that $W$ is a
Bessel function for a representation $\mathrm{rep}$ of $\Us{1}$, reproduced by the coefficient of~$l$, non-zero
at~$x$, and that two integrals over the Bessel subgroup are bounded in terms of \m{P\,\nrm{W(x)}}.

\statement{arith}{\phantomsection\label{t:5}\textbf{Theorem 5} (a Bessel function with a peak).\enspace{\itshape Let $\fs$
be a shape, $I$ a finite set, \m{\varepsilon\colon\Ind{\fs}\to\{\pm1\}} signs and \m{\fp\in\frm{\fs}{\varepsilon}} a
frame. Let \m{x\in\Gone{\fs}}, \m{v_0\in\CC^{I}}, let $A$ and $\mathfrak{w}$ be natural numbers, $b$ a real number,
\m{l\in\ZZ^{\Ind{\fs}}}, \m{l'\in\ZZ^{\kps{\fs}}}, and let $\mathrm{rep}$ be a homomorphism from the group
$\Us{1}$ of the unitary matrices of $\CC^{\kps{\fs}}$ to the unitary \m{I\times I} matrices. Suppose that
$\mathrm{rep}$ is continuous; that \m{\ip{v_0}{\mathrm{rep}(h)\,v_0}=\dsc{l'}{h}} for every \m{h\in\Us{1}}; that $l$
has large gap~$A$ for~$\varepsilon$; that \m{16\le b}; and that, with $\phi_{\fs}(\sigma)=\sd{\sigma_1}-\overline{\sd{\sigma_1}}$,
\[\bms{\fs}\bigl\{\sigma\in\bgr{\fs}:\ \op{\bph{x}{\sigma}_1-1}\le1\bigr\}\ \le\ b^{\,\mathfrak{w}}\cdot
\Bigabs{\int_{\bgr{\fs}}\lapint{l}{l'}{\bph{x}{\,\cdot\,}}{\phi_{\fs}}{\sigma}\,d\bms{\fs}(\sigma)}.\]
Then there is a function \m{W\colon\Gone{\fs}\to\CC^{I}} such that
$\peak{l}{\mathrm{rep}}{W}{x}{b^{\,8\rsh{\fs}^{2}+1+\mathfrak{w}}}$ holds.}}
\leanline{\lean{existsisBesselFunction}}

The displayed hypothesis is the conclusion of Theorem~2 for the data of Theorem~4, with \m{b^{\mathfrak{w}}} in place of the
power of~$\beta$. The function is
\[W(y)=\int_{\bgr{\fs}}\dsc{l}{\fp^{-1}x^{-1}\sigma_1y\,\fp}\,e^{-\phi_{\fs}(\sigma)}\,\mathrm{rep}(\sigma_2)^{-1}v_0\,d\bms{\fs}(\sigma),\]
and \m{\ip{v_0}{W(x)}} is the integral in the hypothesis.

The last statement is adelic; §\ref{s:adelic} has the definitions. For a CM field~$E$ with $f_E$ infinite places
and signs~$\varepsilon$ on~$\Omega$, the group $G(E,\varepsilon)=G_\infty(E,\varepsilon)\times G_{\mathrm f}(E,\Omega)$ is
the adelic unitary group in standard coordinates, with $G_\infty(E,\varepsilon)=\prod_v\Us{\varepsilon}$ over the
infinite places~$v$. A hermitian space with a point~$V$ gives the discrete subgroup $\Gamma(V)$ of
$G(E,\varepsilon)$ and a measure $\varpi_V$ that computes integrals over \m{\Gamma(V)\backslash G(E,\varepsilon)};
\m{V\cdot x} is $V$ with its point moved by~$x$, $\lambda_1(V)$ is the length of a shortest non-zero lattice vector
of~$V$, cut off at~$1$, \m{V\oplus\langle1\rangle} is the sum of $V$ and a positive line, $\mathsf{K}_\Omega(a_0)$ is the
principal congruence subgroup of level~$a_0$, and \m{g\oplus1} is the block-diagonal embedding into the group of
\m{\Omega\sqcup\{*\}}.

\statement{arith}{\phantomsection\label{t:6}\textbf{Theorem 6} (the period of a Poincaré series of deep level).\enspace{\itshape Let
$E$ be a CM field with a CM type that contains a primitive root of unity of order~$5040$, let $\Omega$ be a finite
set, \m{\varepsilon\colon\Omega\to\{\pm1\}} signs, and $V$ a hermitian space with a point over~$E$ for~$\varepsilon$. Let
\m{F\colon G(E,\varepsilon)\to\CC} be a function, $C$, $\rho$ and $Q$ real numbers, $a_0$ an integer of~$E$, $A$ a
natural number, \m{l=(l_v)_v} a family of vectors \m{l_v\in\ZZ^{\Omega\sqcup\{*\}}} indexed by the infinite
places~$v$ of~$E$, \m{x\in G(E,\varepsilon)}, and $\omega,\alpha\in G_\infty(E,\varepsilon\oplus1)$. Suppose that
\begin{enumerate}[label=\textup{(\alph*)},leftmargin=2.3em,itemsep=1pt]
\item $F$ is automorphic for~$V$, and \m{\abs{F(g)}\le C} for every \m{g\in G(E,\varepsilon)};
\item \m{a_0\ne0}, and $F$ has level $\mathsf{K}_\Omega(a_0)$;
\item \m{\rho\le\abs{a_0}_v} for every infinite place~$v$, and \m{0\le\rho};
\item \m{2f_E\,(\abs{\Omega}+1)+2\le A} and \m{8\,(\abs{\Omega}+1)^{2}<A};
\item every $l_v$ has large gap~$2A$ for \m{\varepsilon\oplus1};
\item \m{Q^{2}\le\rho\cdot\lambda_1(V\cdot x)}, \m{f_E\le Q} and \m{10\le Q}.
\end{enumerate}
Put $\Psi(g)=\prod_v\overline{\dsc{l_v}{\alpha_v^{-1}g_v\,\omega_v}}$ for \m{g\in G_\infty(E,\varepsilon\oplus1)}, and
\[L_{*}=\prod_v\Bigl(\bigl(2\,(\op{\alpha_v}\op{\omega_v})^{2}\bigr)^{A}\cdot16^{A}\,
\vol\bigl\{r\in\Us{\varepsilon}:\op{r-1}\le1\bigr\}\Bigr).\]
Then \m{r\mapsto\Psi(r\oplus1)} is integrable on $G_\infty(E,\varepsilon)$, and, with
$X_\gamma(y)=(x\oplus1)^{-1}\gamma\,(y\oplus1)$ for \m{\gamma\in\Gamma(V\oplus\langle1\rangle)} and
\m{y\in G(E,\varepsilon)},
\begin{multline*}
\Bigabs{\int\sum_{\gamma\in\Gamma(V\oplus\langle1\rangle)}\Psi\bigl(X_\gamma(y)_\infty\bigr)\,
\mathbf{1}_{\mathsf{K}_{\Omega\sqcup\{*\}}(a_0)}\bigl(X_\gamma(y)_{\mathrm f}\bigr)\,F(y)\,d\varpi_V(y)\\
-\vol\bigl(\mathsf{K}_\Omega(a_0)\bigr)\int_{G_\infty(E,\varepsilon)}\Psi(r\oplus1)\,F\bigl(x\cdot(r,1)\bigr)\,dr}
\ \le\ \vol\bigl(\mathsf{K}_\Omega(a_0)\bigr)\cdot L_{*}\cdot C\cdot\frac{10}{Q}.
\end{multline*}}}
\leanline{\lean{normpoincareperiodsuble}}

The sum over~$\gamma$ is, as a function of~$y$, the restriction to $G(E,\varepsilon)$ of a Poincaré series on the
group of \m{V\oplus\langle1\rangle}, made from the test function \m{\Psi\otimes\mathbf{1}_{\mathsf{K}(a_0)}} and translated
by~$x$. The main term is the contribution of the classes \m{\gamma\in\Gamma(V)\oplus1}. All other classes move the
vector of the added line to rational vectors that are congruent modulo~$a_0$, hence far apart at the infinite places
when $a_0$ is large there and $x$ is not near the cusp; the decay of the coefficients then makes their total at most
the stated error.

\sect{The setting}{s:setting}
\begin{conventions}\phantomsection\label{c:conventions}
$\Omega$, $\Omega'$ and $\Xi$ denote a finite set each (in the formal text: a type with finitely many elements),
and $I$ a finite set that indexes the coordinates of a vector space~$\CC^{I}$. A family of signs on~$\Omega$ is a
map $\varepsilon\colon\Omega\to\{\pm1\}$; $\varepsilon'$ is a family of signs on~$\Omega'$, and $1$ also denotes
the constant family~$+1$. The letter $\mathfrak{g}$ denotes a group, $\fs$ a shape (§\ref{s:shape}) and $\fp$ a frame of it.
All matrices have complex entries; for index sets $\Omega_1,\Omega_2$ the entries of an
\m{\Omega_1\times\Omega_2} matrix~$X$ are $X_{ii'}$, and $X^{*}$ is its conjugate transpose. An integer vector
\m{l\in\ZZ^{\Omega}} is written with entries~$l_i$. Every statement introduces the objects it is about, also those
that the formal text takes as implicit arguments; the instance arguments of the formal text (finiteness of the
index sets, the group structure, the topology and the $\sigma$-algebra of~$\mathfrak{g}$, the frame) are the structures named
here and in the statements. Integrals of complex or vector valued functions are Bochner integrals with the
convention of the library: the integral of a function that is not integrable is~$0$. A sum over an infinite set is
the unconditional sum when it exists and $0$ otherwise. For a measure~$\nu$, the measure of a set that occurs in
an inequality between real numbers is read as $0$ when it is infinite; in this note every such set is shown to
have finite measure. The letters $D$, $P$, $d$, $f$, $r$, $s$, $w$, $x$, $\nu$ and $\mathsf K$ of the handout-wide Conventions~0.1, 0.3, 0.4, 0.6 and~0.7
(the data and the parameters, the prime, the number of infinite places, the rank, the complex variable, a finite place, the scale, a character of the ideles, a
compact open subgroup) are used here only for local objects, and each use says which; $\nu$ is a measure (Convention~VI.8), and $\mathsf K$ is a compact open
subgroup of an adelic group.
\end{conventions}

\subhead{Matrices and unitary groups.} These notions belong to the background library; they are defined here
because every statement uses them.

\begin{definition}[norms, unitary groups]\label{d:mat}
Let $\Omega$ be a finite set and $\varepsilon\colon\Omega\to\{\pm1\}$ signs.
\begin{enumerate}[label=\textup{(\alph*)},leftmargin=2.3em,itemsep=1pt]
\item For an \m{\Omega\times\Omega} matrix~$X$, $\hs{X}=(\sum_{i,i'}\abs{X_{ii'}}^{2})^{1/2}$ is the
  Hilbert--Schmidt (Frobenius) norm and $\op{X}$ the operator norm for the Euclidean norm of~$\CC^{\Omega}$. Both are
  submultiplicative, \m{\op{X}\le\hs{X}} and \m{\hs{XY}\le\hs{X}\op{Y}}; \m{\comm{X}{Y}=XY-YX}.
\item For an integer or sign vector $w$ on~$\Omega$, $\dg{w}$ is the diagonal matrix with entries~$w_i$.
\item $\Us{\varepsilon}$ is the group of the invertible \m{\Omega\times\Omega} matrices $g$ with
  $g^{*}\dg{\varepsilon}\,g=\dg{\varepsilon}$; it is a unitary group $\mathrm{U}(p,q)$ with $p$ signs~$+1$ and $q$
  signs~$-1$. Its subgroup \m{\Ks{\varepsilon}} consists of the $g$ that are moreover unitary matrices,
  equivalently that commute with $\dg{\varepsilon}$; it is compact, isomorphic to $\mathrm{U}(p)\times\mathrm{U}(q)$,
  and \m{dh} denotes its Haar measure of total mass one. On $\Us{\varepsilon}$ itself, $dg$ is the Haar measure
  that the library fixes; for \m{\varepsilon=1} it has total mass one.
\item For an \m{\Omega\times\Omega} matrix~$X$, the matrix \m{\bd{\varepsilon}{X}} has the entries $X_{ii'}$ with
  \m{\varepsilon_i=\varepsilon_{i'}} and zeros elsewhere; $\bd{\varepsilon}{X}=\tfrac12(X+\dg{\varepsilon}X\dg{\varepsilon})$.
  A matrix with \m{X=\bd{\varepsilon}{X}} is called block diagonal, and the pairs \m{(i,i')} with
  \m{\varepsilon_i=\varepsilon_{i'}} compact, the others non-compact.
\item For a map \m{\iota\colon\Xi\to\Omega}, $\inc{\iota}$ is the \m{\Omega\times\Xi} matrix with entry $1$ at
  \m{(\iota(\xi),\xi)} for \m{\xi\in\Xi} and $0$ elsewhere. For a \m{\Xi\times\Xi} matrix~$u$ the matrix
  $\inc{\iota}\,u\,\inc{\iota}^{*}$ is supported on \m{\iota(\Xi)\times\iota(\Xi)}; when $\iota$ is injective it is $u$
  placed in that block.
\end{enumerate}
\end{definition}

\begin{definition}[the logarithm]\label{d:lg}
For an \m{\Omega\times\Omega} matrix~$X$, $\lgo{X}=\int_0^1(1+sX)^{-1}X\,ds$, where the inverse of a singular
matrix is read as~$0$. For \m{\hs{X}\le\tfrac12} it is the logarithm of \m{1+X}: by Input~\inp{3},
\m{\det(1+X)=\exp\tr\lgo{X}}.
\end{definition}

\begin{definition}[coordinate projections, principal minors]\label{d:minor}
For a subset \m{\Pi\subseteq\Omega}, $\cpr{\Pi}$ is the diagonal matrix with entries $1$ on~$\Pi$ and $0$
elsewhere, an orthogonal projection, and for an \m{\Omega\times\Omega} matrix~$y$ the number $\pmn{\Pi}{y}$ is the
determinant of the submatrix of~$y$ with rows and columns in~$\Pi$; so \m{\pmn{\Omega}{y}=\det y} and
$\pmn{\Pi}{y}=\det\bigl(1+\cpr{\Pi}(y-1)\bigr)$. For an integer vector $w$ on~$\Omega$ and \m{\tau\in\ZZ} we write
\m{\{w\ge\tau\}} for the set of the \m{i\in\Omega} with \m{w_i\ge\tau}. (The formal text speaks of decidable
predicates on~$\Omega$ where we speak of subsets.)
\end{definition}

\subhead{Parameters and Flensted-Jensen's function.} Also from the background library.

\begin{definition}[parameters]\label{d:param}
Let $\varepsilon\colon\Omega\to\{\pm1\}$ be signs and \m{l\in\ZZ^{\Omega}}. The integer vectors $\ups{l}$ and
$\lowk{l}$ on~$\Omega$ are
\begin{align*}
\ups{l}_i&=\Bigl\lfloor\tfrac12\bigl(l_i+\#\{i':l_{i'}<l_i\}-\#\{i':l_i<l_{i'}\}\bigr)\Bigr\rfloor,\\
\lowk{l}_i&=\ups{l}_i-\#\{i':\varepsilon_{i'}=\varepsilon_i,\ l_{i'}<l_i\}+\#\{i':\varepsilon_{i'}=\varepsilon_i,\ l_i<l_{i'}\}.
\end{align*}
The vector $\lowk{l}$ depends on~$\varepsilon$, which the notation does not show; the statements say for which signs
it is taken.
\end{definition}

One reads $l$ as a doubled Harish-Chandra parameter \m{l=2\lambda} of $\Us{\varepsilon}$. If $\rho$ and $\rho_c$ are
the half sums of all roots, respectively of the compact roots, that are positive on~$\lambda$, then
\m{\ups{l}=\lambda+\rho} and \m{\lowk{l}=\lambda+\rho-2\rho_c}: the infinitesimal character shifted to an
integral vector, and the highest weight of the lowest $K$-type (the Blattner parameter). The comments of the
formal text write $\mu$ for~$\ups{l}$. Since the two counted sets have at most $\abs{\Omega}$ elements,
\begin{equation}\label{e:rhoc}
\bigabs{\ups{l}_i-\lowk{l}_i}\le\abs{\Omega}\qquad\text{for every }i\in\Omega .
\end{equation}

\begin{definition}[gap, large gap]\label{d:gap}
Let $\varepsilon\colon\Omega\to\{\pm1\}$ be signs, \m{l\in\ZZ^{\Omega}} and $H$ a real number. The vector $l$
\kw{has gap~$H$ for~$\varepsilon$} if $l$ is injective, \m{l_i+\abs{\Omega}-1} is even for every~$i$, and
\m{\abs{l_i-l_{i'}}\ge2H} whenever $\varepsilon_i\ne\varepsilon_{i'}$. For a natural number~$A$, $l$ \kw{has large
gap~$A$ for~$\varepsilon$} if it has gap \m{500\abs{\Omega}^{3}+A} for~$\varepsilon$. If $l$ has gap~$H$ then so
has~$-l$, and the half in the definition of $\ups{l}$ is exact.
\end{definition}

\begin{definition}[products of principal minors]\label{d:wm}
Let $w$ be an integer vector on~$\Omega$ and $g$ an invertible \m{\Omega\times\Omega} matrix. Choose integers
\m{c_1\le c_2} with \m{c_1\le w_i\le c_2} for all~$i$ and put
\[\wm{w}{g}=(\det g)^{c_1}\prod_{c_1<\tau\le c_2}\pmn{\{w\ge\tau\}}{g}.\]
This does not depend on $c_1$ and~$c_2$. (The library defines $\wm{w}{g}$ for all matrices as a product over all
\m{\tau\in\ZZ}; for invertible $g$ it is this number, Input~\inp{4}.) On diagonal matrices $\wm{w}{\,\cdot\,}$ is the
character \m{\prod_i g_{ii}^{w_i}}.
\end{definition}

\begin{definition}[Flensted-Jensen's function]\label{d:coeff}
Let $\varepsilon\colon\Omega\to\{\pm1\}$ be signs and \m{l\in\ZZ^{\Omega}}; take $\lowk{l}$ for~$\varepsilon$. Put
$\ktn{\varepsilon}{\lowk{l}}=\int_{\Ks{\varepsilon}}\abs{\wm{\lowk{l}}{h}}^{2}\,dh$, a positive number
(Input~\inp{8}). For \m{\Omega\times\Omega} matrices $g$ and $z$,
\[\begin{gathered}\fj{l}{g}{z}=\frac{1}{\ktn{\varepsilon}{\lowk{l}}}\int_{\Ks{\varepsilon}}
\frac{\wm{\lowk{l}}{g\,h}\;\overline{\wm{\lowk{l}}{h}}}{\wm{-\ups{l}}{h^{*}z\,h}}\;dh,\\
\dsc{l}{g}=\fj{l}{\bd{\varepsilon}{g}}{\bd{\varepsilon}{g}^{-1}g}.\end{gathered}\]
The \kw{formal degree} is $\fdeg{l}=\bigl(\int_{\Us{\varepsilon}}\abs{\dsc{l}{g}}^{2}\,dg\bigr)^{-1}$. Both
$\dsc{l}{\,\cdot\,}$ and $\fdeg{l}$ depend on~$\varepsilon$.
\end{definition}

For \m{g\in\Us{\varepsilon}} the matrix $\bd{\varepsilon}{g}$ is invertible (Lemma~\ref{l:fourth}), and for $l$ with
a gap the integrand is defined. By the header of the library's folder, $\dsc{l}{\,\cdot\,}$ restricted to
$\Us{\varepsilon}$ is the matrix coefficient of the lowest $K$-type of the discrete series with Harish-Chandra
parameter~$\lambda$, normalised by \m{\dsc{l}{1}=1}; this identification is a result of the background library that
the present note does not use, and it is the reason for the name. What the note uses are the properties listed in
§\ref{s:inputs}.

\subhead{The definitions of the unit for two coefficients.}

\begin{definition}[the gap-weighted square norm]\label{d:gw}
Let $\Omega$ be a finite set, $w$ an integer vector on~$\Omega$ and $X$ an \m{\Omega\times\Omega} matrix. Then
\[\gwn{w}{X}=\sum_{i\in\Omega}\sum_{i'\in\Omega}\abs{w_i-w_{i'}}\,\abs{X_{ii'}}^{2}.\]
\end{definition}
\leanline{\lean{gapWeightedNormSq}}

\begin{definition}[the phase]\label{d:phase}
Let $\Omega$, $\Omega'$ be finite sets, $\varepsilon$, $\varepsilon'$ signs on them, \m{l\in\ZZ^{\Omega}},
\m{l'\in\ZZ^{\Omega'}}, and let \m{x=(x_1,x_2)} be a pair of an \m{\Omega\times\Omega} matrix and an
\m{\Omega'\times\Omega'} matrix. Then
\[\phs{l}{l'}{x}=\tr\bigl(\dg{\lowk{l}}\,x_1\bigr)-\tr\bigl(\dg{\lowk{l'}}\,x_2\bigr),\]
with $\lowk{l}$ taken for~$\varepsilon$ and $\lowk{l'}$ for~$\varepsilon'$. It is linear in~$x$.
\end{definition}
\leanline{\lean{phase}}

\begin{definition}[steep weights]\label{d:steep}
Let $\Omega$ and $\Xi$ be finite sets, $\varepsilon\colon\Omega\to\{\pm1\}$ signs, \m{l\in\ZZ^{\Omega}}, $B$ a real
number, $M$ a natural number and $e\colon\Xi\sqcup\{0,\dots,M-1\}\to\Omega$ a bijection. We say that
\m{\steep{l}{B}{e}} \kw{holds for~$\varepsilon$} if
\begin{enumerate}[label=\textup{(S\arabic*)},leftmargin=2.9em,itemsep=1pt]
\item \m{\varepsilon_{e(\xi)}=1} for every \m{\xi\in\Xi}, and \m{\varepsilon_{e(n)}=(-1)^{n}} for \m{0\le n<M};
\item \m{\lowk{l}_{e(\xi)}=0} for every \m{\xi\in\Xi};
\item the map \m{n\mapsto\lowk{l}_{e(n)}} is strictly decreasing on \m{\{0,\dots,M-1\}}, and
  $-3B\le\lowk{l}_{e(n)}\le-B$ for \m{0\le n<M}.
\end{enumerate}
\end{definition}
\leanline{\lean{IsSteep}}

\begin{definition}[the pair $j-1$]\label{d:pair}
Let $\mathfrak{g}$ be a group, $\varepsilon$ and $\varepsilon'$ signs on finite sets $\Omega$ and~$\Omega'$,
$j\colon \mathfrak{g}\to\Us{\varepsilon}\times\Us{\varepsilon'}$ a homomorphism of groups with components \m{j=(j_1,j_2)}, and
\m{\sigma\in \mathfrak{g}}. The formal text abbreviates the pair of matrices $\bigl(j_1(\sigma)-1,\ j_2(\sigma)-1\bigr)$; this
note writes the pair out.
\end{definition}
\leanline{\lean{subOnePair}}

\begin{definition}[the integrand]\label{d:theta}
Let $\mathfrak{g}$ be a group, $\varepsilon$ and $\varepsilon'$ signs on finite sets $\Omega$ and~$\Omega'$, \m{l\in\ZZ^{\Omega}},
\m{l'\in\ZZ^{\Omega'}}, $j=(j_1,j_2)\colon \mathfrak{g}\to\Us{\varepsilon}\times\Us{\varepsilon'}$ a homomorphism of groups,
\m{\phi\colon \mathfrak{g}\to\CC} a function and \m{\sigma\in \mathfrak{g}}. Then
\[\lapint{l}{l'}{j}{\phi}{\sigma}=\dsc{l}{j_1(\sigma)}\cdot\overline{\dsc{l'}{j_2(\sigma)}}\cdot\exp\bigl(-\phi(\sigma)\bigr).\]
\end{definition}
\leanline{\lean{laplaceIntegrand}}

\begin{definition}[two-factor data]\label{d:twofactor}
Let $\Omega$, $\Omega'$, $\Xi$ be finite sets, $\mathfrak{g}$ a group, $\varepsilon$ and $\varepsilon'$ signs on $\Omega$ and
$\Omega'$, \m{l\in\ZZ^{\Omega}}, \m{l'\in\ZZ^{\Omega'}},
$j=(j_1,j_2)\colon \mathfrak{g}\to\Us{\varepsilon}\times\Us{\varepsilon'}$ a homomorphism of groups, \m{e'\colon\Xi\to\Omega'} a
map, $B$ a real number, $M$ a natural number, $e\colon\Xi\sqcup\{0,\dots,M-1\}\to\Omega$ a bijection and
\m{\phi\colon \mathfrak{g}\to\CC} a function. We say that $\twof{l}{l'}{j}{e'}{B}{e}{\phi}$ \kw{holds} if \m{\steep{l}{B}{e}} holds
for~$\varepsilon$ and
\begin{enumerate}[label=\textup{(T\arabic*)},leftmargin=2.9em,itemsep=1pt]
\item $\phs{l}{l'}{(j_1(\sigma)-1,\ j_2(\sigma)-1)}=\phi(\sigma)$ for every \m{\sigma\in \mathfrak{g}};
\item $\tr(j_1(\sigma)-1)=\tr(j_2(\sigma)-1)$ for every \m{\sigma\in \mathfrak{g}};
\item $\phi(\sigma\sigma')=\phi(\sigma)+\phi(\sigma')$ for all \m{\sigma,\sigma'\in \mathfrak{g}};
\item \m{\re\phi(\sigma)=0} for every \m{\sigma\in \mathfrak{g}};
\item \m{\varepsilon'_{e'(\xi)}=1} and \m{\lowk{l'}_{e'(\xi)}=0} for every \m{\xi\in\Xi};
\item \m{0\le\lowk{l'}_i} for every \m{i\in\Omega'}, and \m{\lowk{l'}_i\le3B} for every \m{i\in\Omega'} such that
  $\varepsilon'_i\ne\varepsilon'_{i'}$ for some \m{i'\in\Omega'};
\item for every unitary \m{\Xi\times\Xi} matrix~$u$ there is \m{\sigma\in \mathfrak{g}} with
  \[j_1(\sigma)-1=\inc{e|_\Xi}\,(u-1)\,\inc{e|_\Xi}^{*}\qquad\text{and}\qquad
  j_2(\sigma)-1=\inc{e'}\,(u-1)\,\inc{e'}^{*}.\]
\end{enumerate}
\end{definition}
\leanline{\lean{IsTwoFactor}}

In (T6) the second bound is asked of the indices whose sign differs from some other sign; if $\varepsilon'$ is
constant, as for a compact second factor, it is empty. Condition (T7) says that $j(\mathfrak{g})$ contains the unitary group of
$\CC^{\Xi}$, placed in the block $e(\Xi)$ of the first factor and along $e'$ in the second. No topology on~$\mathfrak{g}$ is
part of the definition.

\subsection*{\normalsize The shape and its Bessel subgroup}\label{s:shape}
\refstepcounter{subsection}
The objects of the next definition belong to the background library and to the units of Chapter~VI on the real
place; the letters are those of the notes of these units.

\begin{definition}[shape]\label{d:shape}
A \kw{shape} $\fs$ consists of a natural number $\rhs{\fs}$, a finite set $\kps{\fs}$ and an element of
$\kps{\fs}$. Its \kw{basis index set} $\Ind{\fs}$ is the disjoint union of \m{\{0,\dots,\rhs{\fs}-1\}}, of
\m{\kps{\fs}\sqcup\{z\}} and of a second copy of \m{\{0,\dots,\rhs{\fs}-1\}}; the standard basis of
$\CC^{\Ind{\fs}}$ is written
\[e_0,\dots,e_{\rhs{\fs}-1};\qquad x_i\ (i\in\kps{\fs}),\ z;\qquad e^{*}_0,\dots,e^{*}_{\rhs{\fs}-1}.\]
Put \m{\nsh{\fs}=\rhs{\fs}+1} and $\rsh{\fs}=\abs{\Ind{\fs}}=2\rhs{\fs}+\abs{\kps{\fs}}+1$.
\begin{enumerate}[label=\textup{(\alph*)},leftmargin=2.3em,itemsep=1pt]
\item $\Jone{\fs}$ is the Gram matrix of the hermitian form on $\CC^{\Ind{\fs}}$ for which $e_a$ and $e^{*}_a$ span
  a hyperbolic plane with product~$1$, the $x_i$ are orthonormal, $z$ has square~$-1$, and all other products of
  basis vectors vanish. $\Gone{\fs}$ is its unitary group, with the Haar measure $dg$ of the library.
\item The \kw{level} of a basis index is given by
  \[\lev{e_a}=a,\qquad\lev{x_i}=\lev{z}=\rhs{\fs},\qquad\lev{e^{*}_a}=2\rhs{\fs}-a.\]
  A matrix~$X$ is \kw{block upper triangular for the levels} if \m{X_{ii'}=0} whenever
  \m{\lev{i'}<\lev{i}}. The \kw{frame} is \m{f_a=e_a} for \m{a<\rhs{\fs}} and \m{f_{\rhs{\fs}}=z}.
\item $\Uu{\fs}$ is the subgroup of the \m{u\in\Gone{\fs}} with \m{u_{ii'}=\delta_{ii'}} whenever
  \m{\lev{i'}\le\lev{i}}, a unipotent group; $\sd{g}=\sum_{a<\rhs{\fs}}g_{f_a,f_{a+1}}$ is the sum of the entries
  of~$g$ just above the diagonal on the frame.
\item $\Us{1}$, for the signs $1$ on $\kps{\fs}$, is the group of the unitary matrices of $\CC^{\kps{\fs}}$, and
  $\emb{\fs}{h}$, for \m{h\in\Us{1}}, is the element of $\Gone{\fs}$ that acts as $h$ on the span of the $x_i$ and
  fixes the other basis vectors. $\Hh{\fs}$ is the image of $\emb{\fs}{\,\cdot\,}$, and $\Qq{\fs}$, the \kw{Bessel
  subgroup}, is the group \m{\Uu{\fs}\,\Hh{\fs}}; it normalizes $\Uu{\fs}$.
\item $\hU{\fs}$ is the Haar measure on $\Uu{\fs}$ that the library fixes, and $\hq{\fs}$ is the measure on
  $\Gone{\fs}$ of the library that computes integrals over \m{\Uu{\fs}\backslash\Gone{\fs}}: for measurable
  \m{f\ge0} on $\Gone{\fs}$, $\int f(g)\,dg=\iint f(u\,y)\,d\hU{\fs}(u)\,d\hq{\fs}(y)$.
\item For \m{x\in\Gone{\fs}}, $\bsz{x}=\inf_{q\in\Qq{\fs}}\op{q\,x}+1$.
\item Let $I$ be a finite set and $\mathrm{rep}$ a homomorphism from $\Us{1}$ to the unitary \m{I\times I}
  matrices. A function \m{W\colon\Gone{\fs}\to\CC^{I}} is a \kw{Bessel function for $\mathrm{rep}$} if it is bounded
  and strongly measurable, $W(u\,y)=e^{\sd{u}-\overline{\sd{u}}}\,W(y)$ for \m{u\in\Uu{\fs}}, and
  $W(\emb{\fs}{h}\,y)=\mathrm{rep}(h)\,W(y)$ for \m{h\in\Us{1}}, for all \m{y\in\Gone{\fs}}.
\end{enumerate}
\end{definition}

The signature of $\Jone{\fs}$ is $(\rhs{\fs}+\abs{\kps{\fs}},\,\rhs{\fs}+1)$. The function $\sd{\,\cdot\,}$ is
additive on $\Qq{\fs}$ and vanishes on $\Hh{\fs}$ (Input~\inp{13}), so $u\mapsto e^{\sd{u}-\overline{\sd{u}}}$ is a
unitary character of $\Uu{\fs}$, the generic character of this note. Item~(f) is a definition of the unit on
the second moment (Filling), item~(g) one of the library.

\begin{definition}[frames]\label{d:frame}
Let $\fs$ be a shape and $\varepsilon\colon\Ind{\fs}\to\{\pm1\}$ signs. The set \m{\frm{\fs}{\varepsilon}} of the
\kw{frames of $\fs$ for~$\varepsilon$} consists of the unitary \m{\Ind{\fs}\times\Ind{\fs}} matrices~$\fp$ with
$\fp^{*}\,\Jone{\fs}\,\fp=\dg{\varepsilon}$. For such $\fp$, \m{g\mapsto\fp^{-1}g\,\fp} is an isomorphism of
topological groups from $\Gone{\fs}$ to $\Us{\varepsilon}$ that preserves the Haar measures.
\end{definition}
\leanline{\lean{Frame}}

In the formal text a frame is a structure that the declarations below take as an instance argument, and the signs
are determined by it; here a statement names its frame. The columns of~$\fp$ are an orthonormal basis of
$\CC^{\Ind{\fs}}$ that is also orthogonal for $\Jone{\fs}$.

\begin{definition}[the Bessel graph]\label{d:graph}
Let $\fs$ be a shape. Then $\bgr{\fs}$ is the set of the pairs \m{(g,h)\in\Gone{\fs}\times\Us{1}}, $\Us{1}$ the
unitary group of $\CC^{\kps{\fs}}$, with \m{\emb{\fs}{h}^{-1}g\in\Uu{\fs}}. It is a subgroup, because
$\emb{\fs}{h}$ normalizes $\Uu{\fs}$. An element is written \m{\sigma=(\sigma_1,\sigma_2)}.
\end{definition}
\leanline{\lean{besselGraph}}

The first projection maps $\bgr{\fs}$ isomorphically onto \m{\Hh{\fs}\,\Uu{\fs}=\Qq{\fs}}; the graph records the
factor in $\Us{1}$, which the second coefficient of the integrand needs.

\begin{definition}[the chart]\label{d:chart}
Let $\fs$ be a shape. Then $\bch{\fs}$ is the homeomorphism from \m{\Uu{\fs}\times\Us{1}} to $\bgr{\fs}$ that sends
\m{(u,h)} to \m{(\emb{\fs}{h}\,u,\ h)}; its inverse sends \m{(g,h)} to \m{(\emb{\fs}{h}^{-1}g,\ h)}.
\end{definition}
\leanline{\lean{besselGraphChart}}

\begin{definition}[the measure of the Bessel graph]\label{d:gmeasure}
Let $\fs$ be a shape. Then $\bms{\fs}$ is the image under the chart $\bch{\fs}$ of the product of the measure
$\hU{\fs}$ on $\Uu{\fs}$ and the Haar measure of total mass one on $\Us{1}$.
\end{definition}
\leanline{\lean{besselGraphMeasure}}

\begin{definition}[the pair homomorphism]\label{d:pairhom}
Let $\fs$ be a shape, $\varepsilon\colon\Ind{\fs}\to\{\pm1\}$ signs, \m{\fp\in\frm{\fs}{\varepsilon}} a frame and
\m{x\in\Gone{\fs}}. For $\sigma=(\sigma_1,\sigma_2)\in\bgr{\fs}$,
\[\bph{x}{\sigma}=\bigl(\fp^{-1}x^{-1}\sigma_1\,x\,\fp,\ \sigma_2\bigr)\ \in\ \Us{\varepsilon}\times\Us{1}.\]
The map $\bph{x}{\,\cdot\,}$ is a homomorphism of groups. It depends on $\fs$ and on the frame~$\fp$, which the
notation does not show (they are implicit in the formal text).
\end{definition}
\leanline{\lean{besselPairHom}}

\begin{definition}[peak]\label{d:peak}
Let $\fs$ be a shape, $\varepsilon\colon\Ind{\fs}\to\{\pm1\}$ signs, \m{\fp\in\frm{\fs}{\varepsilon}} a frame, $I$ a
finite set, \m{l\in\ZZ^{\Ind{\fs}}}, $\mathrm{rep}$ a homomorphism from $\Us{1}$ (of $\CC^{\kps{\fs}}$) to the unitary
\m{I\times I} matrices, \m{W\colon\Gone{\fs}\to\CC^{I}} a function, \m{x\in\Gone{\fs}} and $P$ a real number. We say
that $\peak{l}{\mathrm{rep}}{W}{x}{P}$ \kw{holds} if $W$ is a Bessel function for $\mathrm{rep}$
(Definition~\ref{d:shape}(g)) and
\begin{enumerate}[label=\textup{(P\arabic*)},leftmargin=2.9em,itemsep=1pt]
\item $W(y)=\int_{\Gone{\fs}}\fdeg{l}\,\dsc{l}{\fp^{-1}g^{-1}\fp}\;W(y\,g)\,dg$ for every \m{y\in\Gone{\fs}};
\item \m{W(x)\ne0};
\item $\nrm{W(y)}\le\int_{\bgr{\fs}}\bigabs{\dsc{l}{\fp^{-1}x^{-1}\sigma_1\,y\,\fp}}\,d\bms{\fs}(\sigma)$ for every
  \m{y\in\Gone{\fs}};
\item for every \m{y\in\Gone{\fs}},
  \[\int_{\bgr{\fs}}\Bigl(\frac{2}{\op{\fp^{-1}x^{-1}\sigma_1\,y\,\fp}}\Bigr)^{8\,\rsh{\fs}^{2}+1}d\bms{\fs}(\sigma)
  \ \le\ P\cdot\nrm{W(x)}.\]
\end{enumerate}
The integrals in \textup{(P3)} and \textup{(P4)} are integrals of non-negative functions, with values in
\m{[0,\infty]}, and $\nrm{\,\cdot\,}$ is the Euclidean norm of~$\CC^{I}$.
\end{definition}
\leanline{\lean{HasPeakAt}}

Condition (P1) says that $W$ is fixed by the right convolution with the function
$g\mapsto\fdeg{l}\,\dsc{l}{\fp^{-1}g\,\fp}$, the kernel of the projection to the discrete series of~$l$
(Input~\inp{12}). In the formal text the predicate has the five
arguments \m{l,\mathrm{rep},W,x,P}, and $\fs$, $\varepsilon$, $\fp$, $I$ are implicit.

\begin{definition}[type of a Bessel system]\label{d:bestype}
Let $\fs$ be a shape, $\varepsilon\colon\Ind{\fs}\to\{\pm1\}$ signs, \m{\fp\in\frm{\fs}{\varepsilon}} a frame, $b$ a
real number, $f$ a natural number, and $\mathcal{B}$ a Bessel system for \m{(\fs,b,f)} in the sense of the unit on
the second moment: among its components are a finite index set $I(\mathcal{B})$, a continuous function
$W_{\mathcal{B}}\colon\Gone{\fs}\to\CC^{I(\mathcal{B})}$ called its full Bessel function, and a positive number
\m{\deg\mathcal{B}}. Let \m{l\in\ZZ^{\Ind{\fs}}} and let $w$ be a dominant weight of the unitary group of $\CC^{\kps{\fs}}$, in the encoding of the library. We say that
\m{\isbes{\mathcal{B}}{l}{w}} \kw{holds} if
\begin{enumerate}[label=\textup{(B\arabic*)},leftmargin=2.9em,itemsep=1pt]
\item $\deg\mathcal{B}=\prod_{i,i':\,l_{i'}<l_i}\tfrac12(l_i-l_{i'})\big/\prod_{n<\rsh{\fs}}n!$, the value of Weyl's
  dimension polynomial at~$l$;
\item there are a matrix $Y$ with rows indexed by the standard tableaux of the weight~$w$ for $\kps{\fs}$ and
  columns indexed by $I(\mathcal{B})$, and a matrix $Y'$ with \m{Y\,Y'=1} and \m{Y'\,Y=1}, such that
  $Y\,W_{\mathcal{B}}(\emb{\fs}{h}\,y)=\mathbb{S}_w(h)\,Y\,W_{\mathcal{B}}(y)$ for all \m{h\in\Us{1}} and
  \m{y\in\Gone{\fs}}, where $\mathbb{S}_w(h)$ is the matrix of $h$ in the Schur module of weight~$w$ in the basis of
  the tableaux, as the library defines it;
\item $\int_{\Us{\varepsilon}}\fdeg{l}\,\overline{\dsc{l}{g}}\;W_{\mathcal{B}}(y\,\fp\,g\,\fp^{-1})_i\,dg=W_{\mathcal{B}}(y)_i$
  for every \m{y\in\Gone{\fs}} and \m{i\in I(\mathcal{B})}.
\end{enumerate}
\end{definition}
\leanline{\lean{IsBessel}}

This predicate records that a Bessel system has the degree, the $\Us{1}$-type and the discrete series type that
belong to~$l$ and~$w$. The unit proves nothing about it; the units on the second moment and on the real place use
it (§\ref{s:use}). In (B2) the formal text writes the Schur matrix as the transpose of the Schur functor applied
to the transpose of~$h$; we have not unfolded the library's Schur functor further than this.

\subsection*{\normalsize Adelic unitary groups}\label{s:adelic}
\refstepcounter{subsection}
Theorem~6 uses the following objects of the background library, in its coordinates. Let $E$ be a CM field with a
CM type: for every infinite place~$v$ an embedding \m{E\to\CC} that induces~$v$ is chosen, and \m{\abs{a}_v} is the
absolute value of the image of \m{a\in E} under it. The number of infinite places is~$f_E$. Let $\Omega$ be a finite
set and $\varepsilon$ signs on~$\Omega$.

\begin{definition}[the adelic group and a space with a point]\label{d:adelic}
\begin{enumerate}[label=\textup{(\alph*)},leftmargin=2.3em,itemsep=1pt]
\item \m{G_\infty(E,\varepsilon)} is the product of copies of $\Us{\varepsilon}$, one for each infinite place~$v$,
  with the product of the Haar measures; \m{G_{\mathrm f}(E,\Omega)} is the unitary group of the identity matrix
  over the finite adeles of~$E$; and $G(E,\varepsilon)=G_\infty(E,\varepsilon)\times G_{\mathrm f}(E,\Omega)$, with
  components \m{g=(g_\infty,g_{\mathrm f})} and the product Haar measure of the library. For an integer \m{a_0\ne0}
  of~$E$, $\mathsf{K}_\Omega(a_0)$ is the open compact subgroup of the \m{k\in G_{\mathrm f}(E,\Omega)} with integral entries
  such that all entries of \m{k-1} are multiples of~$a_0$ by integral finite adeles.
\item A \kw{hermitian space with a point over~$E$ for~$\varepsilon$} is a pair \m{V=(J_V,x_V)} of an
  \m{\Omega\times\Omega} matrix $J_V$ over~$E$ and an invertible adelic matrix~$x_V$ with
  \m{x_V^{*}J_Vx_V} equal to $\dg{\varepsilon}$ at the infinite places and to~$1$ at the finite places. Then
  $\Gamma(V)=x_V^{-1}\,\mathrm{U}(J_V)(E)\,x_V$ is a discrete subgroup of $G(E,\varepsilon)$, and $\varpi_V$ is the
  measure of the library on $G(E,\varepsilon)$ that computes integrals over $\Gamma(V)\backslash G(E,\varepsilon)$.
  For \m{x\in G(E,\varepsilon)}, \m{V\cdot x=(J_V,x_Vx)}.
\item The \kw{lattice} of~$V$ consists of the vectors \m{\eta\in E^{\Omega}} for which \m{x_V^{-1}\eta} is integral
  at all finite places; the \kw{length} $\nrm{\eta}_V$ is the Euclidean norm of the vector of the components of
  \m{x_V^{-1}\eta} at the infinite places, a vector of $\CC^{f_E\abs{\Omega}}$; and $\lambda_1(V)$ is the infimum
  of~$1$ and of the lengths of the non-zero lattice vectors.
\item \m{\langle1\rangle} is the positive line: the space with matrix $(1)$ and point~$1$ on the index set
  \m{\{*\}}; \m{V\oplus\langle1\rangle} has block-diagonal matrix and point, for the signs \m{\varepsilon\oplus1} on
  \m{\Omega\sqcup\{*\}}; and \m{g\mapsto g\oplus1} is the block-diagonal embedding of $G(E,\varepsilon)$ into
  \m{G(E,\varepsilon\oplus1)}, of $\Us{\varepsilon}$ into \m{\Us{\varepsilon\oplus1}}, and of \m{G_\infty(E,\varepsilon)}
  into \m{G_\infty(E,\varepsilon\oplus1)}.
\item A function \m{F\colon G(E,\varepsilon)\to\CC} is \kw{automorphic for~$V$} if it is continuous, invariant under
  left translations by $\Gamma(V)$, and invariant under right translations by some open subgroup of
  \m{G_{\mathrm f}(E,\Omega)}. It \kw{has level~$\mathsf{K}$}, for a subgroup $\mathsf{K}$ of \m{G_{\mathrm f}(E,\Omega)}, if
  \m{F(g\cdot(1,k))=F(g)} for all \m{g} and all \m{k\in \mathsf{K}}.
\end{enumerate}
\end{definition}
The formal text requires of~$E$, for the positive line to exist as a space with a point, that it contains a
primitive root of unity of order~$5040$.

\sect{What is taken for granted}{s:inputs}
The proofs below use the following results of other units. All but \inp{14} are proved in the background library
of the artifact (the unit numbers are those of the reading project's ledger); they are stated as the dossier of the
unit reads them, and the statements \inp{7}, \inp{9}, \inp{10}, \inp{11} and \inp{14} were compared here with their
Lean text. A reference is given where the second reader of the project or the catalogue of the write-up has one;
``reference to be confirmed'' means that none is recorded. Facts of Mathlib (Sylvester's identity
\m{\det(1+XY)=\det(1+YX)}, the spectral theorem for hermitian matrices, Fubini's and Tonelli's theorems, Markov's
inequality) are used without number.

\subhead{Inputs.} In \inp{1} to \inp{4}, $X$, $Y$, $D$ are
\m{\Omega\times\Omega} matrices; in \inp{7} to \inp{10}, $\varepsilon\colon\Omega\to\{\pm1\}$ are signs,
\m{l\in\ZZ^{\Omega}} and \m{g\in\Us{\varepsilon}}.
\begin{enumerate}[label=\textup{(I\arabic*)},leftmargin=2.7em,itemsep=2.5pt]
\item \emph{Hilbert--Schmidt norm} (unit 112, linear algebra; elementary). For matrices $X$, $Y$:
  $\abs{\tr(XY)}\le\hs{X}\hs{Y}$; $\tr(X^{*}X)=\hs{X}^{2}$; \m{\hs{uX}=\hs{X}} if \m{u^{*}u=1}; \m{\hs{pX}\le\hs{X}} for an
  orthogonal projection~$p$;
  replacing entries of~$X$ by zero does not increase $\hs{X}$.
\item \emph{Inverses near~$1$} (unit 112; elementary). If \m{\hs{X}\le c\le\tfrac12} then \m{1+X} is invertible and
  \m{\hs{(1+X)^{-1}-1}\le2c}.
\item \emph{The logarithm} (unit 112; an integrated form of Jacobi's formula, elementary). If \m{\hs{X}\le\tfrac12},
  then $s\mapsto\tr\bigl(D(1+sX)^{-1}X\bigr)$ is integrable on $[0,1]$ with integral \m{\tr(D\,\lgo{X})};
  \m{\det(1+X)=\exp\tr\lgo{X}}; and \m{\hs{\lgo{X}}\le2\hs{X}}.
\item \emph{Principal minors} (unit 112; elementary). (a) If \m{X_{ii'}=0} or \m{Y_{i'i''}=0} whenever \m{i,i''\in\Pi}
  and \m{i'\notin\Pi}, then $\pmn{\Pi}{XY}=\pmn{\Pi}{X}\,\pmn{\Pi}{Y}$. (b) If \m{o\colon\Omega\to\ZZ} is injective
  and \m{X_{ii'}=0} whenever \m{o_{i'}<o_i}, then \m{\det X=\prod_iX_{ii}}. (c) For invertible~$X$, the library's
  product of principal minors of weight~$w$ is the number $\wm{w}{X}$ of Definition~\ref{d:wm}.
\item \emph{Lagrange weights} (unit 112; Lagrange interpolation, elementary). For integers
  $\tau_0>\tau_1>\dots>\tau_{n_0-1}$ there are real numbers \m{c_n\ge1} with
  $c_n^{2}=(-1)^{n}\prod_{n'\ne n}(\tau_n-\tau_{n'})$ such that for every real polynomial $\pi$ of degree less
  than~$n_0$ the coefficient of degree \m{n_0-1} of~$\pi$ is $\sum_n(-1)^{n}\pi(\tau_n)/c_n^{2}$.
\item \emph{Witt's theorem} (unit 120). Let $J$ be an invertible hermitian \m{\Omega\times\Omega} matrix and $Y_1$, $Y_2$
  matrices with $\Omega$ rows, the same set of columns, and left inverses, such that \m{Y_1^{*}JY_1=Y_2^{*}JY_2}. Then
  \m{g\,Y_1=Y_2} for some~$g$ with \m{g^{*}Jg=J}. Reference: U.~First, \emph{Witt's extension theorem for quadratic
  spaces over semiperfect rings}, Remark~2.1 and Corollary~4.9; relation: the theorem itself, for a field with a
  non-trivial involution (second reader's verdict).
\item \emph{Flensted-Jensen's formula and its invariance} (units 116 and 112). Suppose that $l$ has gap~$H$
  for~$\varepsilon$ for some real~$H$. Call a matrix~$Y$ \kw{accretive} if \m{Y+Y^{*}} is positive definite; this
  holds if \m{\hs{Y-1}<1}. (a) If $a$ is an \m{\Omega\times\Omega} matrix commuting with $\dg{\varepsilon}$ such that
  \m{a^{-1}g} is accretive, then \m{\dsc{l}{g}=\fj{l}{a}{a^{-1}g}}. (b) If $u'$ is a matrix with \m{u'_{ii'}=0} unless
  \m{l_i<l_{i'}} and \m{\varepsilon_i=\varepsilon_{i'}}, and $Y$ and \m{(1-u')Y} are accretive, then
  $\fj{l}{1}{((1-u')Y)^{*}}=\overline{\fj{l}{1}{Y}}$. Reference for (a): M.~Flensted-Jensen, \emph{Discrete series for
  semisimple symmetric spaces}, Ann.\ of Math.\ 111 (1980); place and relation to be confirmed. For (b): reference
  to be confirmed.
\item \emph{Elementary properties of the coefficient} (units 116, 100, 18). (a) \m{\ktn{\varepsilon}{w}>0} for every integer
  vector~$w$, and $\dsc{l}{\,\cdot\,}$ is continuous on $\Us{\varepsilon}$. Suppose that $l$ has gap~$H$ for~$\varepsilon$
  for some~$H$. Then: (b) \m{\dsc{l}{h}=\wm{\lowk{l}}{h}} for \m{h\in\Ks{\varepsilon}}; (c)
  $\dsc{-l}{g}=\overline{\dsc{l}{g}}$ and \m{\lowk{-l}=-\lowk{l}}; (d) \m{\ups{l}_i<\ups{l}_{i'}} if and only if
  \m{l_i<l_{i'}}.
\item \emph{Decay} (unit 116). If $l$ has gap~$H$ for~$\varepsilon$ and $c$ is a natural number with \m{c\le H}, then
  $\abs{\dsc{l}{g}}^{2}\,\abs{\det\bd{\varepsilon}{g}}^{c}\le1$; if moreover \m{\Omega\ne\emptyset}, then
  $\abs{\dsc{l}{g}}\le(2/\op{g})^{c}$. In particular \m{\abs{\dsc{l}{g}}\le1}. The catalogue relates the first bound
  to Flensted-Jensen (1980) and to B.~Sun, \emph{Matrix coefficients of cohomologically induced representations},
  Compositio Math.\ 143 (2007), and does not identify it with a published statement: reference to be confirmed.
\item \emph{The coefficient near~$1$} (unit 116; elementary from a Lipschitz bound of the library). Let
  \m{\Omega\ne\emptyset}, let $l$ have large gap~$A_0$ for~$\varepsilon$ for some natural~$A_0$, and \m{\abs{l_i}\le L}
  for all~$i$. (a) If \m{\beta_0\ge\abs{\Omega}+L+1} and $\op{g-1}\le(200\beta_0^{5})^{-1}$, then
  \m{\abs{\dsc{l}{g}}\ge\tfrac12}. (b) If $\nu$ is a right-invariant measure on a group~$\mathfrak{g}$ with measurable
  translations, \m{j\colon \mathfrak{g}\to\Us{\varepsilon}} a homomorphism with $\abs{\dsc{l}{j(\,\cdot\,)}}^{2}$ integrable, and
  \m{R\ge1}, then
  $\nu\{\op{j(\sigma)}\le R\}\le(2000(\abs{\Omega}+L+1)^{6}R^{2})^{2\abs{\Omega}^{2}}\int4\abs{\dsc{l}{j(\sigma)}}^{2}d\nu$.
\item \emph{Balls and decay on a group over $\Us{\varepsilon}$} (unit 115, Haar measure; the catalogue cites Hua, Macdonald
  and Gross for the volume growth and has no place: reference to be confirmed). Let \m{\Omega\ne\emptyset}, let $\mathfrak{g}$ be a
  group with a topology, its Borel $\sigma$-algebra and measurable translations, $\nu$ a right-invariant measure
  on~$\mathfrak{g}$, \m{j\colon \mathfrak{g}\to\Us{\varepsilon}} a homomorphism and \m{y\in\Us{\varepsilon}}. (a) For real \m{D,R\ge1} and
  \m{K\ge3DR^{2}}:
  $\nu\{\op{j(\sigma)y}\le R\}\le K^{2\abs{\Omega}^{2}}\,\nu\{\op{j(\sigma)-1}\le D^{-1}\}$. (b) If $j$ is proper and
  $\nu$ is finite on compact sets, then $\nu\{\op{j(\sigma)-1}\le R\}<\infty$ for every~$R$, and for every natural
  \m{c>8\abs{\Omega}^{2}} the function $\sigma\mapsto(2/\op{j(\sigma)y})^{c}$ is $\nu$-integrable with
  $\int(2/\op{j(\sigma)y})^{c}\,d\nu\le16^{c}\,\nu\{\op{j(\sigma)-1}\le1\}$. (c) If $\nu$ is positive on non-empty open
  sets and \m{f\colon \mathfrak{g}\to\RR} is continuous with \m{f(1)<R}, then \m{\nu\{f\le R\}\ne0}. (d) \m{\op{g}\ge1} and
  \m{\op{g^{-1}}=\op{g}} for \m{g\in\Us{\varepsilon}}, and \m{g\mapsto g\oplus1} is a proper map from $\Us{\varepsilon}$ to
  \m{\Us{\varepsilon\oplus1}} (unit 43).
\item \emph{Schur orthogonality} (units 112 and 117). Let $\fs$ be a shape, $\fp$ a frame of~$\fs$ for signs~$\varepsilon$,
  and let \m{l\in\ZZ^{\Ind{\fs}}} have large gap~$A$ for~$\varepsilon$, for some natural~$A$. Then
  $k(g)=\fdeg{l}\,\dsc{l}{\fp^{-1}g\,\fp}$ is continuous and integrable on $\Gone{\fs}$, satisfies
\m{k(g^{-1})=\overline{k(g)}} for every $g$ and \m{k(1)\neq0}, \m{\fdeg{l}>0}, and
  $\int k(g^{-1})\,k(y_1\,g\,y_2)\,dg=k(y_1)\,k(y_2)/k(1)$ for all \m{y_1,y_2\in\Gone{\fs}}. Reference:
  Harish-Chandra, \emph{Discrete series for semisimple Lie groups~II}, Acta Math.\ 116 (1966), pp.~87--88, for the
  orthogonality of square-integrable coefficients with the formal degree; relation unresolved (second reader's
  verdict): the cited pages do not give the integrability for the explicit gap nor the normalisation of the
  measure, which the library proves.
\item \emph{The Bessel subgroup} (units 1, 118, 119). Let $\fs$ be a shape. (a) $\sd{q\,q'}=\sd{q}+\sd{q'}$ and
  \m{\sd{q^{-1}}=-\sd{q}} for \m{q,q'\in\Qq{\fs}}, and \m{\sd{q}=0} for \m{q\in\Hh{\fs}}. (b) Let $C_{\fs}$ be the \kw{chain
  matrix}: the \m{\Ind{\fs}\times\Ind{\fs}} matrix with \m{C_{\fs}e_a=f_{a+1}} for \m{a<\rhs{\fs}}, \m{C_{\fs}x_i=0},
  \m{C_{\fs}z=-e^{*}_{\rhs{\fs}-1}} (if \m{\rhs{\fs}\ge1}), \m{C_{\fs}e^{*}_{a+1}=e^{*}_{a}} and \m{C_{\fs}e^{*}_0=0}. Then
  $\tr(C_{\fs}\,g)=\sd{g}+\overline{\sd{g^{-1}}}$ for \m{g\in\Gone{\fs}}. (c) Conjugation by an element of $\Qq{\fs}$
  and right translation by an element of $\Uu{\fs}$ preserve~$\hU{\fs}$. (d) The formula of
  Definition~\ref{d:shape}(e) for~$\hq{\fs}$; it is an instance of Weil's formula for the integral over a quotient, which
  the library proves for a measure with a cutoff function. Reference: Kallenberg, Probab.\ Theory Related Fields 139
  (2007), as cited by Gentner and Last; relation unresolved (second reader's verdict: the library's statement is
  more general than the cited one and follows from Tonelli's theorem and the two invariances).
\item \emph{A bound for the size} (unit 132 of the argument, the complex chart at the real place). Let $\fs$ be a
  shape, \m{x\in\Gone{\fs}}, \m{M_0\ge1} real, and \m{V_0,V_1,V_2,\dots} vectors of $\CC^{\Ind{\fs}}$ with
  \m{V_a=x^{-1}f_a} for \m{a\le\rhs{\fs}}, with Euclidean norms \m{\nrm{V_a}\le M_0} for \m{a\le2\rhs{\fs}}, and with
  \m{V_a^{*}\Jone{\fs}V_{a'}=-1} if \m{a+a'=2\rhs{\fs}} and \m{=0} if \m{a+a'<2\rhs{\fs}}. Then
  $\bsz{x}\le8\,(\rhs{\fs}+1)\,M_0$.
\end{enumerate}

\noindent The adelic inputs of Theorem~6 are listed in §\ref{s:period}, where they are used.

\sect{Determinants of compressions, to second order}{s:minors}
In this section and the next, norms without mention are Hilbert--Schmidt norms. The first result compares
\m{\log\det(1+pX)}, for an orthogonal projection~$p$, with \m{\tr(q\,\lgo{X})} for a second projection~$q$ close to~$p$.
To first order the two agree when \m{q=p}; the second-order term is a quarter of the square of the commutator, and
the error is cubic: one factor $\hs{X}$ and two factors of the size of commutators.

\begin{proposition}[compressed determinants]\label{p:compress}
Let $\Omega$ be a finite set, let $p$, $q$, $X$, $Y$ be \m{\Omega\times\Omega} matrices and $\delta$, $c$ real numbers.
Suppose that $p$ and $q$ are orthogonal projections \textup{(}\m{p=p^{2}=p^{*}} and \m{q=q^{2}=q^{*}}\textup{)}, that
\m{\delta\le\tfrac12}, and that
\[\begin{gathered}\hs{X}\le\delta,\qquad \tr\bigl(\comm{p}{q}\,Y\bigr)=0,\qquad \hs{\comm{p}{X}-Y}\le\delta\,c,\\
\hs{Y}\le c,\qquad \hs{q-p}\le c.\end{gathered}\]
Then there is a complex number~$\eta$ with
\[\det(1+p\,X)=\exp\Bigl(\tr\bigl(q\,\lgo{X}\bigr)-\tfrac14\tr(Y^{2})+\eta\Bigr)\qquad\text{and}\qquad \abs{\eta}\le100\,\delta\,c^{2}.\]
\end{proposition}
\begin{proof}
The numbers $\delta$ and $c$ are non-negative, being bounds of norms. Put \m{A=\comm{p}{X}} and \m{D=q-p}; then
$\hs{A}\le\hs{Y}+\hs{A-Y}\le c+\delta c\le\tfrac32c$. Since \m{\hs{pX}\le\delta\le\tfrac12} by~\inp{1}, Input~\inp{3} gives
\m{\det(1+pX)=\exp\tr\lgo{pX}}, and it writes \m{\tr\lgo{pX}} and \m{\tr(q\,\lgo{X})} as integrals over \m{s\in[0,1]}.
With \m{\tfrac14=\int_0^1\tfrac s2\,ds} the claimed identity holds for
\[\eta=\int_0^1\mathcal{I}(s)\,ds,\qquad \mathcal{I}(s)=\tr(V\,pX)-\tr(q\,W\,X)+\tfrac s2\tr(Y^{2}),\]
where \m{W=(1+sX)^{-1}} and \m{V=(1+s\,pX)^{-1}},
and it suffices to show $\abs{\mathcal{I}(s)}\le100\,\delta c^{2}$ for each \m{s\in[0,1]}. Fix~$s$. By~\inp{2},
\m{\hs{W-1}\le2\delta} and \m{\hs{V-1}\le2\delta}, so \m{\op{W}\le2} and \m{\op{V}\le2}. We split
$\mathcal{I}(s)=\mathcal{I}_1+\mathcal{I}_2+\mathcal{I}_3$ with
\[\begin{gathered}\mathcal{I}_1=\tr(V\,pX)-\tr(p\,W\,X)+\tfrac s2\tr(A^{2}),\qquad\mathcal{I}_2=-\tr(D\,W\,X),\\
\mathcal{I}_3=-\tfrac s2\tr\bigl((A-Y)(A+Y)\bigr),\end{gathered}\]
using $\tr(A^{2}-Y^{2})=\tr((A-Y)(A+Y))$. At once $\abs{\mathcal{I}_3}\le\tfrac12\,\delta c\,(\tfrac32c+c)=\tfrac54\delta c^{2}$.

\emph{The term $\mathcal{I}_1$: one projection.} We use
\[\begin{gathered}pAp=0=(1-p)A(1-p),\qquad pA=pX(1-p),\\
Ap=-(1-p)Xp,\qquad\tr(pA^{2})=\tfrac12\tr(A^{2}),\end{gathered}\]
the last because \m{\tr(pA^{2})=\tr((1-p)A^{2})}. From \m{p(1+sX)-(1+s\,pX)p=s\,pA} we get, multiplying by $V$ on the
left and by $W$ on the right, \m{Vp-pW=s\,V\,pA\,W}, so
\[\tr(V\,pX)-\tr(p\,WX)=s\,\tr\bigl(pA\cdot WXV\bigr)=s\,\tr\bigl(pA\,(1-p)\,WXV\,p\bigr).\]
Now \m{\comm{p}{W}=-s\,WAW}, so \m{(1-p)W=W(1-p)+s\,WAW}; and \m{Vp=pVp}, because \m{(1-p)V=1-p} by
\m{V=1-s\,pXV}. Therefore
\[(1-p)\,WXV\,p=W(1-p)Xp\,Vp+s\,WAW\,XVp=-\,WA\,Vp+s\,WAW\,XVp,\]
and, writing \m{WAV-A=(W-1)AV+A(V-1)},
\[\mathcal{I}_1=-s\,\tr\bigl(pA\,\bigl((W-1)AV+A(V-1)\bigr)\,p\bigr)+s^{2}\,\tr\bigl(pA\cdot WAW\cdot XV\,p\bigr).\]
By \inp{1} and the bounds for $W$ and~$V$,
$\abs{\mathcal{I}_1}\le\hs{A}^{2}(2\delta\cdot2+2\delta)+\hs{A}\cdot4\hs{A}\cdot\delta\cdot2=14\,\delta\hs{A}^{2}\le\tfrac{63}{2}\,\delta c^{2}$.

\emph{The term $\mathcal{I}_2$: the change of projection.} A computation with \m{p^{2}=p}, \m{q^{2}=q} shows that
$D^{2}$ commutes with~$p$ and that \m{pDp=-pD^{2}}, \m{(1-p)D(1-p)=(1-p)D^{2}}. So \m{D=D_0+D_1} with
\[\begin{gathered}D_0=(1-2p)\,D^{2},\qquad D_1=pq(1-p)+(1-p)qp,\\
\text{while}\quad \comm{p}{q}=pq(1-p)-(1-p)qp.\end{gathered}\]
For every matrix~$Z$, since $(1-p)Zp=-(1-p)\comm{p}{Z}p$ and $pZ(1-p)=p\,\comm{p}{Z}(1-p)$, the cyclicity of the
trace gives $\tr(D_1Z)=-\tr\bigl(\comm{p}{q}\,\comm{p}{Z}\bigr)$. For \m{Z=WX} we have
\m{\comm{p}{WX}=WA-s\,WAW\,X}, and \m{\tr(\comm{p}{q}\,Y)=0} by hypothesis, so
\[\mathcal{I}_2=-\tr(D_0\,WX)+\tr\Bigl(\comm{p}{q}\,\bigl((A-Y)+(W-1)A-s\,WAW\,X\bigr)\Bigr).\]
Here \m{\hs{D_0}=\hs{D^{2}}\le c^{2}}, because \m{1-2p} is unitary; \m{\hs{WX}\le2\delta}; $\hs{\comm{p}{q}}=\hs{D_1}\le\hs{D}\le c$;
and the second factor has norm at most $\delta c+2\delta\hs{A}+4\hs{A}\delta\le10\,\delta c$. Hence
$\abs{\mathcal{I}_2}\le2\,\delta c^{2}+10\,\delta c^{2}$.

Together $\abs{\mathcal{I}(s)}\le(\tfrac{63}{2}+12+\tfrac54)\,\delta c^{2}\le100\,\delta c^{2}$.
\end{proof}
\leanline{\lean{existsdetcompresseqexplogOneAdd}}

The hypothesis \m{\tr(\comm{p}{q}Y)=0} removes the only term of second order that the change from $p$ to~$q$ could
produce: the pairing of the off-diagonal part of \m{q-p} with the off-diagonal part of~$X$.

The second result applies this to Flensted-Jensen's formula. For \m{g\in\Us{\varepsilon}} near~$1$ it finds a
matrix~$y$ that may replace $g$ in the formula and whose principal minors, conjugated by any element of the
maximal compact subgroup, are Gaussian in \m{g-1} with the \emph{same} sign for compact and non-compact pairs. For
$g$ itself the signs differ, because \m{X=g-1} is, to first order, anti-hermitian on the compact pairs and
hermitian on the non-compact ones; the passage from $g$ to~$y$ removes the triangular parts of the compact block
of~$X$, which costs nothing in the formula and flips the sign.

\begin{proposition}[principal minors for Flensted-Jensen's formula]\label{p:minors}
Let $\Omega$ be a finite set, $\varepsilon\colon\Omega\to\{\pm1\}$ signs, $H$ a real number, \m{l\in\ZZ^{\Omega}} a
vector that has gap~$H$ for~$\varepsilon$, and \m{g\in\Us{\varepsilon}}. Put \m{X=g-1} and suppose that
\m{\hs{X}<\tfrac1{1000}}. Then there is an \m{\Omega\times\Omega} matrix~$y$ with \m{\dsc{l}{g}=\fj{l}{1}{y}} such that
for every subset \m{\Pi\subseteq\Omega} with \m{l_i<l_{i'}} for all \m{i\in\Pi} and \m{i'\notin\Pi}, and for every
\m{h\in\Ks{\varepsilon}}, there is a complex number~$\eta$ with
\begin{gather*}
\pmn{\Pi}{h^{*}\,y\,h}=\exp\Bigl(\tr\bigl(\cpr{\Pi}\,\lgo{X}\bigr)+\tfrac14\hs{\comm{\cpr{\Pi}}{X}}^{2}+\eta\Bigr),\\
\abs{\eta}\le200000\,\hs{X}\,\Bigl(\hs{\comm{\cpr{\Pi}}{X}}^{2}+\hs{\comm{\cpr{\Pi}}{h}}^{2}\Bigr).
\end{gather*}
\end{proposition}
\begin{proof}
For an \m{\Omega\times\Omega} matrix~$Z$ let $Z^{\uparrow}$ be the matrix that keeps the entries $Z_{ii'}$ with
\m{\varepsilon_i=\varepsilon_{i'}} and \m{l_i<l_{i'}} and replaces the others by zero, and $Z^{\downarrow}$ the same with
\m{l_i>l_{i'}}. Put \m{\delta_0=\hs{X}} and
\[y=(1-X^{\uparrow})\,g\,(1-X^{\downarrow}).\]
With $X_1=(1-X^{\uparrow})g-1=X-X^{\uparrow}-X^{\uparrow}X$ we have $y-1=X_1-X^{\downarrow}-X_1X^{\downarrow}$, and by~\inp{1}
$\hs{X^{\uparrow}},\hs{X^{\downarrow}}\le\delta_0$, so
\begin{equation}\label{e:ynear}
\hs{X_1}\le2\delta_0+\delta_0^{2},\qquad \hs{y-1}\le3\delta_0+3\delta_0^{2}+\delta_0^{3}\le4\,\delta_0<\tfrac12 .
\end{equation}

\emph{The formula.} The matrices $g$, \m{(1-X^{\uparrow})g} and $y$ are within Hilbert--Schmidt distance less
than~$1$ of~$1$, hence accretive, as are their adjoints. By \inp{7}(a) with \m{a=1},
\m{\dsc{l}{g}=\fj{l}{1}{g}}. By \inp{7}(b) for \m{u'=X^{\uparrow}} and \m{Y=g},
$\fj{l}{1}{((1-X^{\uparrow})g)^{*}}=\overline{\fj{l}{1}{g}}$. The matrix \m{u'=(X^{\downarrow})^{*}} also has its non-zero
entries at compact pairs with \m{l_i<l_{i'}}, and $(1-u')\,((1-X^{\uparrow})g)^{*}=y^{*}$; so \inp{7}(b) for this $u'$ and
\m{Y=((1-X^{\uparrow})g)^{*}} gives $\fj{l}{1}{y}=\overline{\fj{l}{1}{((1-X^{\uparrow})g)^{*}}}=\fj{l}{1}{g}$.

\emph{Upper sets.} Let \m{\Pi'\subseteq\Omega} be such that \m{l_{i'}<l_i} for all \m{i\in\Pi'}, \m{i'\notin\Pi'}. Then
\begin{equation}\label{e:upper}
\pmn{\Pi'}{y}=\pmn{\Pi'}{g}.
\end{equation}
Indeed \m{(X^{\uparrow})_{ii'}=0} for \m{i\in\Pi'}, \m{i'\notin\Pi'}, and \m{(X^{\downarrow})_{i'i}=0} for such indices, so by
\inp{4}(a) the minor is multiplicative on the three factors of~$y$; and $\pmn{\Pi'}{1-X^{\uparrow}}=1=\pmn{\Pi'}{1-X^{\downarrow}}$
by \inp{4}(b), these being determinants of matrices that are triangular for the injective function~$l$, or~$-l$,
with diagonal entries~$1$. In particular \m{\det y=\det g}.

Now fix $\Pi$ and $h$ as in the statement, and write \m{p=\cpr{\Pi}}, \m{A=\comm{p}{X}}, \m{\alpha_0=\hs{A}},
\m{\kappa_0=\hs{\comm{p}{h}}}. The map \m{Z\mapsto\comm{p}{Z}} multiplies the entry at \m{(i,i')} by $1$, $-1$ or~$0$; so
it commutes with the passage to $Z^{\uparrow}$, $Z^{\downarrow}$ and \m{\bd{\varepsilon}{Z}}, and it kills diagonal
entries. Let $A_c=\bd{\varepsilon}{A}=A^{\uparrow}+A^{\downarrow}$ be the compact part of~$A$ and \m{A_n=A-A_c} its
non-compact part; they have norms at most~$\alpha_0$, and \m{\tr(A_cA_n)=0}.

\emph{The commutator of $y$.} From the formulas for \m{y-1} and $X_1$,
\begin{equation}\label{e:commy}
\comm{p}{y-1}=A_n-R,\qquad R=G+X_1A^{\downarrow}+(A-A^{\uparrow}-G)\,X^{\downarrow},\qquad G=A^{\uparrow}X+X^{\uparrow}A,
\end{equation}
and $\hs{R}\le2\delta_0\alpha_0+(2\delta_0+\delta_0^{2})\alpha_0+(2\alpha_0+2\delta_0\alpha_0)\delta_0\le7\,\delta_0\alpha_0$.

\emph{Three expansions.} Let $p'=1-p=\cpr{\Omega\setminus\Pi}$ and \m{q=h\,p\,h^{*}}, an orthogonal projection with
$\hs{p-q}=\hs{\comm{p}{h}h^{*}}=\kappa_0$. We apply Proposition~\ref{p:compress} three times; the six slots are
\m{(p,q,X,Y,\delta,c)}.
\begin{enumerate}[label=\textup{(\roman*)},leftmargin=2.3em,itemsep=1pt]
\item With \m{(p',p',X,-A,\delta_0,\alpha_0)}: here \m{\comm{p'}{X}=-A}. So
  \[\det(1+p'X)=\exp\bigl(\tr(p'\lgo{X})-\tfrac14\tr(A^{2})+\eta_1\bigr)\qquad\text{with}\quad\abs{\eta_1}\le100\,\delta_0\alpha_0^{2}.\]
\item With $(p',p',y-1,-A_n,4\delta_0,4\alpha_0)$: by \eqref{e:commy}, \m{\comm{p'}{y-1}+A_n=R} has norm at most
  \m{16\,\delta_0\alpha_0}. So $\det(1+p'(y-1))=\exp\bigl(\tr(p'\lgo{y-1})-\tfrac14\tr(A_n^{2})+\eta_2\bigr)$ with
  $\abs{\eta_2}\le6400\,\delta_0\alpha_0^{2}$.
\item With $(q,p,y-1,A_n,4\delta_0,4(\alpha_0+\kappa_0))$: the matrices $h$ and~$p$ commute with $\dg{\varepsilon}$, so
  \m{\comm{q}{p}} is block diagonal, and its trace against~$A_n$, which is supported on the non-compact pairs,
  vanishes; and $\comm{q}{y-1}-A_n=\comm{q-p}{y-1}-R$ has norm at most
  $2\kappa_0\cdot4\delta_0+7\delta_0\alpha_0\le16\,\delta_0(\alpha_0+\kappa_0)$. So
  $\det(1+q(y-1))=\exp\bigl(\tr(p\,\lgo{y-1})-\tfrac14\tr(A_n^{2})+\eta_3\bigr)$ with
  $\abs{\eta_3}\le6400\,\delta_0(\alpha_0+\kappa_0)^{2}$.
\end{enumerate}
The left sides are minors: $\det(1+p'X)=\pmn{\Omega\setminus\Pi}{g}$, $\det(1+p'(y-1))=\pmn{\Omega\setminus\Pi}{y}$, and by
Sylvester's identity $\det(1+q(y-1))=\det(1+p\,h^{*}(y-1)h)=\pmn{\Pi}{h^{*}yh}$. The complement of~$\Pi$ is an upper
set, so the first two are equal by \eqref{e:upper}; and $\exp\tr\lgo{y-1}=\det y=\det g=\exp\tr\lgo{X}$ by~\inp{3}. Multiplying
(iii) by (ii), using \m{p+p'=1}, and dividing by~(i), which is not zero, we find
\begin{equation}\label{e:three}
\pmn{\Pi}{h^{*}yh}=\exp\Bigl(\tr(p\,\lgo{X})+\tfrac14\tr(A^{2})-\tfrac12\tr(A_n^{2})+\eta_2+\eta_3-\eta_1\Bigr).
\end{equation}

\emph{The sign.} Put \m{V=g^{-1}-1} and \m{E=VA+AV+VAV}. The relation
$g^{*}=\dg{\varepsilon}g^{-1}\dg{\varepsilon}$ gives $X^{*}=\dg{\varepsilon}V\dg{\varepsilon}$, and
\[\comm{p}{V}=-g^{-1}A\,g^{-1}=-(A+E),\qquad A^{*}=-\comm{p}{X^{*}}=\dg{\varepsilon}(A+E)\dg{\varepsilon}.\]
Conjugation by $\dg{\varepsilon}$ fixes $A_c$ and changes the sign of~$A_n$, so
\[\hs{A}^{2}=\tr(A^{*}A)=\tr\bigl((A+E)(A_c-A_n)\bigr)=\tr(A_c^{2})-\tr(A_n^{2})+\tr\bigl(E(A_c-A_n)\bigr).\]
With $\tr(A^{2})=\tr(A_c^{2})+\tr(A_n^{2})$ this turns the exponent of \eqref{e:three} into
$\tr(p\,\lgo{X})+\tfrac14\hs{A}^{2}+\eta$ for
\[\eta=\eta_2+\eta_3-\eta_1+\tfrac14\tr\bigl(E\,(A_n-A_c)\bigr).\]
By \inp{2}, \m{\hs{V}\le2\delta_0}, so $\hs{E}\le(4\delta_0+4\delta_0^{2})\alpha_0\le5\,\delta_0\alpha_0$, and \m{\hs{A_n-A_c}=\alpha_0}. Therefore
\[\abs{\eta}\le\delta_0\bigl(6400\,\alpha_0^{2}+6400\,(\alpha_0+\kappa_0)^{2}+100\,\alpha_0^{2}+2\,\alpha_0^{2}\bigr)\le19302\,\delta_0\,(\alpha_0^{2}+\kappa_0^{2}),\]
which is less than the bound claimed.
\end{proof}
\leanline{\lean{existsprincipalMinoreqexp}}

\sect{The Gaussian shape: proof of Theorem 1}{s:gauss}
For an integer vector $w$ on~$\Omega$ write \m{\abs{w}_1=\sum_i\abs{w_i}}; then \m{-\abs{w}_1\le w_i\le\abs{w}_1} for all~$i$,
so that Definition~\ref{d:wm} applies with \m{c_1=-\abs{w}_1} and \m{c_2=\abs{w}_1}.

\begin{lemma}[sums over the level sets]\label{l:levelsums}
Let $\Omega$ be a finite set, $w$ an integer vector on~$\Omega$ and $X$ an \m{\Omega\times\Omega} matrix. Then, the sums
being over the integers~$\tau$ with \m{-\abs{w}_1<\tau\le\abs{w}_1},
\[\sum_{\tau}\hs{\comm{\cpr{\{w\ge\tau\}}}{X}}^{2}=\gwn{w}{X}\qquad\text{and}\qquad
\sum_{\tau}\tr\bigl(\cpr{\{w\ge\tau\}}\,X\bigr)=\sum_{i\in\Omega}\bigl(w_i+\abs{w}_1\bigr)\,X_{ii}.\]
\end{lemma}
\begin{proof}
The entry of \m{\comm{\cpr{\{w\ge\tau\}}}{X}} at \m{(i,i')} is $X_{ii'}$ times the difference of the indicators of
\m{w_i\ge\tau} and \m{w_{i'}\ge\tau}. This difference is $\pm1$ for the integers~$\tau$ with
$\min(w_i,w_{i'})<\tau\le\max(w_i,w_{i'})$, which are \m{\abs{w_i-w_{i'}}} in number and lie in the range of summation, and
it is $0$ for the others. This gives the first identity. For the second, the diagonal entry of \m{\cpr{\{w\ge\tau\}}X}
at~$i$ is $X_{ii}$ for the \m{w_i+\abs{w}_1} values \m{\tau\le w_i} of the range, and $0$ for the others.
\end{proof}
\leanline{\lean{sumnormliecoordProjsq}}

\begin{lemma}[the products of minors decay on the compact group]\label{l:minorsK}
Let $\Omega$ be a finite set, $\varepsilon\colon\Omega\to\{\pm1\}$ signs, $w$ an integer vector on~$\Omega$ and
\m{h\in\Ks{\varepsilon}}. Then $\abs{\wm{w}{h}}^{2}\le\exp\bigl(-\tfrac12\,\gwn{w}{h}\bigr)$.
\end{lemma}
\begin{proof}
Since \m{\abs{\det h}=1}, Definition~\ref{d:wm} and Lemma~\ref{l:levelsums} reduce the claim to one factor: for a
subset \m{\Pi\subseteq\Omega} and \m{p=\cpr{\Pi}} we show $\abs{\pmn{\Pi}{h}}^{2}\le\exp(-\tfrac12\hs{\comm{p}{h}}^{2})$. By
Sylvester's identity $\pmn{\Pi}{h}=\det(1-p(1-h))=\det Z$ with \m{Z=(1-p)+p\,h\,p}. A positive semidefinite
matrix has determinant at most the exponential of its trace minus its size, by the inequality \m{a\le e^{a-1}} for its
eigenvalues~$a$; for $Z^{*}Z$ this reads
\[\abs{\det Z}^{2}\le\exp\bigl(\tr(Z^{*}Z)-\abs{\Omega}\bigr)=\exp\bigl(\hs{php}^{2}-\tr p\bigr).\]
As $h$ is unitary, \m{\hs{ph}^{2}=\tr p=\hs{hp}^{2}}, hence
$\hs{\comm{p}{h}}^{2}=\hs{ph(1-p)}^{2}+\hs{(1-p)hp}^{2}=2\,(\tr p-\hs{php}^{2})$, and the exponent is
\m{-\tfrac12\hs{\comm{p}{h}}^{2}}.
\end{proof}
\leanline{\lean{normweightMinorssqleexp}}

\begin{proof}[Proof of Theorem 1]
Put \m{X=g-1}, so \m{\hs{X}\le\delta\le10^{-12}} because \m{T\ge1}, and abbreviate \m{\Lambda=\lowk{l}},
\m{\Upsilon=\ups{l}}, $c_\Lambda=\ktn{\varepsilon}{\lowk{l}}$ and
\[E=\tr\bigl(\dg{\Lambda}\,\lgo{X}\bigr),\qquad q_0=\tfrac14\,\gwn{-\Upsilon}{X}\ \ge0,\qquad \kappa_{\mathrm{G}}=e^{-q_0}\in(0,1].\]
The vector $l$ has a gap, so Proposition~\ref{p:minors} gives a matrix~$y$ with
\m{\dsc{l}{g}=\fj{l}{1}{y}}. By Definition~\ref{d:coeff}, with the factor $\wm{\Lambda}{h}\overline{\wm{\Lambda}{h}}$ in the numerator,
\begin{equation}\label{e:fjy}
\dsc{l}{g}=\frac1{c_\Lambda}\int_{\Ks{\varepsilon}}\frac{\abs{\wm{\Lambda}{h}}^{2}}{\wm{-\Upsilon}{h^{*}y\,h}}\,dh .
\end{equation}

\emph{The denominator.} Fix \m{h\in\Ks{\varepsilon}}. For an integer~$\tau$ the set \m{\{-\Upsilon\ge\tau\}} satisfies the
condition of Proposition~\ref{p:minors}: if $-\Upsilon_i\ge\tau>-\Upsilon_{i'}$ then \m{\Upsilon_i<\Upsilon_{i'}}, so
\m{l_i<l_{i'}} by \inp{8}(d). Write $p_\tau=\cpr{\{-\Upsilon\ge\tau\}}$. The proposition gives numbers $\eta_\tau$ with
\[\pmn{\{-\Upsilon\ge\tau\}}{h^{*}yh}=\exp\Bigl(\tr(p_\tau\lgo{X})+\tfrac14\hs{\comm{p_\tau}{X}}^{2}+\eta_\tau\Bigr),\]
\[\abs{\eta_\tau}\le2\cdot10^{5}\,\delta\,\bigl(\hs{\comm{p_\tau}{X}}^{2}+\hs{\comm{p_\tau}{h}}^{2}\bigr),\]
and, for the set~$\Omega$ itself, where the commutators vanish, \m{\det(h^{*}yh)=\exp\tr\lgo{X}}. We insert this into
Definition~\ref{d:wm} for the weight~$-\Upsilon$ and sum over~$\tau$ by Lemma~\ref{l:levelsums}: the traces add up to
$-\tr(\dg{\Upsilon}\lgo{X})+\abs{\Upsilon}_1\tr\lgo{X}$, the squares to~$4q_0$. Hence
\begin{equation}\label{e:denominator}
\wm{-\Upsilon}{h^{*}yh}=\exp\bigl(q_0-E-r(h)\bigr),\qquad r(h)=\tr\bigl(\dg{\Upsilon-\Lambda}\lgo{X}\bigr)-\sum_\tau\eta_\tau .
\end{equation}
To bound $r(h)$ we use \eqref{e:rhoc}, the bound $\abs{\lgo{X}_{ii}}\le\hs{\lgo{X}}\le2\delta$ of~\inp{3}, Lemma~\ref{l:levelsums}
for $X$ and for~$h$, and, again by \eqref{e:rhoc},
$\gwn{-\Upsilon}{h}\le\gwn{\Lambda}{h}+2\abs{\Omega}\hs{h}^{2}=\gwn{\Lambda}{h}+2\abs{\Omega}^{2}$:
\begin{equation}\label{e:rbound}
\abs{r(h)}\le2\abs{\Omega}^{2}\delta+2\cdot10^{5}\,\delta\,\bigl(4q_0+\gwn{\Lambda}{h}+2\abs{\Omega}^{2}\bigr)
\le10^{6}\,\delta\,\bigl(G_0(h)+\abs{\Omega}^{2}\bigr),
\end{equation}
where $G_0(h)=q_0+\tfrac12\gwn{\Lambda}{h}$.

\emph{The comparison.} By \eqref{e:fjy} and \eqref{e:denominator}, and since $\int\abs{\wm{\Lambda}{h}}^{2}dh=c_\Lambda$ and
$dh$ has mass one,
\begin{align*}
\frac{\dsc{l}{g}}{e^{E}}-\kappa_{\mathrm{G}}&=\frac1{c_\Lambda}\int_{\Ks{\varepsilon}}\abs{\wm{\Lambda}{h}}^{2}\,e^{-q_0}\bigl(e^{r(h)}-1\bigr)\,dh,\\
\frac{\kappa_{\mathrm{G}}}{10}+e^{-100T}&=\frac1{c_\Lambda}\int_{\Ks{\varepsilon}}\Bigl(\frac{\kappa_{\mathrm{G}}}{10}\abs{\wm{\Lambda}{h}}^{2}+e^{-100T}c_\Lambda\Bigr)dh.
\end{align*}
So it suffices to show, for every \m{h\in\Ks{\varepsilon}},
\begin{equation}\label{e:pointwise}
\abs{\wm{\Lambda}{h}}^{2}\,e^{-q_0}\,\bigabs{e^{r(h)}-1}\ \le\ \frac{\kappa_{\mathrm{G}}}{10}\,\abs{\wm{\Lambda}{h}}^{2}+e^{-100T}\,c_\Lambda .
\end{equation}
If \m{\abs{r(h)}<\tfrac1{20}}, then $\abs{e^{r(h)}-1}\le2\abs{r(h)}<\tfrac1{10}$ and \eqref{e:pointwise} holds without
the last term. Suppose \m{\abs{r(h)}\ge\tfrac1{20}}. Since $10^{6}\delta\le10^{-6}/T\le10^{-6}$ and \m{\abs{\Omega}^{2}\le T},
\eqref{e:rbound} gives $\tfrac1{20}\le10^{-6}(G_0(h)+T)/T$, that is \m{G_0(h)\ge(5\cdot10^{4}-1)\,T}, and also
\m{\abs{r(h)}\le10^{-6}(G_0(h)+T)}. By Lemma~\ref{l:minorsK} and \m{\abs{e^{r}-1}\le e^{\abs{r}}} the left side of
\eqref{e:pointwise} is at most
\[\exp\bigl(-G_0(h)+\abs{r(h)}\bigr)\le\exp\bigl(-(1-10^{-6})\,G_0(h)+10^{-6}T\bigr)\le e^{-49998\,T}.\]
It remains to bound $c_\Lambda$ from below. Apply \inp{10}(b) to the compact group \m{\mathfrak{g}=\Ks{\varepsilon}} with its Haar
measure, the inclusion~$j$ and \m{R=1}: every \m{h\in\Ks{\varepsilon}} has \m{\op{h}=1}, and
\m{\dsc{l}{h}=\wm{\Lambda}{h}} by \inp{8}(b), so
\begin{align*}
1&\le\bigl(2000\,(\abs{\Omega}+L+1)^{6}\bigr)^{2\abs{\Omega}^{2}}\cdot4\,c_\Lambda,\\
\log c_\Lambda&\ge-\log4-2\abs{\Omega}^{2}\log2000-12\abs{\Omega}^{2}\log(\abs{\Omega}+L+1)\ge-30\,T,
\end{align*}
by the three lower bounds for~$T$. Hence $e^{-49998T}\le e^{-100T}c_\Lambda$, which proves \eqref{e:pointwise}, and the
theorem with \m{\kappa_{\mathrm{G}}=e^{-q_0}}.
\end{proof}
\leanline{\lean{dsCoeffgaussianshape}\sep \lean{gapWeightedNormSq}}

\sect{Two-factor data: the integrand near the block-diagonal matrices}{s:twofactor}
Norms without mention are Hilbert--Schmidt norms in this section too. The aim is Proposition~\ref{p:nearone}: where
both $j_1(\sigma)$ and $j_2(\sigma)$ are close to block-diagonal matrices, the point $\sigma$ can be moved, without
changing the integrand, to a point where both are close to~$1$.

\begin{lemma}[a trace inequality]\label{l:trace}
Let $\Omega$ be a finite set, $\varepsilon\colon\Omega\to\{\pm1\}$ signs, \m{g\in\Us{\varepsilon}}, let $c$ and $C$ be
real numbers and $v$, $w$ integer vectors on~$\Omega$ such that, for all \m{i,i'\in\Omega},
\[c\,v_i^{2}\le w_i,\qquad\text{and}\qquad w_i\le C\ \text{ if }\ \varepsilon_i\ne\varepsilon_{i'}.\]
Then $c\,\hs{\dg{v}(g-1)}^{2}+2\re\tr\bigl(\dg{w}(g-1)\bigr)\le2C\,\hs{g-\bd{\varepsilon}{g}}^{2}$.
\end{lemma}
\begin{proof}
Since $g^{-1}=\dg{\varepsilon}g^{*}\dg{\varepsilon}$ and $\dg{\varepsilon}$ commutes with $\dg{w}$, the number
$\tr(\dg{w}(g^{-1}-1))=\tr(\dg{w}(g^{*}-1))$ is the complex conjugate of \m{\tr(\dg{w}(g-1))}. As
$(g-1)(1-g^{-1})=(g-1)+(g^{-1}-1)$, this gives
$2\re\tr(\dg{w}(g-1))=\re\tr\bigl(\dg{w}(g-1)(1-g^{-1})\bigr)=\sum_{i,i'}w_i\re\bigl((g-1)_{ii'}(1-g^{-1})_{i'i}\bigr)$. Now
$(1-g^{-1})_{i'i}=\delta_{ii'}-\varepsilon_i\varepsilon_{i'}\overline{g_{ii'}}$: for a compact pair the term is
\m{-w_i\abs{(g-1)_{ii'}}^{2}}, for a non-compact pair it is \m{w_i\abs{g_{ii'}}^{2}}. Hence the left side of the claim is
\[\sum_{\varepsilon_i=\varepsilon_{i'}}(c\,v_i^{2}-w_i)\,\abs{(g-1)_{ii'}}^{2}+\sum_{\varepsilon_i\ne\varepsilon_{i'}}(c\,v_i^{2}+w_i)\,\abs{g_{ii'}}^{2}.\]
The first sum is at most~$0$, and in the second \m{c\,v_i^{2}+w_i\le2w_i\le2C}, while
$\hs{g-\bd{\varepsilon}{g}}^{2}=\sum_{\varepsilon_i\ne\varepsilon_{i'}}\abs{g_{ii'}}^{2}$.
\end{proof}
\leanline{\lean{normsqaddretracele}}

\begin{lemma}[a unitary block]\label{l:block}
Let $\Omega$, $\Xi$ and $\Xi'$ be finite sets, $\varepsilon\colon\Omega\to\{\pm1\}$ signs, $u$ a \m{\Xi\times\Xi}
matrix, and \m{e\colon\Xi\sqcup\Xi'\to\Omega} a bijection such that \m{\varepsilon_{e(\xi)}=1} for every \m{\xi\in\Xi}.
Suppose that \m{u^{*}u=1}. Then there is \m{g\in\Us{\varepsilon}} with $g-1=\inc{e|_\Xi}\,(u-1)\,\inc{e|_\Xi}^{*}$.
\end{lemma}
\begin{proof}
Put \m{Q=\inc{e|_\Xi}}. Then \m{Q^{*}Q=1} because $e$ is injective, and \m{\dg{\varepsilon}Q=Q} because \m{\varepsilon=1}
on~$e(\Xi)$. For \m{g=1+Q(u-1)Q^{*}} this gives
$g^{*}\dg{\varepsilon}g=\dg{\varepsilon}+Q\bigl((u^{*}-1)+(u-1)+(u^{*}-1)(u-1)\bigr)Q^{*}=\dg{\varepsilon}+Q(u^{*}u-1)Q^{*}=\dg{\varepsilon}$.
\end{proof}
\leanline{\lean{existsunitarysuboneeqinclMatrix}}

\begin{lemma}[a block of weight zero does not change the coefficient]\label{l:zeroblock}
Let $\Omega$ and $\Xi$ be finite sets, $\varepsilon\colon\Omega\to\{\pm1\}$ signs, \m{l\in\ZZ^{\Omega}}, $u$ a
\m{\Xi\times\Xi} matrix, and \m{\iota\colon\Xi\to\Omega} a map with \m{\varepsilon_{\iota(\xi)}=1} and
\m{\lowk{l}_{\iota(\xi)}=0} for every \m{\xi\in\Xi}. Let $h,g,\tilde g\in\Us{\varepsilon}$ with
$h-1=\inc{\iota}\,u\,\inc{\iota}^{*}$ and \m{\tilde g=h\,g}. Then
\begin{enumerate}[label=\textup{(\arabic*)},leftmargin=2.3em,itemsep=1pt]
\item \m{\dg{\lowk{l}}\,(h-1)=0};
\item \m{\dsc{l}{\tilde g}=\dsc{l}{g}};
\item $\hs{\tilde g-1}^{2}+2\re\tr(\tilde g-1)\le2\,\hs{g-\bd{\varepsilon}{g}}^{2}$.
\end{enumerate}
\end{lemma}
\begin{proof}
The matrix \m{h-1} is supported on \m{\iota(\Xi)\times\iota(\Xi)}, where \m{\lowk{l}=0} and \m{\varepsilon=1}. This gives (1),
and it shows that $h$ commutes with $\dg{\varepsilon}$; so \m{h\in\Ks{\varepsilon}}, $h$ is unitary, and
$\bd{\varepsilon}{hg}=h\,\bd{\varepsilon}{g}$.

For (2) we show $\wm{\lowk{l}}{h\,y}=\wm{\lowk{l}}{y}$ for every invertible~$y$. Let \m{\tau\in\ZZ} and
\m{\Pi=\{\lowk{l}\ge\tau\}}. An entry $h_{ii'}$ with \m{i\ne i'} vanishes unless \m{i,i'\in\iota(\Xi)}, and then
\m{\lowk{l}_i=0=\lowk{l}_{i'}}; so no such entry joins $\Pi$ to its complement, and $\pmn{\Pi}{hy}=\pmn{\Pi}{h}\pmn{\Pi}{y}$
by \inp{4}(a). If \m{\tau\le0} then \m{\iota(\Xi)\subseteq\Pi} and \m{\pmn{\Pi}{h}=\det h}; if \m{\tau>0} then $\Pi$ does
not meet $\iota(\Xi)$ and \m{\pmn{\Pi}{h}=1}. In the product of Definition~\ref{d:wm}, taken with \m{c_1\le0\le c_2}, the
factors \m{\det h} are \m{-c_1} in number and cancel against \m{(\det h)^{c_1}}. Now
$\bd{\varepsilon}{\tilde g}^{-1}\tilde g=\bd{\varepsilon}{g}^{-1}g$, and in the integral that defines
$\dsc{l}{\tilde g}=\fj{l}{h\,\bd{\varepsilon}{g}}{\bd{\varepsilon}{g}^{-1}g}$ the factor
$\wm{\lowk{l}}{h\,\bd{\varepsilon}{g}\,h_1}$, with $h_1$ the variable of integration, equals $\wm{\lowk{l}}{\bd{\varepsilon}{g}\,h_1}$. This is~(2).

For (3) apply Lemma~\ref{l:trace} to~$\tilde g$ with \m{c=C=1} and \m{v=w=1}; its right side is
$2\hs{\tilde g-\bd{\varepsilon}{\tilde g}}^{2}=2\hs{h(g-\bd{\varepsilon}{g})}^{2}=2\hs{g-\bd{\varepsilon}{g}}^{2}$ by~\inp{1}.
\end{proof}
\leanline{\lean{dsCoeffmuleqofzeroweightblock}}

\begin{lemma}[decay away from the block-diagonal matrices]\label{l:fourth}
Let $\Omega$ be a finite set, $\varepsilon\colon\Omega\to\{\pm1\}$ signs, \m{l\in\ZZ^{\Omega}} and $A$ a natural number
such that $l$ has large gap~$A$ for~$\varepsilon$, and let \m{g\in\Us{\varepsilon}}. Then $\bd{\varepsilon}{g}$ is
invertible and
\[\abs{\dsc{l}{g}}^{4}\le\bigl(1+\hs{g-\bd{\varepsilon}{g}}^{2}\bigr)^{-A}.\]
\end{lemma}
\begin{proof}
Write \m{J=\dg{\varepsilon}}, $g_0=\bd{\varepsilon}{g}=\tfrac12(g+JgJ)$ and \m{X=g-g_0=\tfrac12(g-JgJ)}. Expanding and
using \m{g^{*}Jg=J} and \m{J^{2}=1} one finds $g_0^{*}g_0=\tfrac14(g^{*}g+Jg^{*}gJ+2)$ and
$X^{*}X=\tfrac14(g^{*}g+Jg^{*}gJ-2)$, so
\[g_0^{*}\,g_0=1+X^{*}X.\]
Hence $g_0$ is invertible, and $\abs{\det g_0}^{2}=\det(1+X^{*}X)\ge1+\tr(X^{*}X)=1+\hs{X}^{2}$, the eigenvalues of
$X^{*}X$ being non-negative. By \inp{9} with the natural number \m{A\le500\abs{\Omega}^{3}+A},
$\abs{\dsc{l}{g}}^{2}\le\abs{\det g_0}^{-A}$; squaring gives the claim.
\end{proof}
\leanline{\lean{normdsCoeffpowfourle}}

\begin{setting}[two-factor data]\label{set:two}
$\Omega$, $\Omega'$ and $\Xi$ are finite sets, $\mathfrak{g}$ is a group, $\varepsilon\colon\Omega\to\{\pm1\}$ and
$\varepsilon'\colon\Omega'\to\{\pm1\}$ are signs, \m{l\in\ZZ^{\Omega}}, \m{l'\in\ZZ^{\Omega'}},
$j=(j_1,j_2)\colon \mathfrak{g}\to\Us{\varepsilon}\times\Us{\varepsilon'}$ is a homomorphism of groups, \m{e'\colon\Xi\to\Omega'} is
a map, $B$ is a real number, $M$ a natural number, $e\colon\Xi\sqcup\{0,\dots,M-1\}\to\Omega$ a bijection and
\m{\phi\colon \mathfrak{g}\to\CC} a function, and $\twof{l}{l'}{j}{e'}{B}{e}{\phi}$ holds.
\end{setting}

\begin{lemma}[the phase of the logarithms]\label{l:phaselog}
In Setting~\ref{set:two}, let \m{\sigma\in \mathfrak{g}} with \m{\hs{j_1(\sigma)-1}\le\tfrac12} and
\m{\hs{j_2(\sigma)-1}\le\tfrac12}. Then
\[\phs{l}{l'}{\bigl(\lgo{j_1(\sigma)-1},\ \lgo{j_2(\sigma)-1}\bigr)}=\phi(\sigma).\]
\end{lemma}
\begin{proof}
Pairs of matrices are added and multiplied componentwise, and for \m{\sigma'\in \mathfrak{g}} we write
$Y_{\sigma'}=(j_1(\sigma')-1,\ j_2(\sigma')-1)$. Since $j$ is a homomorphism,
$Y_\sigma(1+Y_{\sigma'})=Y_{\sigma\sigma'}-Y_{\sigma'}$. We claim that for every natural number~$n$ and every
\m{\sigma'\in \mathfrak{g}}
\[\phs{l}{l'}{Y_\sigma^{\,n}\,Y_\sigma\,(1+Y_{\sigma'})}=\phi(\sigma)\ \text{ if }n=0,\qquad\text{and}\quad=0\ \text{ if }n\ge1.\]
For \m{n=0} the left side is $\phi(\sigma\sigma')-\phi(\sigma')=\phi(\sigma)$ by (T1), the linearity of the phase,
and~(T3). For the step, $Y_\sigma^{\,n+1}Y_\sigma(1+Y_{\sigma'})=Y_\sigma^{\,n}Y_\sigma(1+Y_{\sigma\sigma'})-Y_\sigma^{\,n}Y_\sigma(1+Y_{\sigma'})$,
and the two terms have the same phase by the claim for~$n$. With \m{\sigma'=1} the claim says that the phase of
$Y_\sigma$ is $\phi(\sigma)$ and the phase of every higher power of~$Y_\sigma$ is~$0$. Now for \m{s\in[0,1]} the
Neumann series $(1+sY_\sigma)^{-1}Y_\sigma=\sum_{n\ge0}(-s)^{n}Y_\sigma^{\,n+1}$ converges in both components, the norms
being at most~$\tfrac12$; so the phase of \m{(1+sY_\sigma)^{-1}Y_\sigma} is $\phi(\sigma)$ for every~$s$. Integrating over~$s$
and using \inp{3} for each of the two traces gives the lemma.
\end{proof}
\leanline{\lean{phaselogOneAdd}\sep \lean{phase}\sep \lean{subOnePair}}

\begin{proposition}[moving a point near the identity]\label{p:nearone}
In Setting~\ref{set:two}, let $\delta$ be a real number with \m{1\le B}, \m{0<\delta} and \m{\delta\le\tfrac1{100}}, and
let \m{\sigma\in \mathfrak{g}} be such that
\[\hs{j_1(\sigma)-\bd{\varepsilon}{j_1(\sigma)}}^{2}+\hs{j_2(\sigma)-\bd{\varepsilon'}{j_2(\sigma)}}^{2}\le\delta^{2}.\]
Then there is \m{\sigma'\in \mathfrak{g}} with
\[\lapint{l}{l'}{j}{\phi}{\sigma}=\lapint{l}{l'}{j}{\phi}{\sigma'},\qquad \hs{j_1(\sigma')-1}\le10\,\delta,\qquad
\hs{j_2(\sigma')-1}\le10\,\delta .\]
\end{proposition}
\begin{proof}
Write \m{g=j_1(\sigma)}, \m{g'=j_2(\sigma)}, \m{\Lambda=\lowk{l}}, \m{\Lambda'=\lowk{l'}}, and let $\chi$ be the integer vector
on~$\Omega$ that is $0$ on $e(\Xi)$ and $1$ elsewhere. The matrix \m{c_0=\dg{\chi}(g-1)} consists of the rows of
\m{g-1} outside~$e(\Xi)$.

\emph{Step 1: \m{\hs{c_0}^{2}\le6\,\delta^{2}}.} We use Lemma~\ref{l:trace} three times; each time
we give its data \m{(c,C,v,w)}. For $g$ with \m{(B,3B,\chi,-\Lambda)}: the hypothesis holds by (S2) on $e(\Xi)$ and by (S3)
elsewhere, and the lemma gives
\[B\hs{c_0}^{2}-2\re\tr(\dg{\Lambda}(g-1))\le6B\,\hs{g-\bd{\varepsilon}{g}}^{2}.\]
For $g'$ with \m{(0,3B,0,\Lambda')}: the
hypothesis is (T6), and $2\re\tr(\dg{\Lambda'}(g'-1))\le6B\,\hs{g'-\bd{\varepsilon'}{g'}}^{2}$. By (T1) and (T4) the two
real parts are equal. Adding, \m{B\hs{c_0}^{2}\le6B\,\delta^{2}}, and \m{B\ge1}. For $g$ with \m{(-1,0,\chi,-\chi)} the
lemma gives
\begin{equation}\label{e:trc}
-2\re\tr c_0\le\hs{c_0}^{2}.
\end{equation}

\emph{Step 2: a unitary matrix.} We write \m{\Omega\times\Omega} matrices in blocks for the partition of $\Omega$ into
$e(\Xi)$, identified with~$\Xi$ by~$e$, and its complement~$\Omega_1$:
\[g=\begin{pmatrix}g_{11}&g_{12}\\ g_{21}&g_{22}\end{pmatrix},\qquad
c_0=\begin{pmatrix}0&0\\ g_{21}&g_{22}-1\end{pmatrix},\qquad
\dg{\varepsilon}=\begin{pmatrix}1&0\\0&J_2\end{pmatrix}.\]
By Step~1, \m{\hs{\tfrac12c_0}\le\tfrac12}; so by \inp{2} the matrix \m{W=(1+\tfrac12c_0)^{-1}} exists with
\m{\hs{W-1}\le\hs{c_0}}, the block \m{1+g_{22}} is invertible, and
\[W=\begin{pmatrix}1&0\\-(1+g_{22})^{-1}g_{21}&2\,(1+g_{22})^{-1}\end{pmatrix}.\]
Let $u$ be the \m{\Xi\times\Xi} matrix with $u^{*}=g_{11}-g_{12}(1+g_{22})^{-1}g_{21}$. It is unitary: for
\m{v\in\CC^{\Xi}} put $v'=-(1+g_{22})^{-1}g_{21}v\in\CC^{\Omega_1}$; then $g$ maps the vector \m{(v,v')} to
\m{(u^{*}v,\,-v')}, and since $g$ preserves the form \m{\nrm{v}^{2}+v'^{*}J_2v'}, we get \m{\nrm{u^{*}v}=\nrm{v}}.

\emph{Step 3: the new point.} By (T7) there is \m{\sigma_0\in \mathfrak{g}} with
$j_1(\sigma_0)-1=\inc{e|_\Xi}(u-1)\inc{e|_\Xi}^{*}$ and $j_2(\sigma_0)-1=\inc{e'}(u-1)\inc{e'}^{*}$; put
\m{\sigma'=\sigma_0\sigma}, \m{T=j_1(\sigma')-1} and \m{T'=j_2(\sigma')-1}. Lemma~\ref{l:zeroblock} applies to
$(\iota,u,h,g,\tilde g)=(e|_\Xi,\,u-1,\,j_1(\sigma_0),\,g,\,j_1(\sigma'))$ by (S1) and (S2), and to
$(e',\,u-1,\,j_2(\sigma_0),\,g',\,j_2(\sigma'))$ by (T5). Its part~(1) and (T1) give \m{\phi(\sigma_0)=0}, so
\m{\phi(\sigma')=\phi(\sigma)} by (T3); with part~(2) for both factors, the integrand has the same value at $\sigma'$
and at~$\sigma$. Part~(3) and (T2) give
\begin{equation}\label{e:TT}
\hs{T}^{2}+2\re\tr T\le2\,\hs{g-\bd{\varepsilon}{g}}^{2},\qquad \hs{T'}^{2}+2\re\tr T\le2\,\hs{g'-\bd{\varepsilon'}{g'}}^{2}.
\end{equation}

\emph{Step 4: the trace of~$T$.} In blocks \m{j_1(\sigma_0)} is $u$ on the first block and $1$ on the second, and
$u\,g_{11}=1+u\,g_{12}(1+g_{22})^{-1}g_{21}$ by the definition of~$u$. Hence
\begin{align*}
T-c_0&=\begin{pmatrix}u\,g_{12}(1+g_{22})^{-1}g_{21}&u\,g_{12}\\0&0\end{pmatrix},\\
\tr\bigl((T-c_0)(W-1)\bigr)&=-\tr\bigl(u\,g_{12}(1+g_{22})^{-1}g_{21}\bigr)=\tr c_0-\tr T,
\end{align*}
and by \inp{1}, $\re\tr c_0-\re\tr T\le(\hs{T}+\hs{c_0})\,\hs{c_0}$.

\emph{Step 5: the bounds.} By \eqref{e:trc} and Step~4, $-2\re\tr T\le\hs{c_0}^{2}+2(\hs{T}+\hs{c_0})\hs{c_0}$. With
\eqref{e:TT} and the hypothesis on~$\sigma$ this gives $(\hs{T}-\hs{c_0})^{2}\le2\delta^{2}+4\hs{c_0}^{2}\le26\,\delta^{2}$, so
$\hs{T}\le(\sqrt6+\sqrt{26})\,\delta<7.6\,\delta$, and then
$\hs{T'}^{2}\le2\delta^{2}+6\delta^{2}+2\,(7.6+2.5)\cdot2.5\,\delta^{2}<59\,\delta^{2}$. Both norms are at most~$10\,\delta$.
\end{proof}
\leanline{\lean{existsnearone}\sep \lean{laplaceIntegrand}}

The unitary $u$ of Step~2 is the transfer matrix of the matrix~$g$ with its second block closed by the
reflection~$-1$; the formal proof writes it as $1+\inc{e|_\Xi}^{*}W^{*}(g-1)^{*}\inc{e|_\Xi}$ and verifies its unitarity,
and the identity of Step~4, by linear combinations of the defining relations.

\sect{Laplace's method: proof of Theorem 2}{s:laplace}
Throughout the proof the data are those of Theorem~2, and we abbreviate
\[\mathsf{n}=(\abs{\Omega}+\abs{\Omega'})^{2},\qquad T=\mathsf{n}\,\log\beta,\qquad \gamma(\sigma)=\frac{2}{\op{j_1(\sigma)}}\qquad(\sigma\in \mathfrak{g}).\]
Since \m{\beta\ge10^{12}} we have \m{\log\beta\ge1}, so \m{T\ge\mathsf{n}\ge4}. The map $j_1$ is continuous, being proper.
By (T4), $\abs{\lapint{l}{l'}{j}{\phi}{\sigma}}=\abs{\dsc{l}{j_1(\sigma)}}\,\abs{\dsc{l'}{j_2(\sigma)}}$.

\begin{proof}[Proof of Theorem 2]
\emph{Step 1: decay.} For every natural number \m{c\le18\,\mathsf{n}} and every \m{\sigma\in \mathfrak{g}},
\begin{equation}\label{e:decay}
\abs{\lapint{l}{l'}{j}{\phi}{\sigma}}\le\gamma(\sigma)^{c}.
\end{equation}
Indeed \m{\abs{\dsc{l'}{\,\cdot\,}}\le1} and $\abs{\dsc{l}{g}}\le(2/\op{g})^{c}$ by \inp{9}, since \m{c\le18\,\mathsf{n}\le A}
by hypothesis~(d), and $A$ is at most the gap of~$l$.

\emph{Step 2: the real part.} For every \m{\sigma\in \mathfrak{g}},
\begin{equation}\label{e:cone}
4\,\abs{\lapint{l}{l'}{j}{\phi}{\sigma}}\ \le\ 8\,\Bigl(\re\lapint{l}{l'}{j}{\phi}{\sigma}+16\,\beta^{-50\mathsf{n}}\,\gamma(\sigma)^{8\mathsf{n}}\Bigr).
\end{equation}
Note first that an inequality $4\abs{z}\le8(\re z+16\,\epsilon)$ for a complex number~$z$ follows from
\m{\abs{z}\le\epsilon}, because \m{\re z\ge-\abs{z}}. If $\gamma(\sigma)^{8\mathsf{n}}\le\beta^{-50\mathsf{n}}$, then
$\beta^{-50\mathsf{n}}\gamma(\sigma)^{8\mathsf{n}}\ge\gamma(\sigma)^{16\mathsf{n}}\ge\abs{\lapint{l}{l'}{j}{\phi}{\sigma}}$ by \eqref{e:decay}, and
\eqref{e:cone} holds. Suppose now that $\gamma(\sigma)^{8\mathsf{n}}\ge\beta^{-50\mathsf{n}}$. Then
$\beta^{-50\mathsf{n}}\gamma(\sigma)^{8\mathsf{n}}\ge\beta^{-100\mathsf{n}}=e^{-100T}$, and it suffices to show
\begin{equation}\label{e:cone2}
4\,\abs{\lapint{l}{l'}{j}{\phi}{\sigma}}\le8\,\bigl(\re\lapint{l}{l'}{j}{\phi}{\sigma}+16\,e^{-100T}\bigr).
\end{equation}
Put \m{\delta=10^{-13}/T}, so \m{0<\delta\le\tfrac1{100}}, and let
\[\operatorname{dist}(\sigma)^{2}=\hs{j_1(\sigma)-\bd{\varepsilon}{j_1(\sigma)}}^{2}+\hs{j_2(\sigma)-\bd{\varepsilon'}{j_2(\sigma)}}^{2}\]
be the squared distance of $j(\sigma)$ from the block-diagonal matrices.

\emph{Far from the block-diagonal matrices: $\operatorname{dist}(\sigma)^{2}\ge\delta^{2}$.} By Lemma~\ref{l:fourth} for both factors,
$\abs{\lapint{l}{l'}{j}{\phi}{\sigma}}^{4}\le(1+\operatorname{dist}(\sigma)^{2})^{-A}\le(1+\delta^{2})^{-A}$. As $\log(1+\delta^{2})\ge\tfrac12\delta^{2}$ and
\m{A\ge10^{30}T^{3}} by~(d), $A\log(1+\delta^{2})\ge5000\,T\ge400\,T$. So $\abs{\lapint{l}{l'}{j}{\phi}{\sigma}}\le e^{-100T}$, which gives
\eqref{e:cone2}.

\emph{Near them: $\operatorname{dist}(\sigma)^{2}\le\delta^{2}$.} By (a), \m{B\ge1} and Proposition~\ref{p:nearone} there is \m{\sigma'\in \mathfrak{g}} with
$\lapint{l}{l'}{j}{\phi}{\sigma}=\lapint{l}{l'}{j}{\phi}{\sigma'}$ and $\hs{g_1-1},\hs{g_2-1}\le10\,\delta$, where \m{g_1=j_1(\sigma')} and \m{g_2=j_2(\sigma')}.
Theorem~1 applies to \m{(l,g_1)} with the numbers \m{(10\delta,T,L,A)}: indeed \m{10\delta\cdot T=10^{-12}},
\m{\abs{\Omega}^{2}\le T}, and $\abs{\Omega}^{2}\log(\abs{\Omega}+L+1)\le\mathsf{n}\log\beta$ because
$\abs{\Omega}+L+1\le\abs{\Omega}+\abs{\Omega'}+L\le\beta$. It applies in the same way to \m{(-l',g_2)}, for the
signs~$\varepsilon'$: the vector $-l'$ has large gap~$A$ and entries bounded by~$L$. With
\[E_1=\tr\bigl(\dg{\lowk{l}}\lgo{g_1-1}\bigr),\qquad E_2=\tr\bigl(\dg{\lowk{l'}}\lgo{g_2-1}\bigr),\]
and \m{\lowk{-l'}=-\lowk{l'}}, $\dsc{-l'}{g_2}=\overline{\dsc{l'}{g_2}}$ from \inp{8}(c), we get
\m{\kappa_1,\kappa_2\in[0,1]} with
\[\Bigabs{\frac{\dsc{l}{g_1}}{e^{E_1}}-\kappa_1}\le\frac{\kappa_1}{10}+e^{-100T},\qquad
\Bigabs{\frac{\overline{\dsc{l'}{g_2}}}{e^{-E_2}}-\kappa_2}\le\frac{\kappa_2}{10}+e^{-100T}.\]
By Lemma~\ref{l:phaselog}, \m{\phi(\sigma')=E_1-E_2}: the phase of the integrand cancels the two phases of the
coefficients. So the complex number $z=\lapint{l}{l'}{j}{\phi}{\sigma'}$ is \m{(\kappa_1+z_1)(\kappa_2+z_2)} with complex numbers
$\abs{z_n}\le\tfrac1{10}\kappa_n+\epsilon$, \m{\epsilon=e^{-100T}\le1}. Expanding,
$\re z\ge0.79\,\kappa_1\kappa_2-3.2\,\epsilon$ and $\abs{z}\le1.21\,\kappa_1\kappa_2+3.2\,\epsilon$, whence
$8(\re z+16\epsilon)\ge6.32\,\kappa_1\kappa_2+102\,\epsilon\ge4\abs{z}$. This is \eqref{e:cone2}, and Step~2 is
complete.

\emph{Step 3: integrability.} By \inp{11}(b) for the proper homomorphism~$j_1$, the point \m{y=1} and the exponent
$8\,\mathsf{n}>8\abs{\Omega}^{2}$, the number $\mathsf{m}=\nu\{\sigma:\op{j_1(\sigma)-1}\le1\}$ is finite, and
\begin{equation}\label{e:gammaint}
\int_{\mathfrak{g}}\gamma(\sigma)^{8\mathsf{n}}\,d\nu(\sigma)\le16^{\,8\mathsf{n}}\,\mathsf{m} .
\end{equation}
The function $\lapint{l}{l'}{j}{\phi}{\,\cdot\,}$ is continuous: the coefficients are continuous by \inp{8}(a), $j_1$ and $j_2$ are continuous,
and $\phi$ is continuous by~(T1). By \eqref{e:decay} it is dominated by~$\gamma^{8\mathsf{n}}$, so it is $\nu$-integrable.

\emph{Step 4: a small ball.} Let $U_0$ be the set of the \m{\sigma\in \mathfrak{g}} with $\op{j_1(\sigma)-1}\le(200\beta^{5})^{-1}$
and $\op{j_2(\sigma)-1}\le(200\beta^{5})^{-1}$. By \inp{10}(a), with \m{\beta_0=\beta}, both coefficients have absolute
value at least~$\tfrac12$ on~$U_0$, so $4\abs{\lapint{l}{l'}{j}{\phi}{\sigma}}\ge1$ for \m{\sigma\in U_0}. The right side of \eqref{e:cone} is therefore a
non-negative integrable function that is at least~$1$ on~$U_0$, and Markov's inequality with \eqref{e:gammaint} gives
\[\nu(U_0)\le8\,\re\!\int_{\mathfrak{g}}\lapint{l}{l'}{j}{\phi}{\sigma}\,d\nu(\sigma)+128\,\beta^{-50\mathsf{n}}\!\int_{\mathfrak{g}}\gamma^{8\mathsf{n}}\,d\nu
\le8\,\mathsf{J}+128\,\beta^{-50\mathsf{n}}\,16^{\,8\mathsf{n}}\,\mathsf{m},\]
where $\mathsf{J}=\bigabs{\int_{\mathfrak{g}}\lapint{l}{l'}{j}{\phi}{\sigma}\,d\nu(\sigma)}$.

\emph{Step 5: covering.} The map $\sigma\mapsto j_1(\sigma)\oplus j_2(\sigma)$ is a homomorphism from $\mathfrak{g}$ to the unitary
group of the signs \m{\varepsilon\oplus\varepsilon'} on \m{\Omega\sqcup\Omega'}, and $U_0$ is the set where it is within
operator distance \m{(200\beta^{5})^{-1}} of~$1$. By \inp{11}(a) with \m{y=1}, \m{D=200\beta^{5}}, \m{R=\beta+1} and
\m{K=\beta^{8}\ge3DR^{2}},
\[\nu\bigl\{\sigma:\op{j_1(\sigma)\oplus j_2(\sigma)}\le\beta+1\bigr\}\le\beta^{16\mathsf{n}}\,\nu(U_0).\]
The set on the left contains $\{\sigma:\op{j_1(\sigma)-1}\le1\}$, because there \m{\op{j_1(\sigma)}\le2} and
\m{\op{j_2(\sigma)}\le\beta+1} by~(e). Hence
\[\mathsf{m}\le8\,\beta^{16\mathsf{n}}\,\mathsf{J}+128\cdot16^{\,8\mathsf{n}}\,\beta^{-34\mathsf{n}}\,\mathsf{m}
\le\tfrac12\,\beta^{20\mathsf{n}}\,\mathsf{J}+\tfrac12\,\mathsf{m},\]
using \m{\beta\ge16} twice. Since $\mathsf{m}$ is finite, this is the theorem.
\end{proof}
\leanline{\lean{laplacemethod}\sep \lean{IsTwoFactor}\sep \lean{IsSteep}}

\sect{The restriction to a conjugate subgroup: proof of Theorem 3}{s:restriction}
The data are those of Theorem~3. Write \m{\Lambda=\lowk{l}} for~$\varepsilon$ and \m{\Lambda'=\lowk{l'}} for
\m{\varepsilon\oplus1}; since $l'$ has a gap, \m{\lowk{-l'}=-\Lambda'} by \inp{8}(c). Operator norms are used.

\begin{proof}[Proof of Theorem 3]
\emph{Nodes.} Define integers \m{\tau_0,\dots,\tau_{2M}} by $\tau_n=-\Lambda'_{\tilde e(M-n)}$ for \m{0\le n\le M} and
\m{\tau_{M+1+n}=\Lambda_{e(n)}} for \m{0\le n<M}. By (S3) for $l'$ the first \m{M+1} of them decrease strictly and lie in
\m{[B,3B]}; by (S3) for~$l$ the last $M$ decrease strictly and lie in \m{[-3B,-B]}. So the sequence is strictly
decreasing, and \m{\abs{\tau_n-\tau_{n'}}\le6B} for all \m{n,n'}. Let \m{c_0,\dots,c_{2M}} be the weights of \inp{5} for these
\m{2M+1} nodes. For \m{0\le a\le M} put
\[\begin{gathered}d_a=\frac{(-1)^{a}}{c_a}\prod_{n<M}(\tau_a-\tau_{M+1+n}),\qquad
\pi_a(\tau)=d_a\!\!\prod_{n\le M,\,n\ne a}\!\!(\tau-\tau_n),\\
\rho_{na}=\frac{\pi_a(\tau_{M+1+n})}{c_{M+1+n}}\quad(n<M).\end{gathered}\]
The polynomial $\pi_a$ has degree~$M$ and leading coefficient~$d_a$, and by the formula for~$c_a^{2}$,
$\pi_a(\tau_{a'})/c_{a'}=\delta_{aa'}$ for \m{a,a'\le M}. For \m{a,a'\le M} the product \m{\pi_a\pi_{a'}} has degree $2M$ and
leading coefficient \m{d_ad_{a'}}, so \inp{5} gives $\sum_{n\le2M}(-1)^{n}\pi_a(\tau_n)\pi_{a'}(\tau_n)/c_n^{2}=d_ad_{a'}$. The terms
with \m{n\le M} add up to \m{(-1)^{a}\delta_{aa'}}; in the others \m{(-1)^{M+1+n}=-(-1)^{n}}, because $M$ is even. Hence
\begin{equation}\label{e:rho}
\sum_{n<M}(-1)^{n}\rho_{na}\,\rho_{na'}+d_a\,d_{a'}=(-1)^{a}\,\delta_{aa'}\qquad(a,a'\le M).
\end{equation}

\emph{The element~$\omega$.} Let \m{(b_i)} be the standard basis of \m{\CC^{\Omega\sqcup\{*\}}}, and let $\omega$ be the
matrix with
\[\omega\,b_{\tilde e(\xi)}=b_{e(\xi)}\quad(\xi\in\Xi),\qquad
\omega\,b_{\tilde e(a)}=\sum_{n<M}\rho_{n,M-a}\,b_{e(n)}+d_{M-a}\,b_{*}\quad(0\le a\le M).\]
The signs \m{\varepsilon\oplus1} are $+1$ at $e(\xi)$, $\tilde e(\xi)$ and~$*$, \m{(-1)^{n}} at $e(n)$ and \m{(-1)^{a}} at
$\tilde e(a)$, by (S1) for $l$ and for~$l'$. So the product of the columns \m{\tilde e(a)} and \m{\tilde e(a')} of~$\omega$ for
the form \m{\dg{\varepsilon\oplus1}} is the left side of \eqref{e:rho} for \m{(M-a,M-a')}, that is
$(-1)^{M-a}\delta_{aa'}=(-1)^{a}\delta_{aa'}$; the columns \m{\tilde e(\xi)} are distinct basis vectors of sign~$+1$ that do
not occur in the other columns. Therefore $\omega^{*}\dg{\varepsilon\oplus1}\,\omega=\dg{\varepsilon\oplus1}$, and
$\omega\in\Us{\varepsilon\oplus1}$.

\emph{Property (i).} As \m{c_n\ge1} and \m{6B\ge1}, we have \m{\abs{d_a}\le(6B)^{M}} and \m{\abs{\rho_{na}}\le(6B)^{2M}}; so every
entry of~$\omega$ has absolute value at most \m{(6B)^{2M+1}}, and the operator norm of a square matrix of
size \m{\abs{\Omega}+1} is at most \m{(\abs{\Omega}+1)^{2}} times the largest entry.

\emph{Property (iii).} $\omega^{-1}(g\oplus1)\omega-1=\omega^{-1}\bigl((g\oplus1)-1\bigr)\omega$.

\emph{Property (ii).} The condition \m{\steep{l}{B}{e}} is a hypothesis. For (T1) we show that
\begin{equation}\label{e:intertwine}
\text{the rows of }\ \omega\,\dg{-\Lambda'}\,\omega^{-1}\ \text{ with index in }\Omega\text{ are those of }\ \bigl(\dg{\Lambda}\ \big|\ v\bigr)
\end{equation}
for a column~$v$ at~$*$. Then $\tr\bigl(\dg{-\Lambda'}\,\omega^{-1}(Z\oplus0)\,\omega\bigr)=\tr(\dg{\Lambda}Z)$ for every
\m{\Omega\times\Omega} matrix~$Z$, and with (iii) the phase $\phs{l}{-l'}{(g-1,\ j^{\omega}_2(g)-1)}$ vanishes for every~$g$, as
(T1) asks for the function \m{\phi=0}. To see \eqref{e:intertwine}, let \m{v_{e(\xi)}=0} and
$v_{e(n)}=-\prod_{a\le M}(\tau_{M+1+n}-\tau_a)/c_{M+1+n}$, and compare the entries of \m{\omega\,\dg{-\Lambda'}} and of
\m{(\dg{\Lambda}\,|\,v)\,\omega} in the rows of~$\Omega$. In the column \m{\tilde e(a)}, where \m{-\Lambda'=\tau_{M-a}}, and the row
$e(n)$, where \m{\Lambda=\tau_{M+1+n}}, the difference is
\[(\tau_{M-a}-\tau_{M+1+n})\,\rho_{n,M-a}-v_{e(n)}\,d_{M-a}=0,\]
because $(\tau_{a'}-\tau)\,\pi_{a'}(\tau)=-d_{a'}\prod_{a\le M}(\tau-\tau_a)$. In the rows $e(\xi)$ and in the columns
\m{\tilde e(\xi)} both sides vanish, as \m{\Lambda=0} on $e(\Xi)$ and \m{\Lambda'=0} on \m{\tilde e(\Xi)} by (S2).

Conditions (T3) and (T4) hold for the function~$0$, and (T2) because a conjugation keeps the trace. For (T5) and
(T6): the signs \m{\varepsilon\oplus1} are $+1$ on \m{\tilde e(\Xi)}, and \m{\lowk{-l'}=-\Lambda'} is $0$ there and lies in
\m{[B,3B]} at the other indices, by (S1) to (S3) for~$l'$. For (T7), let $u$ be a unitary \m{\Xi\times\Xi} matrix.
Lemma~\ref{l:block} gives \m{g\in\Us{\varepsilon}} with $g-1=\inc{e|_\Xi}(u-1)\inc{e|_\Xi}^{*}$. The matrix $\omega$ maps
$b_{\tilde e(\xi)}$ to $b_{e(\xi)}$, and its row at $e(\xi)$ is the unit row at $\tilde e(\xi)$; so
$\omega^{-1}\bigl((g-1)\oplus0\bigr)\omega=\inc{\tilde e|_\Xi}(u-1)\inc{\tilde e|_\Xi}^{*}$, which by (iii) is the second
equation of~(T7).
\end{proof}
\leanline{\lean{existsisTwoFactorrestriction}}

\sect{The Bessel graph: proofs of Theorems 4 and 5}{s:bessel}
\begin{lemma}[the Bessel graph and its measure]\label{l:gmeasure}
Let $\fs$ be a shape. Then $\bgr{\fs}$ is a closed subset of \m{\Gone{\fs}\times\Us{1}}, and the measure $\bms{\fs}$ is
positive on non-empty open subsets of $\bgr{\fs}$, finite on its compact subsets, and invariant under right
translations.
\end{lemma}
\begin{proof}
The condition \m{\emb{\fs}{h}^{-1}g\in\Uu{\fs}} says that certain entries of \m{\emb{\fs}{h}^{-1}g} are those of the
identity matrix, which is a closed condition on \m{(g,h)}. The chart $\bch{\fs}$ is a homeomorphism, and the product of
the Haar measure $\hU{\fs}$ with the Haar measure of the compact group $\Us{1}$ is positive on open sets and finite
on compact sets; so is its image. For the invariance let $\sigma'=(\emb{\fs}{h_0}u_0,\,h_0)\in\bgr{\fs}$. Then
\[\begin{aligned}\bch{\fs}(u,h)\cdot\sigma'&=\bigl(\emb{\fs}{h}\,u\,\emb{\fs}{h_0}\,u_0,\ h\,h_0\bigr)\\
&=\bch{\fs}\bigl(\emb{\fs}{h_0}^{-1}u\,\emb{\fs}{h_0}\cdot u_0,\ h\,h_0\bigr).\end{aligned}\]
In the chart, right translation by~$\sigma'$ is therefore the product of two maps: on $\Uu{\fs}$ the conjugation by
\m{\emb{\fs}{h_0}^{-1}\in\Qq{\fs}} followed by the right translation by~$u_0$, which preserve $\hU{\fs}$ by \inp{13}(c); and on
$\Us{1}$ the right translation by~$h_0$.
\end{proof}
\leanline{\lean{besselGraph}\sep \lean{besselGraphChart}\sep \lean{besselGraphMeasure}\sep
\lean{besselGraphMeasurefiniteOnCompacts}\sep \lean{besselGraphMeasureisMulRightInvariant}}

The closedness and the positivity on open sets are two unnamed instances of the formal text.

\begin{lemma}\label{l:fst}
Let $\fs$ be a shape and $\sigma=(\sigma_1,\sigma_2)\in\bgr{\fs}$. Then \m{\sigma_1\in\Qq{\fs}}.
\end{lemma}
\begin{proof}
$\sigma_1=\emb{\fs}{\sigma_2}\cdot\emb{\fs}{\sigma_2}^{-1}\sigma_1$ is the product of an element of $\Hh{\fs}$ and one of $\Uu{\fs}$.
\end{proof}
\leanline{\lean{besselGraphfstmembesselS}}

\begin{lemma}[unfolding along the Bessel graph]\label{l:unfold}
Let $\fs$ be a shape, \m{x\in\Gone{\fs}}, and let $f$ be a function from $\Gone{\fs}$ to \m{[0,\infty]} that is
measurable. Then
\[\int_{\Gone{\fs}}f(g)\,dg=\int_{\Gone{\fs}}\int_{\bgr{\fs}}f\bigl(x\,\sigma_1\,y\bigr)\,d\bms{\fs}(\sigma)\,d\hq{\fs}(y).\]
\end{lemma}
\begin{proof}
By the definition of $\bms{\fs}$ the inner integral is
\[\int_{\Us{1}}\int_{\Uu{\fs}}f(x\,\emb{\fs}{h}\,u\,y)\,d\hU{\fs}(u)\,dh.\]
By Tonelli's theorem and Definition~\ref{d:shape}(e), applied to \m{g\mapsto f(x\,\emb{\fs}{h}\,g)} for each~$h$, the right side is
$\int_{\Us{1}}\int_{\Gone{\fs}}f(x\,\emb{\fs}{h}\,g)\,dg\,dh$. The inner integral is \m{\int f(g)\,dg} by the left invariance of
the Haar measure, and $dh$ has mass one.
\end{proof}
\leanline{\lean{lintegralunfoldbesselGraphMeasure}}

\begin{lemma}[traces against the graph]\label{l:gtrace}
Let $\fs$ be a shape, $\sigma=(\sigma_1,\sigma_2)\in\bgr{\fs}$, and let $Z$ be an \m{\Ind{\fs}\times\Ind{\fs}} matrix that is
block upper triangular for the levels. Then $\tr\bigl(Z\,(\sigma_1-\emb{\fs}{\sigma_2})\bigr)=0$.
\end{lemma}
\begin{proof}
With $u=\emb{\fs}{\sigma_2}^{-1}\sigma_1\in\Uu{\fs}$ the trace is $\tr\bigl(Z\,\emb{\fs}{\sigma_2}\,(u-1)\bigr)$. The matrix
\m{Z\,\emb{\fs}{\sigma_2}} is block upper triangular, since $\emb{\fs}{\sigma_2}$ differs from~$1$ only inside the middle
level; and \m{(u-1)_{i'i}=0} whenever \m{\lev{i}\le\lev{i'}}. So in $\sum_{i,i'}(Z\,\emb{\fs}{\sigma_2})_{ii'}(u-1)_{i'i}$ every term
vanishes.
\end{proof}
\leanline{\lean{besselGraphtracemulsubeqzero}}

\begin{lemma}[properness]\label{l:proper}
Let $\fs$ be a shape, $\varepsilon$ signs on $\Ind{\fs}$, \m{\fp\in\frm{\fs}{\varepsilon}} a frame and \m{x\in\Gone{\fs}}. Then
the map \m{\sigma\mapsto\bph{x}{\sigma}_1} from $\bgr{\fs}$ to $\Us{\varepsilon}$ is proper.
\end{lemma}
\begin{proof}
It is the first projection of the closed subset $\bgr{\fs}$ of \m{\Gone{\fs}\times\Us{1}}, which is proper because
$\Us{1}$ is compact, followed by the homeomorphism $g\mapsto\fp^{-1}x^{-1}g\,x\,\fp$.
\end{proof}
\leanline{\lean{isProperMapfstbesselPairHom}\sep \lean{besselPairHom}\sep \lean{Frame}}

\begin{proof}[Proof of Theorem 4]
Write \m{\varrho=\rhs{\fs}} and \m{\Lambda=\lowk{l}}, and let \m{\fp_i} be the column of~$\fp$ at \m{i\in\Ind{\fs}}; these
columns are orthonormal, and $\fp_i^{*}\Jone{\fs}\fp_{i'}=\varepsilon_i\delta_{ii'}$.

\emph{Nodes and vectors.} Put \m{\tau_0=0} and \m{\tau_{n+1}=\Lambda_{e(n)}-a} for \m{0\le n\le2\varrho+1}. By (S3) these
\m{2\varrho+3} integers decrease strictly, and $\abs{\tau_{n+1}}\le3B+a\le2b^{2}$, so $\tau_{n+1}^{2}\le4b^{4}\le b^{5}$. Let
\m{c_0,\dots,c_{2\varrho+2}} be the weights of~\inp{5}, and define vectors of $\CC^{\Ind{\fs}}$
\[V^{(0)}=c_0\sum_{n\le2\varrho+1}\frac{1}{c_{n+1}}\,\fp_{e(n)},\qquad
V^{(k)}=\sum_{n\le2\varrho+1}\frac{\tau_{n+1}^{\,k}}{c_{n+1}}\,\fp_{e(n)}\quad(k\ge1).\]
Since \m{\varepsilon_{e(n)}=(-1)^{n}}, the products of these vectors for $\Jone{\fs}$ are, up to the factors~$c_0$,
the sums
\[\sum_{n\le2\varrho+1}(-1)^{n}\,\frac{\tau_{n+1}^{\,k+k'}}{c_{n+1}^{2}}.\]
Input~\inp{5} for the polynomial \m{\tau^{k+k'}}, of
degree at most \m{2\varrho+2}, evaluates them: the node \m{\tau_0=0} contributes only for \m{k+k'=0}, and the coefficient
of degree \m{2\varrho+2} is $1$ or~$0$. The result is, for \m{k,k'\ge1} with \m{k+k'\le2\varrho+2} and for
\m{1\le k''\le2\varrho+1},
\begin{equation}\label{e:gramV}\begin{gathered}
V^{(0)*}\Jone{\fs}V^{(0)}=1,\qquad V^{(0)*}\Jone{\fs}V^{(k'')}=0,\\
V^{(k)*}\Jone{\fs}V^{(k')}=\begin{cases}-1&\text{if }k+k'=2\varrho+2,\\0&\text{otherwise.}\end{cases}
\end{gathered}\end{equation}
The vectors \m{\fp_{e(\xi)}}, \m{\xi\in\Xi}, are orthonormal for $\Jone{\fs}$ with squares~$+1$ by (S1), and orthogonal to
all~$V^{(k)}$.

\emph{The point~$x$.} Consider the two families of vectors
\[\bigl(\fp_{e(\xi)}\ (\xi\in\Xi),\ V^{(0)},\ V^{(1)},\dots,V^{(\varrho+1)}\bigr)\qquad\text{and}\qquad
\bigl(x_{\tilde e(\xi)}\ (\xi\in\Xi),\ x_{\tilde e(*)},\ f_0,\dots,f_{\varrho}\bigr).\]
By \eqref{e:gramV} and Definition~\ref{d:shape} they have the same Gram matrix for $\Jone{\fs}$: the identity on the
first \m{\abs{\Xi}+1} vectors, and $-1$ at the last vector, $0$ elsewhere. The second family consists of distinct basis
vectors. The first is linearly independent, because its matrix of products with the vectors \m{\fp_{e(\xi)}},
\m{V^{(0)}}, $V^{(2\varrho+1)},\dots,V^{(\varrho+1)}$ is triangular with diagonal entries~$\pm1$, again by \eqref{e:gramV}. By
Witt's theorem \inp{6} there is \m{x\in\Gone{\fs}} with
\[x\,\fp_{e(\xi)}=x_{\tilde e(\xi)},\qquad x\,V^{(0)}=x_{\tilde e(*)},\qquad x\,V^{(a+1)}=f_a\quad(0\le a\le\varrho).\]

\emph{Its size.} Apply \inp{14} to~$x$ and the vectors \m{V_a=V^{(a+1)}}, \m{a\ge0}. For \m{a\le\varrho}, \m{V_a=x^{-1}f_a}. For
\m{a+a'\le2\varrho} the product \m{V_a^{*}\Jone{\fs}V_{a'}} is $-1$ if \m{a+a'=2\varrho} and $0$ otherwise, by \eqref{e:gramV}. As $\fp$
is unitary and \m{c_{n+1}\ge1}, for \m{a\le2\varrho} we have
$\nrm{V_a}^{2}\le(2\varrho+2)\,b^{5(2\varrho+1)}\le b^{10\varrho+6}$. So $\bsz{x}\le8(\varrho+1)\,b^{5\varrho+3}\le b^{5\varrho+4}\le b^{5\nsh{\fs}}$.

\emph{The conjugated weight matrix.} Put \m{q=x\,\fp}, so that $q^{*}\Jone{\fs}q=\dg{\varepsilon}$ and
$\bph{x}{\sigma}=(q^{-1}\sigma_1q,\ \sigma_2)$, and let \m{D=q\,\dg{\Lambda}\,q^{-1}}. Because \m{\Lambda_{e(\xi)}=0} and
\m{\Lambda_{e(n)}=\tau_{n+1}+a}, the matrix \m{\fp\,\dg{\Lambda}\,\fp^{-1}} kills $\fp_{e(\xi)}$, sends $V^{(k)}$ to
\m{a\,V^{(k)}+V^{(k+1)}} for \m{k\ge1}, and $V^{(0)}$ to \m{a\,V^{(0)}+c_0V^{(1)}}. Applying~$x$:
\begin{equation}\label{e:Dframe}
D\,x_{\tilde e(\xi)}=0,\qquad D\,x_{\tilde e(*)}=a\,x_{\tilde e(*)}+c_0\,f_0,\qquad D\,f_k=a\,f_k+f_{k+1}\quad(k<\varrho).
\end{equation}
Moreover $q^{-1}=\dg{\varepsilon}q^{*}\Jone{\fs}$, so $\Jone{\fs}D=\Jone{\fs}q\,\dg{\Lambda}\dg{\varepsilon}\,q^{*}\Jone{\fs}$ is hermitian:
\m{\Jone{\fs}D=D^{*}\Jone{\fs}}.

\emph{The matrix \m{D-C_{\fs}} is block upper triangular for the levels}, where $C_{\fs}$ is the chain matrix of \inp{13}(b).
By \eqref{e:Dframe}, \m{(D-C_{\fs})e_k=a\,e_k}, and \m{(D-C_{\fs})x_i=D\,x_i} lies in the span of \m{x_{\tilde e(*)}} and~$f_0$, of
levels at most~$\varrho$. The rows of level higher than that of a column $z$ or $e^{*}_{b}$ are rows $e^{*}_{c}$, and the
coefficient of $e^{*}_{c}$ in a vector~$w$ is \m{e_c^{*}\Jone{\fs}w}. By the symmetry of~$D$ and \eqref{e:Dframe},
$e_c^{*}\Jone{\fs}Dw=(a\,e_c+f_{c+1})^{*}\Jone{\fs}w$. For \m{w=z} this is $-1$ if \m{c=\varrho-1} and $0$ for the other
\m{c<\varrho}; for \m{w=e^{*}_b} and \m{c<b} it is $1$ if \m{c=b-1} and $0$ otherwise. These are the coefficients of
$e^{*}_c$ in \m{C_{\fs}z} and in \m{C_{\fs}e^{*}_b}.

\emph{Conditions (T1) to (T7)}, for \m{j=\bph{x}{\,\cdot\,}} and \m{\phi=\phi_{\fs}}; \m{\steep{l}{B}{e}} is a hypothesis. Let
\m{\sigma\in\bgr{\fs}}. By Lemma~\ref{l:gtrace}, $\tr\bigl((D-C_{\fs})(\sigma_1-\emb{\fs}{\sigma_2})\bigr)=0$. By \inp{13}(a),(b)
and Lemma~\ref{l:fst}, $\tr(C_{\fs}\sigma_1)=\sd{\sigma_1}-\overline{\sd{\sigma_1}}$ and $\tr(C_{\fs}\emb{\fs}{\sigma_2})=0$. And
by \eqref{e:Dframe} the block of~$D$ on the vectors~$x_i$ is diagonal, with entry $a$ at \m{\tilde e(*)} and $0$ elsewhere,
which is \m{\dg{\lowk{l'}}} by hypothesis; so $\tr\bigl(D(\emb{\fs}{\sigma_2}-1)\bigr)=\tr\bigl(\dg{\lowk{l'}}(\sigma_2-1)\bigr)$.
Adding the three facts,
\[\phs{l}{l'}{\bigl(q^{-1}(\sigma_1-1)q,\ \sigma_2-1\bigr)}=\tr\bigl(D(\sigma_1-1)\bigr)-\tr\bigl(\dg{\lowk{l'}}(\sigma_2-1)\bigr)
=\sd{\sigma_1}-\overline{\sd{\sigma_1}},\]
which is (T1). For (T2), $\tr(q^{-1}(\sigma_1-1)q)=\tr(\sigma_1-1)=\tr(\emb{\fs}{\sigma_2}-1)=\tr(\sigma_2-1)$, by
Lemma~\ref{l:gtrace} for \m{Z=1}. (T3) holds by \inp{13}(a) and Lemma~\ref{l:fst}, and (T4) is clear. (T5) and (T6) hold
because the signs on $\kps{\fs}$ are~$+1$ and $\lowk{l'}$ is $0$ on \m{\tilde e(\Xi)} and \m{a\ge0} at \m{\tilde e(*)}. For
(T7) let $u$ be a unitary \m{\Xi\times\Xi} matrix. By Lemma~\ref{l:block}, for the bijection~$\tilde e$ and the signs~$1$, there
is \m{h\in\Us{1}} with $h-1=\inc{\tilde e|_\Xi}(u-1)\inc{\tilde e|_\Xi}^{*}$. Take $\sigma=\bch{\fs}(1,h)=(\emb{\fs}{h},h)$. Its
second component gives the second equation of~(T7). For the first, \m{\emb{\fs}{h}-1} is \m{u-1} placed in the block
of the vectors \m{x_{\tilde e(\xi)}}; the matrix $q$ maps the basis vector at $e(\xi)$ to \m{x_{\tilde e(\xi)}}, and
$q^{-1}=\dg{\varepsilon}q^{*}\Jone{\fs}$ has the unit row at \m{x_{\tilde e(\xi)}} in its row $e(\xi)$; hence
$q^{-1}(\emb{\fs}{h}-1)q=\inc{e|_\Xi}(u-1)\inc{e|_\Xi}^{*}$.
\end{proof}
\leanline{\lean{existsisTwoFactorbesselPairHom}}

\begin{proof}[Proof of Theorem 5]
Write $\psi^{\fp}_{l}(g)=\dsc{l}{\fp^{-1}g\,\fp}$ for \m{g\in\Gone{\fs}}, \m{\mathsf{c}=8\,\rsh{\fs}^{2}+1}, and
$\mathsf{m}=\bms{\fs}\{\sigma:\op{\bph{x}{\sigma}_1-1}\le1\}$. For \m{y\in\Gone{\fs}} and \m{\sigma\in\bgr{\fs}} let
\[F(y,\sigma)=\psi^{\fp}_{l}\bigl(x^{-1}\sigma_1y\bigr)\,e^{-\phi_{\fs}(\sigma)}\,\mathrm{rep}(\sigma_2)^{-1}v_0\ \in\CC^{I},\qquad
W(y)=\int_{\bgr{\fs}}F(y,\sigma)\,d\bms{\fs}(\sigma).\]

\emph{Sizes.} From the hypothesis on~$v_0$ at \m{h=1} and \m{\dsc{l'}{1}=1} (the integral of Definition~\ref{d:coeff} at
\m{g=z=1} is the normalising constant), \m{\nrm{v_0}=1}. As $\mathrm{rep}$ is unitary and $\phi_{\fs}$ is imaginary,
$\nrm{F(y,\sigma)}=\abs{\psi^{\fp}_{l}(x^{-1}\sigma_1y)}$. The map \m{\sigma\mapsto\bph{x}{\sigma}_1} is a proper homomorphism
(Lemma~\ref{l:proper}), $\bms{\fs}$ is right invariant and finite on compact sets (Lemma~\ref{l:gmeasure}), and
$\bph{x}{\sigma}_1\cdot\fp^{-1}x^{-1}y\,\fp=\fp^{-1}x^{-1}\sigma_1y\,\fp$. So \inp{11}(b), with the exponent
$\mathsf{c}>8\abs{\Ind{\fs}}^{2}$, shows that $\mathsf{m}$ is finite and that for every~$y$
\begin{equation}\label{e:P4}
\int_{\bgr{\fs}}\Bigl(\frac{2}{\op{\fp^{-1}x^{-1}\sigma_1y\,\fp}}\Bigr)^{\mathsf{c}}d\bms{\fs}(\sigma)\le16^{\mathsf{c}}\,\mathsf{m} .
\end{equation}
By \inp{9}, with $\mathsf{c}\le500\,\rsh{\fs}^{3}+A$, the integrand of \eqref{e:P4} dominates \m{\nrm{F(y,\sigma)}}. Hence
\m{F(y,\,\cdot\,)} is integrable, $F$ being continuous, and
\begin{equation}\label{e:P3}
\nrm{W(y)}\le\int_{\bgr{\fs}}\bigabs{\psi^{\fp}_{l}(x^{-1}\sigma_1y)}\,d\bms{\fs}(\sigma)\le16^{\mathsf{c}}\,\mathsf{m}\qquad\text{for every }y .
\end{equation}

\emph{$W$ is a Bessel function for $\mathrm{rep}$.} It is bounded by \eqref{e:P3} and measurable. Let
\m{\sigma'\in\bgr{\fs}}. Replacing $\sigma$ by \m{\sigma\sigma'^{-1}} in the integral for \m{W(\sigma'_1y)}, which does not change
the measure, and using $\phi_{\fs}(\sigma\sigma'^{-1})=\phi_{\fs}(\sigma)-\phi_{\fs}(\sigma')$ from \inp{13}(a), we get
\[W(\sigma'_1\,y)=e^{\phi_{\fs}(\sigma')}\,\mathrm{rep}(\sigma'_2)\,W(y).\]
For \m{\sigma'=\bch{\fs}(u,1)=(u,1)} and for $\sigma'=\bch{\fs}(1,h)=(\emb{\fs}{h},h)$, where \m{\phi_{\fs}(\sigma')=0}, these are
the two transformation laws of Definition~\ref{d:shape}(g).

\emph{The value at~$x$.} Since $\ip{v_0}{\mathrm{rep}(h)^{-1}v_0}=\overline{\ip{v_0}{\mathrm{rep}(h)v_0}}=\overline{\dsc{l'}{h}}$,
\[\begin{aligned}\ip{v_0}{W(x)}&=\int_{\bgr{\fs}}\dsc{l}{\fp^{-1}x^{-1}\sigma_1x\,\fp}\;\overline{\dsc{l'}{\sigma_2}}\;e^{-\phi_{\fs}(\sigma)}\,d\bms{\fs}(\sigma)\\
&=\int_{\bgr{\fs}}\lapint{l}{l'}{\bph{x}{\,\cdot\,}}{\phi_{\fs}}{\sigma}\,d\bms{\fs}(\sigma).\end{aligned}\]
So the hypothesis of the theorem and the Cauchy--Schwarz inequality give
\begin{equation}\label{e:m1W}
\mathsf{m}\le b^{\,\mathfrak{w}}\,\nrm{W(x)}.
\end{equation}

\emph{The four conditions of Definition~\ref{d:peak}.} (P3) is the first inequality of \eqref{e:P3}. (P4) follows from
\eqref{e:P4}, \eqref{e:m1W} and \m{16\le b}, with \m{P=b^{\mathsf{c}+\mathfrak{w}}}. (P2): the set in the definition of~$\mathsf{m}$ contains a
neighbourhood of~$1$, and $\bms{\fs}$ is positive on open sets, so \m{\mathsf{m}>0} by \inp{11}(c), and \m{W(x)\ne0} by
\eqref{e:m1W}. (P1): by Fubini's theorem, which \eqref{e:P3} and the integrability of the kernel~$k$ of \inp{12} allow,
\[\begin{aligned}\int_{\Gone{\fs}}\!k(g^{-1})\,W(y\,g)\,dg=\int_{\bgr{\fs}}\!\Bigl(&\int_{\Gone{\fs}}\!k(g^{-1})\,\psi^{\fp}_{l}\bigl(x^{-1}\sigma_1y\,g\bigr)\,dg\Bigr)\\
&\cdot e^{-\phi_{\fs}(\sigma)}\,\mathrm{rep}(\sigma_2)^{-1}v_0\,d\bms{\fs}(\sigma),\end{aligned}\]
and the inner integral is $\psi^{\fp}_{l}(x^{-1}\sigma_1y)$ by \inp{12} with \m{y_2=1}, because \m{k=\fdeg{l}\,\psi^{\fp}_{l}} and
\m{k(1)=\fdeg{l}}.
\end{proof}
\leanline{\lean{existsisBesselFunction}\sep \lean{HasPeakAt}}

\sect{The period of a Poincaré series of deep level: proof of Theorem 6}{s:period}
The data are those of Theorem~6. We abbreviate \m{G=G(E,\varepsilon)}, $G_\infty=G_\infty(E,\varepsilon)$,
\m{V^{+}=V\oplus\langle1\rangle}, \m{\Gamma^{+}=\Gamma(V^{+})}, \m{\mathsf{K}=\mathsf{K}_\Omega(a_0)} and $\mathsf{K}^{+}=\mathsf{K}_{\Omega\sqcup\{*\}}(a_0)$,
write $b_{*}$ for the basis vector of the added line, and
\[k_\gamma(y)=\Psi\bigl(X_\gamma(y)_\infty\bigr)\,\mathbf{1}_{\mathsf{K}^{+}}\bigl(X_\gamma(y)_{\mathrm f}\bigr)\qquad(\gamma\in\Gamma^{+},\ y\in G).\]
Norms of matrices are operator norms. Besides \inp{9} and \inp{11} the proof uses these results of the background
library on adelic unitary groups (units 114, 120, 121, 122; all elementary except the first):
\begin{enumerate}[label=\textup{(A\arabic*)},leftmargin=2.9em,itemsep=1pt]
\item \emph{Unfolding.} If $f$ is integrable on~$G$ for the Haar measure, then \m{y\mapsto f(\delta y)} is
  $\varpi_V$-integrable for every \m{\delta\in\Gamma(V)}, and
  $\sum_{\delta\in\Gamma(V)}\int f(\delta y)\,d\varpi_V(y)=\int_Gf(y)\,dy$. (Weil's formula; reference and relation as
  for \inp{13}(d).)
\item \emph{The first minimum.} \m{\lambda_1(V)\le\nrm{\eta}_V} for every non-zero lattice vector~$\eta$ of~$V$; for
  spaces $V_1$, $V_2$ and points $g_1$, $g_2$ of their groups,
  $\min\bigl(\lambda_1(V_1\cdot g_1),\lambda_1(V_2\cdot g_2)\bigr)\le\lambda_1\bigl((V_1\oplus V_2)\cdot(g_1\oplus g_2)\bigr)$; and
  \m{\lambda_1=1} for a positive line.
\item \emph{Separated points.} If points $f_n$ of a real normed space of dimension~$d_0$ have norms and mutual
  distances at least \m{m_0>0}, and \m{c\ge d_0+2} is an integer, then $\sum_n\nrm{f_n}^{-c}\le(5/m_0)^{c}$.
\item \emph{Structure.} For \m{a_0\ne0}, $\mathsf{K}$ is an open subgroup of finite Haar measure, and $a_0$ is invertible in
  the finite adeles. An element \m{g_1\oplus g_2} lies in \m{\Gamma(V_1\oplus V_2)} if and only if \m{g_1\in\Gamma(V_1)} and
  \m{g_2\in\Gamma(V_2)}; and an element of the unitary group of a block-diagonal form \m{J_1\oplus J_2}, with $J_2$
  invertible, that fixes the vectors of the second block is of the form \m{g_1\oplus1}.
\end{enumerate}

\begin{proof}[Proof of Theorem 6]
\emph{Step 0: the decay of~$\Psi$.} Let $g\in G_\infty(E,\varepsilon\oplus1)$ and \m{t_0>0} be such that
$t_0\le\prod_v\op{g_v(r_v\oplus1)}$ for every \m{r\in G_\infty}. Then \m{r\mapsto\Psi(g\,(r\oplus1))} is integrable on
$G_\infty$ and
\begin{equation}\label{e:Psidecay}
\int_{G_\infty}\bigabs{\Psi\bigl(g\,(r\oplus1)\bigr)}\,dr\le L_{*}\,t_0^{-A}.
\end{equation}
Indeed, for \m{h\in\Us{\varepsilon\oplus1}} and a place~$v$, \inp{9} with the exponent \m{2A}, which is at most the gap
of~$l_v$, and $\op{h}\le\op{\alpha_v}\,\op{\alpha_v^{-1}h\,\omega_v}\,\op{\omega_v}$ from \inp{11}(d) give
\[\bigabs{\dsc{l_v}{\alpha_v^{-1}h\,\omega_v}}\le\Bigl(\frac{2\op{\alpha_v}\op{\omega_v}}{\op{h}}\Bigr)^{2A}
=\bigl(2(\op{\alpha_v}\op{\omega_v})^{2}\bigr)^{A}\Bigl(\frac{2}{\op{h}}\Bigr)^{A}\op{h}^{-A}.\]
Take \m{h=g_v(r_v\oplus1)} and multiply over~$v$: the last factors give at most~$t_0^{-A}$, and the integral over
$G_\infty$ of the product of the middle factors is the product over~$v$ of the integrals
$\int_{\Us{\varepsilon}}(2/\op{g_v(r\oplus1)})^{A}dr$. By \inp{11}(d) and the invariance of~$dr$ under inversion, such an
integral equals $\int(2/\op{(r\oplus1)g_v^{-1}})^{A}dr$, which by \inp{11}(b), for the proper homomorphism
\m{r\mapsto r\oplus1} and the exponent \m{A>8(\abs{\Omega}+1)^{2}}, is at most
\m{16^{A}\vol\{r:\op{r-1}\le1\}}. This proves \eqref{e:Psidecay}. Since \m{\op{h}\ge1} on unitary groups, it applies with
\m{t_0=1} to every~$g$; for \m{g=1} this is the first assertion of the theorem.

\emph{Step 1: the vector of a class.} Let \m{\gamma\in\Gamma^{+}}, say $\gamma=x_{V^{+}}^{-1}\gamma_0\,x_{V^{+}}$ with a rational
isometry~$\gamma_0$, and put $\eta_\gamma=\gamma_0\,b_{*}\in E^{\Omega\sqcup\{*\}}$. As the point of the positive line is~$1$
and \m{y\oplus1} fixes~$b_*$, the adelic vector of~$\eta_\gamma$ at the point of \m{V^{+}\cdot(x\oplus1)} is
\[\bigl(x_{V^{+}}(x\oplus1)\bigr)^{-1}\eta_\gamma=X_\gamma(y)\,b_{*}\qquad\text{for every }y\in G:\]
the last column of $X_\gamma(y)$, which does not depend on~$y$. Let $f(\gamma)\in\CC^{f_E(\abs{\Omega}+1)}$ be the vector of
its components at the infinite places; its Euclidean norm is the length of~$\eta_\gamma$ for \m{V^{+}\cdot(x\oplus1)}.

\emph{Step 2: representatives.} By (A4), \m{\Gamma(V)\oplus1} is the subgroup of the elements of~$\Gamma^{+}$ of the form
\m{g\oplus1}. Choose a set \m{\mathcal{C}\subseteq\Gamma^{+}} of representatives of the cosets \m{\gamma\,(\Gamma(V)\oplus1)}, with
\m{1\in\mathcal{C}}. Every element of $\Gamma^{+}$ is \m{\gamma(\delta\oplus1)} for unique \m{\gamma\in\mathcal{C}} and
\m{\delta\in\Gamma(V)}, and $X_{\gamma(\delta\oplus1)}(y)=X_\gamma(\delta y)$. As $F$ is invariant under~$\Gamma(V)$,
\begin{equation}\label{e:reindex}
\sum_{\gamma'\in\Gamma^{+}}k_{\gamma'}(y)\,F(y)=\sum_{\gamma\in\mathcal{C}}\ \sum_{\delta\in\Gamma(V)}k_\gamma(\delta y)\,F(\delta y),
\end{equation}
as soon as one side converges absolutely. Let $\mathcal{W}$ be the set of the \m{\gamma\in\mathcal{C}} for which
$X_\gamma(y)_{\mathrm f}\in\mathsf{K}^{+}$ for some~$y$; for the other~$\gamma$ the function $k_\gamma$ vanishes. For
\m{\gamma\in\mathcal{W}} fix such a point~$y_\gamma$, with \m{y_1=x}, for which \m{X_1(x)=1}.

\emph{Step 3: one class.} Let \m{\gamma\in\mathcal{W}}. Since $X_\gamma(y_\gamma y)=X_\gamma(y_\gamma)\,(y\oplus1)$, since
$\mathsf{K}^{+}$ is a group, and since $y_{\mathrm f}\oplus1\in\mathsf{K}^{+}$ if and only if \m{y_{\mathrm f}\in\mathsf{K}},
\[k_\gamma(y_\gamma\,y)=\Psi\bigl(X_\gamma(y_\gamma)_\infty\,(y_\infty\oplus1)\bigr)\cdot\mathbf{1}_{\mathsf{K}}(y_{\mathrm f})\qquad(y\in G).\]
By Step~0 and the left invariance of the Haar measure, $k_\gamma$ is integrable on~$G$, and
\begin{equation}\label{e:oneclass}
\int_G\abs{k_\gamma(y)}\,dy=\vol(\mathsf{K})\int_{G_\infty}\bigabs{\Psi\bigl(X_\gamma(y_\gamma)_\infty(r\oplus1)\bigr)}\,dr .
\end{equation}
For \m{\gamma=1}, the same substitution \m{y\mapsto x\,y} and the level of~$F$ give the main term:
\begin{equation}\label{e:main}\begin{aligned}
\int_Gk_1(y)\,F(y)\,dy&=\int_{G_\infty}\int_{\mathsf{K}}\Psi(r\oplus1)\,F\bigl(x\cdot(r,k)\bigr)\,dk\,dr\\
&=\vol(\mathsf{K})\int_{G_\infty}\Psi(r\oplus1)\,F\bigl(x\cdot(r,1)\bigr)\,dr .
\end{aligned}\end{equation}

\emph{Step 4: the vectors $f(\gamma)$, \m{\gamma\in\mathcal{W}}, are pairwise at distance at least~$Q^{2}$.} Let
\m{\gamma\ne\gamma'} in~$\mathcal{W}$. The matrices \m{X_\gamma(y_\gamma)_{\mathrm f}} and $X_{\gamma'}(y_{\gamma'})_{\mathrm f}$ lie
in~$\mathsf{K}^{+}$, so their last columns, which by Step~1 are the finite components of the adelic vectors of
$\eta_\gamma$ and $\eta_{\gamma'}$, are congruent to~$b_*$ modulo~$a_0$. Hence, $a_0$ being invertible in the finite adeles,
$\eta_\gamma-\eta_{\gamma'}=a_0\,\eta$ for a lattice vector~$\eta$ of \m{V^{+}\cdot(x\oplus1)}. It is not zero: otherwise
\m{\gamma_0^{-1}\gamma'_0} would fix~$b_*$ and by (A4) be of the form \m{g\oplus1}, so that $\gamma$ and $\gamma'$ would represent
the same coset. At an infinite place~$v$ the components of \m{a_0\eta} are those of~$\eta$ times the image of~$a_0$ under
the embedding chosen for~$v$, of absolute value \m{\abs{a_0}_v\ge\rho\ge0}. With (A2), hypothesis~(f), and
\m{\lambda_1(V\cdot x)\le1}:
\[\begin{aligned}\nrm{f(\gamma)-f(\gamma')}&\ \ge\ \rho\,\nrm{\eta}_{V^{+}\cdot(x\oplus1)}\ \ge\ \rho\,\lambda_1\bigl(V^{+}\cdot(x\oplus1)\bigr)\\
&\ \ge\ \rho\,\min\bigl(\lambda_1(V\cdot x),1\bigr)=\rho\,\lambda_1(V\cdot x)\ \ge\ Q^{2}.\end{aligned}\]

\emph{Step 5: the lattice sum.} The vector $f(1)$ has the component~$b_*$ at each infinite place, so
$\nrm{f(1)}=\sqrt{f_E}\le f_E\le Q$. By Step~4, $\nrm{f(\gamma)}\ge Q^{2}-Q\ge\tfrac12Q^{2}$ for
$\gamma\in\mathcal{W}\setminus\{1\}$. The real dimension of \m{\CC^{f_E(\abs{\Omega}+1)}} is \m{2f_E(\abs{\Omega}+1)\le A-2}, so
by (A3) with \m{m_0=\tfrac12Q^{2}}
\begin{equation}\label{e:latticesum}
\sum_{\gamma\in\mathcal{W}\setminus\{1\}}\nrm{f(\gamma)}^{-A}\le\Bigl(\frac{10}{Q^{2}}\Bigr)^{A}.
\end{equation}

\emph{Step 6: the other classes.} Let $\gamma\in\mathcal{W}\setminus\{1\}$, and let \m{g=X_\gamma(y_\gamma)_\infty}. For
\m{r\in G_\infty} the matrix \m{g_v(r_v\oplus1)} is the component at~$v$ of \m{X_\gamma(y_\gamma\cdot(r,1))}, whose last column
is the component $f(\gamma)_v$ of~$f(\gamma)$. Hence $\nrm{f(\gamma)_v}\le\op{g_v(r_v\oplus1)}$, and, these norms being at
least~$1$,
\[\nrm{f(\gamma)}\le\sqrt{f_E}\,\prod_v\op{g_v(r_v\oplus1)}\le Q\,\prod_v\op{g_v(r_v\oplus1)}.\]
So \eqref{e:Psidecay} applies with \m{t_0=\nrm{f(\gamma)}/Q}, and with \eqref{e:oneclass} and \m{\abs{F}\le C},
\begin{equation}\label{e:otherclass}
\int_G\abs{k_\gamma(y)\,F(y)}\,dy\le C\cdot\vol(\mathsf{K})\cdot L_{*}\cdot Q^{A}\,\nrm{f(\gamma)}^{-A}.
\end{equation}

\emph{Step 7: conclusion.} For \m{\gamma\in\mathcal{W}} the function \m{k_\gamma F} is integrable on~$G$, so by (A1) the family
$\delta\mapsto\int k_\gamma(\delta y)F(\delta y)\,d\varpi_V(y)$ sums to \m{\int_Gk_\gamma F}, and the same holds for the
absolute values. By \eqref{e:otherclass} and \eqref{e:latticesum} the numbers \m{\int_G\abs{k_\gamma F}}, \m{\gamma\in\mathcal{W}},
have a finite sum. So the double family of \eqref{e:reindex} is absolutely summable in the norm of $L^{1}(\varpi_V)$,
the sum and the integral may be exchanged, and
\[\int\sum_{\gamma'\in\Gamma^{+}}k_{\gamma'}(y)\,F(y)\,d\varpi_V(y)=\sum_{\gamma\in\mathcal{W}}\int_Gk_\gamma(y)\,F(y)\,dy .\]
The term \m{\gamma=1} is the main term, by \eqref{e:main}. The others are bounded by \eqref{e:otherclass} and
\eqref{e:latticesum}: their sum has absolute value at most
\[C\cdot\vol(\mathsf{K})\cdot L_{*}\cdot Q^{A}\Bigl(\frac{10}{Q^{2}}\Bigr)^{A}=C\cdot\vol(\mathsf{K})\cdot L_{*}\Bigl(\frac{10}{Q}\Bigr)^{A}
\le C\cdot\vol(\mathsf{K})\cdot L_{*}\cdot\frac{10}{Q},\]
because \m{10/Q\le1} and \m{A\ge1}.
\end{proof}
\leanline{\lean{normpoincareperiodsuble}}

\sect{Where the results are used}{s:use}
The six theorems are consumed by three other units of the argument; the places were found by a search of the
source for the formal names, and the statements of the consumers are those of their own notes.
\begin{itemize}[leftmargin=1.6em,itemsep=2pt]
\item \emph{Theorem C, the exchange step} (unit 147). The proof that every row pairs with a restriction of explicit
  level uses Theorem~6 for the unfolding of a period, and Theorem~2 for the data of Theorem~3, placed at every
  infinite place, to bound the main term of that period from below. There the exponent~$A$ of Theorem~2 is a power
  of the rank~$r$ of the rows.
\item \emph{Theorem B, the Bessel systems at the real places} (unit 148). A \emph{real place} is an infinite place
  of the maximal totally real subfield of the CM field of the data; the completion of the CM field above it is
  $\CC$, the unitary group of the rows there is $\Gone{\fs}$, and the notes on unit 132 treat it as the place $\CC$ of
  their formalism and call it the complex place. The existence of a Bessel system at such a place uses Theorem~4 for
  a point~$x$ of size at most \m{b^{5\nsh{\fs}}}, Theorem~2 on the Bessel graph, with the measure of
  Definition~\ref{d:gmeasure} and Lemmas \ref{l:gmeasure} and~\ref{l:proper}, for the parameter \m{\beta=b^{3}}, and
  Theorem~5 for the Bessel function. There $\fs$ is the shape of the admissible data $D$ of the overview, with
  \m{\rhs{\fs}=m(D)/2-1} and \m{\abs{\kps{\fs}}=k(D)}, so that \m{\rsh{\fs}=r(D)}, and $b$ is the third scale $b(D)$. The
  predicate of Definition~\ref{d:bestype} is the conclusion of that statement.
\item \emph{Theorem B, the cut of a Bessel function} (unit 132). A peaked Bessel function, in the sense of
  Definition~\ref{d:peak} with a point of size at most \m{b^{5\nsh{\fs}}}, is cut to a pair of a full and a cut Bessel
  function; the proof integrates by Lemma~\ref{l:unfold}.
\end{itemize}
Theorem~1 is used only inside this unit, in the proof of Theorem~2.

\sect{Notes on the formal text}{s:notes}
\subhead{Two norms.} The formal text uses two norms on matrices through scoped instances, and a statement does not
show which one it means. The files on two-factor data and on minors use the Hilbert--Schmidt norm throughout; the
files with Theorems 2, 3, 6 and with the second half of the Bessel statements use the operator norm; the file with
Theorem~1 opens both. That the hypothesis \m{\hs{g-1}\le\delta} of Theorem~1 is a Hilbert--Schmidt bound is forced by
the text: the proof of Theorem~2 passes the conclusion of Proposition~\ref{p:nearone}, which is elaborated with the
Hilbert--Schmidt norm, to Theorem~1 as this hypothesis, and the proof of Theorem~1 passes it on to
Proposition~\ref{p:minors}. The second reader of the project reached the same conclusion. Nothing was compiled for
this note, so this is a reading of the source, not Lean's answer.

\subhead{Statements.} (1) In Theorem~1 the last three hypotheses carry default proofs in the formal text, so that a
caller may omit them; they are hypotheses all the same. (2) The formal Theorem~2 asserts only the inequality,
with the measure of the ball read as a real number; that the ball has finite measure and that the integrand is
integrable is shown in its proof and added to the statement here. (3) In Theorem~3 the hypothesis that $M$ is
even follows from the other hypotheses: $e$ and $\tilde e$ are bijections, so the numbers of positive signs counted
through them, \m{\abs{\Xi}+\lceil M/2\rceil+1} and \m{\abs{\Xi}+\lceil(M+1)/2\rceil}, agree, which forces $M$ to be even.
The gap of~$l'$ enters only through \m{\lowk{-l'}=-\lowk{l'}}, and its size~$H$ is arbitrary. (4) In
Proposition~\ref{p:compress} and Proposition~\ref{p:minors} the constants $100$ and $200000$ are not sharp; the
proofs above give $45$ and $19302$. (5) In Lemma~\ref{l:fourth} the formal text takes the gap as an instance
argument, and it does not state the invertibility of the block-diagonal part. (6) In Lemma~\ref{l:unfold} the
point~$x$ is an implicit argument. (7) In Definition~\ref{d:twofactor} the formal text writes (T6) as a statement
about all pairs of indices; the form above is equivalent. (8) In Theorem~6 the eight auxiliary objects ($\Psi$, its
restriction, $L_*$, the sum of spaces, the embedding, the volume, $X_\gamma$ and the summand) are bound by
\lf{let} inside the statement; the final step loses a factor \m{(10/Q)^{A-1}}. The sum over~$\gamma$ is Lean's
unconditional sum, which is $0$ where the family is not summable; by Step~7 this happens at most on a null set.
(9) The definitions of a frame, of a peak and of the type of a Bessel system are structures; the frame is a
type class, and the signs are an output parameter of it.

\subhead{The arguments as the census reads them.} The reading project's census records the arguments of each
declaration from the text. For ten declarations of this unit its list differs from the one Lean must have, because
a name bound inside the statement coincides with a variable of the file, or because a \lf{let} or an arrow of
homomorphisms was misread: the pair $j-1$, two-factor data, Lemmas \ref{l:trace}, \ref{l:levelsums} and
\ref{l:unfold}, Proposition~\ref{p:minors}, the pair homomorphism, the peak, and Theorems 4 and~5. The statements
above follow Lean's text, as the first reader's notes describe the expected lists; a confirmation by Lean itself
(\lf{\#check}) is pending. In particular the measure of the Bessel graph and the peak do not depend on an element of
the graph, and Theorem~4 does not take a point~$x$ as an argument: it provides one.

\subhead{The dossier.} The first reader's text for the right invariance of the measure of the Bessel graph named
only the conjugation; the second reader pointed out the right translation by~$u_0$, and Lemma~\ref{l:gmeasure} has
both. In the proof of Theorem~2 the source leaves the two quantities of the case distinction to unification;
they are \m{\gamma(\sigma)^{8\mathsf{n}}} and \m{\beta^{-50\mathsf{n}}}, which is the only choice for which the next rewriting
steps of the source apply (the second reader's message makes the arithmetic explicit in a slightly different
way; both lead to the same statement).

\subhead{How this note was written.} From the dossier of the unit, after its program check passed: all
\UnitProofs{} structured proofs were read, and all five Lean files of the unit were read in full.
Line by line, against the structured proofs, we read the proofs of Propositions \ref{p:compress}, \ref{p:minors}
and~\ref{p:nearone} and of Theorems 1, 2, 3, 4 and~5. For Theorem~6 the statement was read in full and the proof
once; in its steps on the choice of representatives and on the exchange of sum and integral we follow the first
reader's reconstruction and the second reader's verdict, and did not trace every rewriting. The four statements
that the census cuts at a \lf{let} (Lemma~\ref{l:levelsums}, Proposition~\ref{p:minors}, Theorems 3 and~6) were read
in full in the source. The definitions of the background library that the statements rest on (parameters,
Flensted-Jensen's function, the shape and its groups, frames, Bessel functions, the size) were read in the
source; those of the adelic groups are taken from the dossier's readings. No Lean was run.

Several arguments here are shorter than, or different from, the formal ones. In Proposition~\ref{p:compress} the
formal proof verifies one identity between traces by a linear combination of the hypotheses; the three terms
$\mathcal{I}_1$, $\mathcal{I}_2$, $\mathcal{I}_3$ above are a derivation of that identity. In
Proposition~\ref{p:nearone} the formal proof checks the unitarity of~$u$ and the trace identity of Step~4 by linear
combinations; the block form and the argument with the vector \m{(v,v')} are ours. Both identities, and the two
of Proposition~\ref{p:compress}, were also tested numerically on random matrices, which is a check of the
transcription and no proof. The numerical estimates in Theorems 1 and~2 and in Proposition~\ref{p:minors} are done
by hand here with other intermediate constants than the formal ones. In Theorem~4 the vectors $V^{(k)}$ and the
description of the chain matrix are our reading of the matrices of the formal proof. The program check of the
unit's dossier covered \UnitSteps{} leaf steps of \UnitProofs{} proofs.

\head{Appendix A. Concordance with the Lean artifact}
{\fontsize{8.6}{11.6}\selectfont\renewcommand{\leanfont}{\ttfamily\fontsize{7.7}{10}\selectfont\color{leancol}}
\setlength{\LTleft}{0pt}\setlength{\LTright}{\fill}\setlength{\tabcolsep}{0pt}
\begin{longtable}{@{}p{30mm}p{56mm}p{43mm}r@{}}
\textbf{In this note} & \textbf{Lean declaration} & \textbf{File} & \textbf{Lines}\\\midrule\endhead
\ConcordanceRows
\end{longtable}}

\subhead{Legend.} The formal names of the variables, and the letters of this note. Index types \lf{n}, \lf{m},
\lf{X}: $\Omega$, $\Omega'$, $\Xi$. Signs \lf{ε}, \lf{η}: $\varepsilon$, $\varepsilon'$. Parameters \lf{l}, \lf{k}: $l$, $l'$; in
the Bessel statements \lf{u} is~$l'$. The map \lf{τ} of two-factor data: $e'$; the bijections \lf{e'} of Theorem~3
and \lf{τ} of Theorem~4: $\tilde e$. Elements \lf{s}, \lf{t} of~$\mathfrak{g}$: $\sigma$, $\sigma'$. The radius \lf{r}: $\delta$. Gaps
\lf{h}: $H$ if real, $A$ if natural. Projections \lf{P}, \lf{Q}: $p$, $q$; errors \lf{r}: $\eta$; elements \lf{k} of the
maximal compact subgroup: $h$. The frame \lf{p}: $\fp$; the vector \lf{v}: $v_0$; in Theorem~5 the exponent \lf{n}: $\mathfrak{w}$. In Theorem~6, \lf{n}, \lf{s}:
$\Omega$, $\varepsilon$; \lf{S}, \lf{a}, \lf{R}, \lf{d}: $C$, $a_0$, $\rho$, $Q$; places \lf{w}: $v$; \lf{φ}: $\Psi$; \lf{L}:
$L_*$; \lf{h}: $y$. Objects of the background library: \lf{USign}, \lf{maxCompactU}, \lf{blockDiagPart},
\lf{inclMatrix}, \lf{logOneAdd}, \lf{coordProj}, \lf{principalMinor} are $\Us{\varepsilon}$, $\Ks{\varepsilon}$,
$\bd{\varepsilon}{\,\cdot\,}$, $\inc{\iota}$, $\lgo{\,\cdot\,}$, $\cpr{\Pi}$, $\pmn{\Pi}{\,\cdot\,}$; \lf{paramAddRho},
\lf{lowestKType}, \lf{weightMinors}, \lf{flenstedJensen}, \lf{dsCoeff}, \lf{formalDegree} are $\ups{l}$, $\lowk{l}$,
$\wm{w}{\,\cdot\,}$, $\fj{l}{\,\cdot\,}{\,\cdot\,}$, $\dsc{l}{\,\cdot\,}$, $\fdeg{l}$; \lf{besselU}, \lf{besselH}, \lf{besselS},
\lf{embedH}, \lf{simpleRootSum}, \lf{besselHaar}, \lf{besselQuotMeasure}, \lf{besselSize} are $\Uu{\fs}$, $\Hh{\fs}$,
$\Qq{\fs}$, $\emb{\fs}{\,\cdot\,}$, $\sd{\,\cdot\,}$, $\hU{\fs}$, $\hq{\fs}$, $\bsz{\,\cdot\,}$.

\subhead{Two statements, as the artifact has them.} Theorem~1:
\quoted{dsCoeffgaussianshape}
Theorem~2:
\quoted{laplacemethod}

\head{Notation}
\begingroup\fontsize{8.6}{10.9}\selectfont\raggedright
\setlength{\columnsep}{9mm}\setlength{\parskip}{0pt}
\begin{multicols}{2}
\nt{$\Omega$, $\Omega'$, $\Xi$, $\Xi'$; $i$, $i'$, $i''$; $\xi$}{finite index sets; elements of $\Omega$ or $\Omega'$; an element of~$\Xi$}
\nt{$\Omega_1$, $\Pi$, $\Pi'$}{subsets of~$\Omega$ (the complement of $e(\Xi)$; lower and upper sets)}
\nt{$\varepsilon$, $\varepsilon'$, $\varepsilon\oplus1$, $1$}{families of signs; $\varepsilon$ extended by $+1$ at~$*$; the constant family}
\nt{$X$, $Y$, $Z$; $X^{*}$, $\comm{X}{Y}$}{matrices; conjugate transpose; commutator}
\nt{$\hs{X}$, $\op{X}$, $\nrm{v}$, $\tr$}{Hilbert--Schmidt norm, operator norm, Euclidean norm of a vector; trace}
\nt{$\dg{w}$, $\bd{\varepsilon}{X}$, $\inc{\iota}$, $\iota$}{diagonal matrix; block-diagonal part; the matrix of a map $\iota\colon\Xi\to\Omega$}
\nt{$\Us{\varepsilon}$, $\Ks{\varepsilon}$; $g$, $\tilde g$, $g_0$, $g_1$, $g_2$; $h$, $h_0$, $h_1$}{unitary group and its maximal compact subgroup; their elements; $g$ also a representative, a point or a vector in Theorem~6}
\nt{$dg$, $dh$, $dr$, $dk$, $\vol$}{Haar measures; volume}
\nt{$\lgo{X}$}{the logarithm of \m{1+X}}
\nt{$p$, $q$, $p'$, $\cpr{\Pi}$, $p_\tau$; $\pmn{\Pi}{y}$}{orthogonal projections; coordinate projections; principal minor; $q$ also the product $x\,\fp$ in §\ref{s:bessel}}
\nt{$v$; $\abs{w}_1$; $\tau$, $c_1$, $c_2$}{integer vectors; sum of absolute values of an integer vector $w$; integers; $v$ also a column in Lemmas and Propositions where this is said}
\nt{$l$, $l'$, $l_v$; $\lambda$, $\rho$, $\rho_c$, $\mu$}{integer vectors, read as doubled Harish-Chandra parameters; $l/2$, half sums of roots, and the name of $\ups{l}$ in the source (§\ref{s:setting} only)}
\nt{$\ups{l}$, $\lowk{l}$; $\Upsilon$, $\Lambda$, $\Lambda'$}{$\lambda+\rho$ and the lowest $K$-type; their values in a proof}
\nt{$H$, $A$, $A_0$, $L$}{a real gap; natural numbers: large gaps, a bound for the entries of~$l$}
\nt{$\wm{w}{g}$, $\ktn{\varepsilon}{w}$, $c_\Lambda$}{product of principal minors; its square norm on $\Ks{\varepsilon}$}
\nt{$\fj{l}{g}{z}$, $\dsc{l}{g}$, $\fdeg{l}$, $\psi^{\fp}_{l}$}{Flensted-Jensen's integral, his function, the formal degree; the function in a frame}
\nt{$\gwn{w}{X}$}{the gap-weighted square norm}
\nt{$\phs{l}{l'}{x}$}{the phase of a pair of matrices}
\nt{$\steep{l}{B}{e}$; $B$, $M$, $e$}{steep weights; a real number, a natural number, a bijection onto~$\Omega$}
\nt{$\mathfrak{g}$; $\sigma$, $\sigma'$, $\sigma_0$; $j=(j_1,j_2)$}{a group; its elements (in §\ref{s:bessel}, \m{\sigma=(\sigma_1,\sigma_2)}); a homomorphism}
\nt{$\phi$, $\phi_{\fs}$, $e'$, $\tilde e$}{phase functions; a map \m{\Xi\to\Omega'}; a bijection extending such a map}
\nt{$\lapint{l}{l'}{j}{\phi}{\sigma}$}{the integrand of two coefficients}
\nt{$\twof{l}{l'}{j}{e'}{B}{e}{\phi}$}{two-factor data}
\nt{$\beta$, $\beta_0$, $T$, $\delta$, $\delta_0$, $\kappa_{\mathrm{G}}$, $\kappa_0$, $\kappa_1$, $\kappa_2$}{real parameters; a radius; Gaussian factors ($\kappa_0$ a norm in §\ref{s:minors}); the bare letter $\kappa$ occurs only in $\kps{\fs}$}
\nt{$C$, $R$, $K$, $t_0$, $m_0$, $d_0$, $n_0$}{real numbers and bounds, natural numbers, as each statement says}
\nt{$\eta$, $\eta_1$, $\eta_2$, $\eta_3$, $\eta_\tau$, $r(h)$, $\epsilon$}{error terms (in §\ref{s:period}, $\eta$, $\eta_\gamma$ are vectors)}
\nt{$\mathcal{I}(s)$, $\mathcal{I}_n$, $W$, $V$, $A$, $A_c$, $A_n$, $E$, $R$, $G$, $X_1$, $X^{\uparrow}$, $X^{\downarrow}$, $\alpha_0$}{in §\ref{s:minors}: integrands, inverses, commutators and their parts, a norm}
\nt{$E$, $q_0$, $G_0(h)$, $a$}{in §\ref{s:gauss}: the exponent of the phase, the exponent of~$\kappa_{\mathrm{G}}$, a sum; an eigenvalue}
\nt{$\chi$, $c_0$, $J$, $J_2$, $g_{11},\dots,g_{22}$, $u$, $T$, $T'$, $Q$, $Y_\sigma$, $v'$}{in §\ref{s:twofactor}: an indicator, rows of \m{g-1}, sign matrices, blocks, a unitary matrix, \m{j(\sigma')-1}, $\inc{e|_\Xi}$, a pair}
\nt{$\mathsf{n}$, $\gamma(\sigma)$, $\operatorname{dist}(\sigma)$, $\mathsf{m}$, $\mathsf{J}$, $U_0$, $z$, $z_n$, $E_1$, $E_2$}{in §\ref{s:laplace}: $(\abs{\Omega}+\abs{\Omega'})^{2}$, \m{2/\op{j_1(\sigma)}}, a distance, the measure of the ball, the integral, a small ball, complex numbers}
\nt{$\tau_n$, $c_n$, $\pi$, $\pi_a$, $d_a$, $\rho_{na}$, $\omega$, $j^{\omega}$, $b_i$, $b_*$}{nodes, Lagrange weights, polynomials and their values; the element of Theorem~3 and its homomorphism; standard basis vectors}
\nt{$\fs$; $\rhs{\fs}$, $\kps{\fs}$, $\nsh{\fs}$, $\rsh{\fs}$, $\Ind{\fs}$, $\varrho$}{a shape and its numbers and index sets}
\nt{$e_a$, $x_i$, $z$, $e^{*}_a$, $f_a$; $a$, $a'$, $b$, $c$, $k$}{basis and frame of a shape; indices of basis vectors; $a$ also a weight in Definition~\ref{d:bestype}; $b$ also a real number in Theorem~6; $c$ also a real number or a natural number in Lemmas and Propositions where this is said}
\nt{$\Jone{\fs}$, $\Gone{\fs}$, $\lev{i}$, $\Uu{\fs}$, $\Hh{\fs}$, $\Qq{\fs}$, $\emb{\fs}{h}$, $\sd{g}$}{Gram matrix, unitary group, level, unipotent group, embedded compact group, Bessel subgroup, embedding, sum above the diagonal}
\nt{$\hU{\fs}$, $\hq{\fs}$, $\bsz{x}$, $C_{\fs}$}{Haar measure of $\Uu{\fs}$, quotient measure, size, chain matrix}
\nt{$\frm{\fs}{\varepsilon}$, $\fp$, $\fp_i$}{frames; a frame, its columns}
\nt{$\bgr{\fs}$, $\bch{\fs}$, $\bms{\fs}$, $\bph{x}{\sigma}$}{the Bessel graph, its chart, its measure, the pair homomorphism}
\nt{$I$, $\mathrm{rep}$, $W$, $v_0$, $y$, $\mathfrak{w}$}{an index set, a unitary representation, a Bessel function, a unit vector, points of $\Gone{\fs}$, the exponent of a peak (the same letter as in the note on the complex place)}
\nt{$n$}{an index of summation or of a node; $\nsh{\fs}=\rhs{\fs}+1$}
\nt{$m(D)$, $k(D)$, $r(D)$, $b(D)$}{numbers and third scale of the admissible data $D$ of the overview (only in §\ref{s:use})}
\nt{$\peak{l}{\mathrm{rep}}{W}{x}{P}$}{peak}
\nt{$\isbes{\mathcal{B}}{l}{w}$; $\mathcal{B}$, $I(\mathcal{B})$, $W_{\mathcal{B}}$, $\mathbb{S}_w$}{type of a Bessel system; its components; a Schur matrix}
\nt{$V^{(k)}$, $V_a$, $M_0$, $D$, $F(y,\sigma)$, $k(g)$, $\mathsf{c}$, $y_1$, $y_2$}{in §\ref{s:bessel}: vectors, a bound, a conjugated weight matrix, an integrand, a kernel, an exponent, points}
\nt{$E$, $f_E$, $v$, $\abs{a}_v$, $a_0$}{in Theorem~6 also a CM field, its number of infinite places, a place $v$, its absolute value, the level}
\nt{$G(E,\varepsilon)$, $G_\infty$, $G_{\mathrm f}$}{adelic groups}
\nt{$\mathsf{K}_\Omega(a_0)$, $\mathsf{K}$, $\mathsf{K}^{+}$}{compact open (congruence) subgroups of the adelic groups}
\nt{$V$, $J_V$, $x_V$, $V\cdot x$, $V^{+}$, $\langle1\rangle$, $\Gamma(V)$, $\Gamma^{+}$, $\varpi_V$, $\lambda_1$, $\nrm{\eta}_V$}{a space with a point, its parts, its translates and sum, discrete groups, quotient measure, first minimum, length}
\nt{$F$, $C$, $\rho$, $Q$, $\alpha$, $\omega$, $\Psi$, $L_*$, $X_\gamma$, $k_\gamma$, $\gamma$, $\gamma_0$, $\delta$, $r$, $k$}{in Theorem~6 also an automorphic form, bounds, positions, the test function, a constant, group elements}
\nt{$\mathcal{C}$, $\mathcal{W}$, $y_\gamma$, $f(\gamma)$, $f_n$, $\eta_\gamma$, $o$}{representatives, points, vectors; in \inp{4}, $o$ is an ordering function}
\nt{$D$, $P$, $d$, $f$, $s$, $w$, $x$, $\nu$}{in the Conventions of this note also local objects, not those of Conventions~0.1, 0.3, 0.4 and~0.6: $D$ a matrix,
$P$ a real number or a bound, $d$ values $d_a$ and bounds $d_0$, $f$ functions, $f_E$ and a Schur matrix, $s$ a real variable of integration in $[0,1]$ (not the complex variable of the
write-up), $w$ an integer vector, in Definition~\ref{d:bestype} a weight, $x$ the pair $(x_1,x_2)$ of matrices of a phase, an index set, a unit vector or a point of $\Gone{\fs}$, $\nu$ a measure}
\end{multicols}
\endgroup

\head{References}
{\fontsize{9.4}{12.6}\selectfont\raggedright
References are given as the second reader of the project and the catalogue of the write-up record them; none
was typed from memory, and the places not marked as confirmed there are to be confirmed.
\begin{itemize}[leftmargin=1.6em,itemsep=1pt]
\item U.~A.~First, \emph{Witt's extension theorem for quadratic spaces over semiperfect rings}, J.~Pure Appl.\ Algebra
  219 (2015); Remark~2.1 and Corollary~4.9 (preprint arXiv:1408.0522). For \inp{6}.
\item M.~Flensted-Jensen, \emph{Discrete series for semisimple symmetric spaces}, Ann.\ of Math.\ 111 (1980). For
  \inp{7}; place to be confirmed.
\item B.~Sun, \emph{Matrix coefficients of cohomologically induced representations}, Compositio Math.\ 143 (2007).
  Named by the catalogue for \inp{9}; not identified with it.
\item Harish-Chandra, \emph{Discrete series for semisimple Lie groups~II}, Acta Math.\ 116 (1966), pp.~87--88. For
  \inp{12}; relation unresolved.
\item O.~Kallenberg, \emph{Invariant measures and disintegrations with applications to Palm and related kernels},
  Probab.\ Theory Related Fields 139 (2007), 285--310; D.~Gentner and G.~Last, \emph{Palm pairs and the general
  mass-transport principle}, §2. For \inp{13}(d) and (A1); relation unresolved.
\item The header of the formal text of the unit names as related work: P.~Nelson and A.~Venkatesh, \emph{The orbit
  method and analysis of automorphic forms} (2021); H.~He, \emph{On the Gan--Gross--Prasad conjecture for
  $\mathrm{U}(p,q)$} (2017); B.~Sun (2007), as above. We have not compared them with the statements of this note.
\end{itemize}}

\end{document}
